% Lesson 14 — Grammar: Adverbs (for, since, ago, after, afterwards, etc.) — Anglais Tle A — Cours signé OnBuch+
% Programme officiel MINESEC. Compiler avec : tectonic EN14-grammar-adverbs-for-since-ago-after-afterwards-etc.tex
\newcommand{\DOCMATIERE}{Anglais}
\newcommand{\DOCNIVEAU}{Tle A}
\newcommand{\DOCMODULE}{Chapter 2 : Economic Life and Occupations}
\newcommand{\DOCLECON}{EN14}
\newcommand{\DOCTITRE}{Lesson 14 — Grammar: Adverbs (for, since, ago, after, afterwards, etc.)}
% OnBuch+ — préambule « cours signés » ENRICHI (compilé avec tectonic / XeLaTeX)
% Système visuel « Pop » d'OnBuch+ : fond crème (jamais blanc), bordures encre
% épaisses, ombres « dures » décalées, titres Archivo Black en capitales.
% Version MATHÉMATIQUES : graphes (pgfplots/tikz), boîtes pédagogiques
% (définition, propriété, méthode, attention, le savais-tu, exercice résolu,
% exercice, TP, bilan), page de garde et signature OnBuch+ ancrée en bas.
%
% Macros attendues dans main.tex AVANT \input{preamble} :
%   \DOCMATIERE  \DOCNIVEAU  \DOCMODULE  \DOCLECON (numéro)  \DOCTITRE
\documentclass[11pt]{article}

\usepackage{fontspec}
\usepackage[english,french]{babel} % français = langue principale ; l'anglais via \eng{..} / textebox
\usepackage[a4paper,margin=2cm,top=3.4cm,bottom=2.8cm,headheight=1.4cm,footskip=1.3cm]{geometry}
\usepackage{amsmath,amssymb,amsfonts}
\usepackage{mathtools} % \xmapsto, \coloneqq... souvent utilisés par les modèles
\usepackage{cancel} % simplifications barrées (bilans de forces)
% \abs{z} (module, valeur absolue) et \norm{u} (norme), utilisés par les modèles
\providecommand{\abs}{}\let\abs\relax\DeclarePairedDelimiter{\abs}{\lvert}{\rvert}
\providecommand{\norm}{}\let\norm\relax\DeclarePairedDelimiter{\norm}{\lVert}{\rVert}
\usepackage[table]{xcolor} % [table] : fournit aussi \rowcolors (alternance de lignes)
\usepackage{pagecolor}
\usepackage{tikz}
\usetikzlibrary{arrows.meta,positioning,calc,shapes.geometric,shapes.misc,shapes.arrows,decorations.pathmorphing,decorations.pathreplacing,decorations.markings,patterns,shadows,mindmap,trees,fit,backgrounds,matrix,intersections,angles,quotes}
\usepackage{pgfplots}
\usepackage[version=4]{mhchem} % équations nucléaires \ce{^{235}_{92}U}
\usepackage[european,siunitx]{circuitikz} % schémas électriques
\pgfplotsset{compat=1.18}
\usepgfplotslibrary{fillbetween,groupplots} % aires entre courbes (calcul intégral)
\usepackage{enumitem}
\usepackage{fancyhdr}
% [most] charge toutes les bibliothèques tcolorbox (listings, minted, raster,
% documentation...) et gonfle inutilement la mémoire TeX sur un document long
% (~300 boîtes) : on ne charge que ce qui est réellement utilisé.
\usepackage[skins,breakable]{tcolorbox}
\usepackage{graphicx}
\usepackage{colortbl}
\usepackage{booktabs}
\usepackage{tabularx}
\usepackage{diagbox} % case d'en-tête diagonale des tableaux à double entrée
% Colonnes extensibles centrée / gauche / droite (C, L, R), souvent utilisées par les modèles
\newcolumntype{C}{>{\centering\arraybackslash}X}
\newcolumntype{L}{>{\raggedright\arraybackslash}X}
\newcolumntype{R}{>{\raggedleft\arraybackslash}X}
\usepackage{array}
\usepackage{multicol}
\usepackage{multirow}
\usepackage{siunitx}
% parse-numbers=false : accepte aussi \SI{2,5 \times 10^{-3}}{...}
\sisetup{locale=FR,output-decimal-marker={,},group-minimum-digits=4,parse-numbers=false}
\DeclareSIUnit\ohm{\ensuremath{\symup{\Omega}}} % U+2126 absent de la police
\DeclareSIUnit\division{div} % réglages d'oscilloscope (ms/div, V/div)
\DeclareSIUnit\tr{tr} % tours (vitesse de rotation, ex. \SI{1500}{\tr\per\minute})
\DeclareSIUnit\molar{M} % concentration molaire (ex. \SI{3}{\molar}, \SI{1}{\micro\molar})
\DeclareSIUnit\an{an} % année (le modèle confond parfois avec \year, un compteur TeX) — ex. \SI{5}{\kilo\meter\per\an}
\DeclareSIUnit\FCFA{FCFA} % franc CFA (ex. \SI{500}{\FCFA\per\kilo\gram})
\DeclareSIUnit\degreeBrix{\textdegree Brix} % teneur en sucre d'un sirop/jus (réfractométrie)
\DeclareSIUnit\month{mois} % le modèle utilise parfois \month, un compteur TeX, comme unité de durée
\usepackage{qrcode}
% \degree : commande fréquemment utilisée par les modèles pour un angle ou une
% température en dehors de \SI{}{} (non fournie par les paquets chargés).
\providecommand{\degree}{^{\circ}}
% \qed (fin de démonstration), utilisé par les modèles sans amsthm
\providecommand{\qed}{\hfill$\square$}
% \barre : double barre des valeurs interdites dans un tableau de variations (tkz-tab)
\providecommand{\barre}{\ensuremath{\|}}
% environnement proof (amsthm non chargé) utilisé par les modèles
\ifdefined\proof\else\newenvironment{proof}[1][Démonstration]{\par\noindent\textit{#1.}\ }{\qed\par}\fi
\usepackage{lastpage}
\usepackage{hyperref}
\hypersetup{hidelinks}

% ============================================================
% Palette « Pop » OnBuch+ — mêmes hex que la classe P (plus_ui.dart)
% ============================================================
\definecolor{poporange}{HTML}{F59321}
\definecolor{popdark}{HTML}{E07A0C}
\definecolor{popcream}{HTML}{FFF6E8}
\definecolor{popink}{HTML}{1C1714}
\definecolor{poppurple}{HTML}{7A5AE0}
\definecolor{popgreen}{HTML}{2E9E6A}
\definecolor{poppink}{HTML}{E0457B}
\definecolor{popblue}{HTML}{2F7DE1}
\definecolor{popgold}{HTML}{C98A12}
\definecolor{popmuted}{HTML}{8A7F76}
\definecolor{popline}{HTML}{EADFD2}
\definecolor{popgray}{HTML}{8A7F76}
\colorlet{popgrey}{popgray}
% Alias tolérants (le modèle nomme parfois les couleurs différemment)
\colorlet{popwhite}{white}
\colorlet{popblack}{popink}
\colorlet{popred}{poppink}
\colorlet{popyellow}{popgold}
% Teintes claires (fonds de tableaux/encadrés) — le modèle nomme parfois
% les couleurs avec un suffixe L (ex. popblueL) sans les avoir définies.
\colorlet{poporangeL}{poporange!20}
\colorlet{popdarkL}{popdark!20}
\colorlet{poppurpleL}{poppurple!20}
\colorlet{popgreenL}{popgreen!20}
\colorlet{poppinkL}{poppink!20}
\colorlet{popblueL}{popblue!20}
\colorlet{popgoldL}{popgold!20}
\colorlet{popredL}{popred!20}
\colorlet{popgrayL}{popgray!20}
% Teintes claires pour fonds de tableaux / zones
\colorlet{popblueL}{popblue!10!white}
\colorlet{poporangeL}{poporange!14!white}
\colorlet{popgreenL}{popgreen!12!white}
\colorlet{poppurpleL}{poppurple!12!white}
\colorlet{poppinkL}{poppink!12!white}
\colorlet{popgoldL}{popgold!14!white}

\pagecolor{popcream}

% ============================================================
% Typographie
% ============================================================
\defaultfontfeatures{Ligatures=TeX}
\setmainfont{PlusJakartaSans-Medium.ttf}[Path=../../../fonts/,BoldFont=PlusJakartaSans-Bold.ttf]
% Maths en sans-serif (Fira Math) avec chiffres et lettres droites de Plus
% Jakarta, pour que les formules se fondent dans le texte.
\usepackage[math-style=ISO,bold-style=ISO]{unicode-math}
\setmathfont{FiraMath-Regular.otf}[Path=../../../fonts/]
\setmathfont{PlusJakartaSans-Medium.ttf}[Path=../../../fonts/,range={up/num,up/{Latin,latin},\mathup}]
\usepackage{icomma} % virgule décimale sans espace en mode math
% Caractères mathématiques Unicode absents des polices de texte (Plus Jakarta,
% Archivo Black) : rendus par leur équivalent mathématique, en texte comme en
% formule — sinon ils s'affichent en carré vide (ex. « Arithmétique dans ℤ »).
\usepackage{newunicodechar}
\newunicodechar{ℝ}{\ensuremath{\mathbb{R}}}
\newunicodechar{ℤ}{\ensuremath{\mathbb{Z}}}
\newunicodechar{ℂ}{\ensuremath{\mathbb{C}}}
\newunicodechar{ℕ}{\ensuremath{\mathbb{N}}}
\newunicodechar{ℚ}{\ensuremath{\mathbb{Q}}}
\newunicodechar{𝔻}{\ensuremath{\mathbb{D}}}
\newunicodechar{α}{\ensuremath{\alpha}}
\newunicodechar{β}{\ensuremath{\beta}}
\newunicodechar{γ}{\ensuremath{\gamma}}
\newunicodechar{δ}{\ensuremath{\delta}}
\newunicodechar{ε}{\ensuremath{\varepsilon}}
\newunicodechar{θ}{\ensuremath{\theta}}
\newunicodechar{λ}{\ensuremath{\lambda}}
\newunicodechar{μ}{\ensuremath{\mu}}
\newunicodechar{σ}{\ensuremath{\sigma}}
\newunicodechar{τ}{\ensuremath{\tau}}
\newunicodechar{φ}{\ensuremath{\varphi}}
\newunicodechar{ψ}{\ensuremath{\psi}}
\newunicodechar{ω}{\ensuremath{\omega}}
\newunicodechar{Σ}{\ensuremath{\Sigma}}
% majuscules produites par \MakeUppercase dans les titres (α-aminés -> Α)
\newunicodechar{Α}{\ensuremath{\alpha}}
\newunicodechar{Β}{\ensuremath{\beta}}
\newunicodechar{Γ}{\ensuremath{\Gamma}}
\newunicodechar{Δ}{\ensuremath{\Delta}}
\newunicodechar{Ε}{\ensuremath{\varepsilon}}
\newunicodechar{Θ}{\ensuremath{\Theta}}
\newunicodechar{Λ}{\ensuremath{\Lambda}}
\newunicodechar{Μ}{\ensuremath{\mu}}
\newunicodechar{Τ}{\ensuremath{\tau}}
\newunicodechar{Φ}{\ensuremath{\Phi}}
\newunicodechar{Ψ}{\ensuremath{\Psi}}
\newunicodechar{Ω}{\ensuremath{\Omega}}
\newunicodechar{↦}{\ensuremath{\mapsto}}
\newunicodechar{∈}{\ensuremath{\in}}
\newunicodechar{∉}{\ensuremath{\notin}}
\newunicodechar{∧}{\ensuremath{\wedge}}
\newunicodechar{⇔}{\ensuremath{\Leftrightarrow}}
\newunicodechar{⟺}{\ensuremath{\Longleftrightarrow}}
\newunicodechar{⇒}{\ensuremath{\Rightarrow}}
\newunicodechar{⟹}{\ensuremath{\Longrightarrow}}
\newunicodechar{∘}{\ensuremath{\circ}}
\newunicodechar{∪}{\ensuremath{\cup}}
\newunicodechar{∩}{\ensuremath{\cap}}
\newunicodechar{≡}{\ensuremath{\equiv}}
\newunicodechar{⊂}{\ensuremath{\subset}}
\newunicodechar{⊥}{\ensuremath{\perp}}
\newunicodechar{∃}{\ensuremath{\exists}}
\newunicodechar{∀}{\ensuremath{\forall}}
\newunicodechar{∛}{\ensuremath{\sqrt[3]{\,}}}
\newunicodechar{∜}{\ensuremath{\sqrt[4]{\,}}}
\newunicodechar{⃗}{\ensuremath{{}^{\to}}}
\newunicodechar{ⁿ}{\ifmmode^{n}\else\textsuperscript{\ensuremath{n}}\fi}
\newunicodechar{ˣ}{\ifmmode^{x}\else\textsuperscript{\ensuremath{x}}\fi}
\newunicodechar{ᵇ}{\ifmmode^{b}\else\textsuperscript{\ensuremath{b}}\fi}
\newunicodechar{ʳ}{\ifmmode^{r}\else\textsuperscript{\ensuremath{r}}\fi}
\newunicodechar{ᵘ}{\ifmmode^{u}\else\textsuperscript{\ensuremath{u}}\fi}
\newunicodechar{ʸ}{\ifmmode^{y}\else\textsuperscript{\ensuremath{y}}\fi}
\newunicodechar{ᵃ}{\ifmmode^{a}\else\textsuperscript{\ensuremath{a}}\fi}
\newunicodechar{ᵉ}{\ifmmode^{e}\else\textsuperscript{\ensuremath{e}}\fi}
\newunicodechar{ᵏ}{\ifmmode^{k}\else\textsuperscript{\ensuremath{k}}\fi}
\newunicodechar{ᵖ}{\ifmmode^{p}\else\textsuperscript{\ensuremath{p}}\fi}
\newunicodechar{ᴺ}{\ifmmode^{N}\else\textsuperscript{\ensuremath{N}}\fi}
\newunicodechar{ᶜ}{\ifmmode^{c}\else\textsuperscript{\ensuremath{c}}\fi}
\newunicodechar{ʰ}{\ifmmode^{h}\else\textsuperscript{\ensuremath{h}}\fi}
\newunicodechar{ᵠ}{\ifmmode^{q}\else\textsuperscript{\ensuremath{q}}\fi}
\newunicodechar{ₙ}{\ifmmode_{n}\else\textsubscript{\ensuremath{n}}\fi}
\newunicodechar{ₐ}{\ifmmode_{a}\else\textsubscript{\ensuremath{a}}\fi}
\newunicodechar{ᵢ}{\ifmmode_{i}\else\textsubscript{\ensuremath{i}}\fi}
\newunicodechar{ᵣ}{\ifmmode_{r}\else\textsubscript{\ensuremath{r}}\fi}
\newunicodechar{ₖ}{\ifmmode_{k}\else\textsubscript{\ensuremath{k}}\fi}
\newunicodechar{ᵦ}{\ifmmode_{\beta}\else\textsubscript{\ensuremath{\beta}}\fi}
\newunicodechar{ₓ}{\ifmmode_{x}\else\textsubscript{\ensuremath{x}}\fi}
\newfontfamily\popdisplay{ArchivoBlack-Regular.ttf}[Path=../../../fonts/]
\newfontfamily\poplogo{SpaceGrotesk-Bold.ttf}[Path=../../../fonts/]
\newfontfamily\popsemi{PlusJakartaSans-SemiBold.ttf}[Path=../../../fonts/]
\newfontfamily\popxbold{PlusJakartaSans-ExtraBold.ttf}[Path=../../../fonts/]
\newfontfamily\popxboldls{PlusJakartaSans-ExtraBold.ttf}[Path=../../../fonts/,LetterSpace=3.0]

\setlist[enumerate]{leftmargin=1.7em,itemsep=5pt,topsep=4pt}
\setlist[itemize]{leftmargin=1.7em,itemsep=4pt,topsep=4pt,label={\color{poporange}\textbullet}}
\setlength{\parindent}{0pt}
\setlength{\parskip}{5pt}
\renewcommand{\arraystretch}{1.25}

% pgfplots : style Pop par défaut
\pgfplotsset{
  every axis/.append style={axis line style={popink,line width=1pt},tick style={popink},
    grid style={popline,line width=0.5pt},label style={font=\small},tick label style={font=\footnotesize}},
  popcurve/.style={poporange,line width=1.8pt,no markers,smooth},
}
% Même style disponible dans un \draw TikZ simple (hors pgfplots)
\tikzset{popcurve/.style={poporange,line width=1.8pt}}
\tikzset{popgrid/.style={popline,very thin}}
\colorlet{popcurve}{poporange} % le nom du style est parfois utilisé comme couleur

% ============================================================
% Pill / squiggle
% ============================================================
\newcommand{\poppill}[3]{%
  \tikz[baseline=-0.55ex]\node[draw=popink,line width=1.2pt,rounded corners=2.6pt,%
    fill=#2,inner xsep=6.5pt,inner ysep=3pt]{%
    \popxboldls\fontsize{7}{7}\selectfont\color{#3}\MakeUppercase{#1}};%
}
\newcommand{\popsquiggle}[1]{%
  \tikz[baseline=-0.4ex]{
    \draw[#1,line width=2.6pt,line cap=round]
      plot [smooth] coordinates {(0,0.09) (0.34,0.01) (0.68,0.21) (1.02,0.01) (1.36,0.10)};
  }%
}
% Mot-clé surligné (marqueur orange sous le texte)
\newcommand{\cle}[1]{{\popxbold\color{popdark}#1}}

% ============================================================
% En-tête / pied
% ============================================================
\pagestyle{fancy}
\fancyhf{}
\renewcommand{\headrulewidth}{0pt}
\renewcommand{\footrulewidth}{0pt}
\lhead{\raisebox{-0.15ex}{\poplogo\fontsize{13}{13}\selectfont\color{popink}OnBuch\color{poporange}+}}
\rhead{\poppill{\DOCMATIERE\ \textbullet\ \DOCNIVEAU\ \textbullet\ Leçon \DOCLECON}{white}{popink}}
\lfoot{\popxbold\fontsize{7.6}{7.6}\selectfont\color{popmuted}\MakeUppercase{Cours signé OnBuch+}\ \textbullet\ {\popsemi\DOCTITRE}}
\rfoot{\tikz[baseline=-0.6ex]\node[rounded corners=8.5pt,fill=popink,minimum height=17pt,minimum width=17pt,inner xsep=5pt,inner sep=0]{\popxbold\fontsize{8}{8}\selectfont\color{popcream}\thepage\,/\,\pageref*{LastPage}};}

% ============================================================
% Titres
% ============================================================
\newcommand{\chaptitre}[2]{%
  \begin{tcolorbox}[enhanced,colback=poporange,colframe=popink,boxrule=2.6pt,arc=11pt,
      drop shadow=popink,left=15pt,right=15pt,top=12pt,bottom=11pt,before skip=3pt,after skip=18pt]
    \poppill{#2}{popink}{popcream}\par\vspace{9pt}
    {\raggedright\hyphenpenalty=10000\exhyphenpenalty=10000\popdisplay\fontsize{22}{24}\selectfont\color{popcream}\MakeUppercase{#1}\par}
  \end{tcolorbox}
}
% Section numérotée : pastille orange avec le numéro + titre Archivo
\newcounter{popsec}
\newcommand{\coursec}[1]{%
  \stepcounter{popsec}\setcounter{popsub}{0}%
  \par\vspace{16pt}\noindent
  \begin{minipage}{\linewidth}
  \tikz[baseline=-0.7ex]\node[draw=popink,line width=1.6pt,fill=poporange,rounded corners=4pt,minimum size=22pt,inner sep=0]{\popdisplay\fontsize{12}{12}\selectfont\color{popcream}\thepopsec};%
  \hspace{8pt}{\popdisplay\fontsize{15}{17}\selectfont\color{popink}\MakeUppercase{#1}}\par
  \vspace{4pt}\noindent{\color{poporange}\rule{36pt}{3.6pt}}
  \end{minipage}\par\nopagebreak\vspace{8pt}
}
\newcounter{popsub}
\newcommand{\courssub}[1]{%
  \stepcounter{popsub}%
  \par\vspace{9pt}\noindent{\popxbold\fontsize{11.5}{13}\selectfont\color{popink}{\color{poporange}\thepopsec.\thepopsub}\hspace{6pt}#1}\par\nopagebreak\vspace{4pt}
}

% ============================================================
% Boîtes pédagogiques — même recette : fond blanc, bordure épaisse colorée,
% ombre dure encre, badge plein. Titre optionnel [#1].
% ============================================================
\tcbset{popbase/.style={breakable,enhanced,colback=white,boxrule=2.3pt,arc=9pt,
  shadow={3pt}{-3pt}{0pt}{popink},left=14pt,right=14pt,top=11pt,bottom=11pt,before skip=11pt,after skip=15pt,
  fontupper=\popsemi\fontsize{10.6}{14.5}\selectfont\color{popink},
  fontlower=\popsemi\fontsize{10.6}{14.5}\selectfont\color{popink}}}
% Coche dessinée (le glyphe ✓ n'existe pas dans les polices du document)
\newcommand{\popcheck}[1]{\tikz[baseline=-0.5ex]\draw[#1,line width=1.6pt,line cap=round,line join=round](-0.13,0.0)--(-0.04,-0.09)--(0.13,0.1);}
\providecommand{\coche}{\popcheck{popgreen}}
\providecommand{\resultat}[1]{\par\smallskip\noindent\fcolorbox{popgreen}{popgreenL}{\parbox{\dimexpr\linewidth-2\fboxsep-2\fboxrule\relax}{\textbf{Résultat.} #1}}\par\smallskip}
\newcommand{\popboxhead}[3]{\poppill{#1}{#2}{white}\ifx\relax#3\relax\else\hspace{6pt}{\popxbold\fontsize{10.6}{12}\selectfont\color{popink}#3}\fi\par\vspace{7pt}}

\newtcolorbox{aretenir}[1][]{popbase,colframe=popgreen,before upper={\popboxhead{À retenir}{popgreen}{#1}}}
% Texte en anglais (document de lecture, dialogue, extrait) : cadre
% d'étude. Un titre optionnel (source, auteur) s'affiche dans l'en-tête.
\newtcolorbox{textebox}[1][]{popbase,colframe=popblue,before upper={\popboxhead{Text}{popblue}{#1}\begin{otherlanguage}{english}},after upper={\end{otherlanguage}}}
% Marques ✓ / ✗ (forme correcte / fautive) souvent utilisées par les modèles
\providecommand{\cross}{\textcolor{poppink}{\ensuremath{\times}}}
\providecommand{\xmark}{\cross}\providecommand{\crossmark}{\cross}\providecommand{\cmark}{\ensuremath{\checkmark}}
% Ligne de réponse / justification à compléter (inventée par les modèles)
\providecommand{\justification}[1]{\par\noindent\textit{Justification~:} \dotfill\par}
% \textipa (package tipa non chargé) : la police ne couvre pas l'API -> simple passage
\providecommand{\textipa}[1]{#1}
% Soulignement (ulem non chargé) : \uline, \ul
\providecommand{\uline}[1]{\underline{#1}}\providecommand{\ul}[1]{\underline{#1}}
% \glossaire{\item ...} inventée par les modèles : liste de glossaire
\providecommand{\glossaire}[1]{\begin{itemize}#1\end{itemize}}
% \alert (beamer) inventée par les modèles : texte mis en évidence
\providecommand{\alert}[1]{\textbf{\textcolor{poppink}{#1}}}
% variantes ASCII d'ordinaux écrites en macro
\providecommand{\iere}{\textsuperscript{re}}\providecommand{\ieme}{\textsuperscript{e}}\providecommand{\ere}{\textsuperscript{re}}\providecommand{\eme}{\textsuperscript{e}}\providecommand{\er}{\textsuperscript{er}}\providecommand{\ie}{\textsuperscript{e}}
% Ordinaux français écrits en macro par les modèles : 1\ière, 3\ième
\newcommand{\ière}{\textsuperscript{re}}\newcommand{\ième}{\textsuperscript{e}}
% \exemple (inventée par les modèles) : étiquette « Exemple : »
\providecommand{\exemple}{\textbf{Exemple~:}\ }
% \square utilisable aussi hors mode math (cases à cocher)
\AtBeginDocument{\let\squareMath\square\renewcommand{\square}{\ensuremath{\squareMath}}}
% Mot ou phrase en anglais à l'intérieur d'un paragraphe français
\newcommand{\eng}[1]{\ifmmode\text{\textit{#1}}\else\begingroup\selectlanguage{english}\itshape #1\endgroup\fi}
\let\esp\eng % alias toléré
\newtcolorbox{exemplebox}[1][]{popbase,colframe=poporange,before upper={\popboxhead{Exemple}{poporange}{#1}}}
\newtcolorbox{definition}[1][]{popbase,colframe=popblue,before upper={\popboxhead{Définition}{popblue}{#1}}}
\newtcolorbox{propriete}[1][]{popbase,colframe=poppurple,before upper={\popboxhead{Propriété}{poppurple}{#1}}}
\newtcolorbox{methode}[1][]{popbase,colframe=popgold,before upper={\popboxhead{Méthode}{popgold}{#1}}}
\newtcolorbox{attention}[1][]{popbase,colframe=poppink,before upper={\popboxhead{Attention}{poppink}{#1}}}
\newtcolorbox{experience}[1][]{popbase,colframe=popblue,colback=popblueL,before upper={\popboxhead{Expérience}{popblue}{#1}}}
% « Le savais-tu ? » : carte violette pleine (contexte, culture, Cameroun)
\newtcolorbox{savaistu}[1][]{popbase,colback=poppurple,colframe=popink,
  fontupper=\popsemi\fontsize{10.4}{14.2}\selectfont\color{white},
  before upper={\poppill{Le savais-tu\,?}{white}{poppurple}\ifx\relax#1\relax\else\hspace{6pt}{\popxbold\color{white}#1}\fi\par\vspace{7pt}}}
% Exercice résolu : énoncé en haut, \tcblower puis la solution sur fond crème
\newcounter{popexo}\newcounter{popexr}
\newtcolorbox{exoresolu}[1][]{popbase,colframe=poporange,colbacklower=popcream,
  segmentation style={popink,line width=1.2pt,dashed},
  before upper={\stepcounter{popexr}\popboxhead{Exercice résolu \thepopexr}{poporange}{#1}},
  before lower={\poppill{Solution}{popink}{popcream}\par\vspace{6pt}}}
% Exercice (non résolu en place) — la correction est dans la partie Corrigés
\newtcolorbox{exercice}[1][]{popbase,colframe=popink,
  before upper={\stepcounter{popexo}\popboxhead{Exercice \thepopexo}{popink}{#1}}}
% Corrigé
\newtcolorbox{corrige}[1][]{popbase,colframe=popgreen,colback=popgreenL,
  before upper={\popboxhead{Corrigé}{popgreen}{#1}}}
% Objectifs / prérequis / bilan
\newtcolorbox{objectifs}{popbase,colframe=popink,colback=poporangeL,
  before upper={\popboxhead{Objectifs de la leçon}{popink}{}}}
\newtcolorbox{prerequis}{popbase,colframe=popink,colback=popblueL,
  before upper={\popboxhead{Prérequis}{popblue}{}}}
\newtcolorbox{bilan}{popbase,colframe=popink,colback=white,boxrule=2.8pt,
  before upper={\popboxhead{Fiche bilan}{poporange}{L'essentiel à connaître pour l'examen}}}
% Figure encadrée avec légende
\newtcolorbox{popfigure}[1][]{enhanced,breakable=false,colback=white,colframe=popline,boxrule=1.4pt,arc=8pt,
  left=10pt,right=10pt,top=10pt,bottom=8pt,before skip=10pt,after skip=12pt,halign=center,
  lower separated=false,halign lower=center,
  fontlower=\popsemi\fontsize{9}{11.5}\selectfont\color{popmuted},
  #1}
\newcommand{\legende}[1]{\par\vspace{6pt}{\popsemi\fontsize{9}{11.5}\selectfont\color{popmuted}#1}}

% ============================================================
% Signature OnBuch+
% ============================================================
\newcommand{\poplogoinline}[1]{{\poplogo\fontsize{#1}{#1}\selectfont\color{popink}OnBuch\color{poporange}+}}
\newcommand{\signatureonbuch}{%
  \begin{tcolorbox}[enhanced,colback=popink,colframe=popink,boxrule=2.2pt,arc=10pt,
      drop shadow=poporange,left=14pt,right=14pt,top=10pt,bottom=10pt,before skip=0pt,after skip=0pt]
    \tikz[baseline=-0.7ex]\node[circle,fill=popgreen,draw=popcream,line width=1.3pt,minimum size=22pt,inner sep=0]{\popcheck{white}};%
    \hspace{9pt}\begin{minipage}[c]{0.62\linewidth}
      {\popxbold\fontsize{12}{13}\selectfont\color{popcream}Signé\ \textbullet\ {\poplogo OnBuch\color{poporange}+}}\par\vspace{2pt}
      {\raggedright\popsemi\fontsize{8}{10.5}\selectfont\color{popline}\MakeUppercase{Équipe pédagogique OnBuch+}\\\MakeUppercase{Conforme au programme officiel MINESEC}\par}
    \end{minipage}\hfill
    {\poplogo\fontsize{22}{22}\selectfont\color{popcream}OnBuch\color{poporange}+}
  \end{tcolorbox}%
}

% ============================================================
% Page de garde
% #1 titre · #2 matière · #3 niveau · #4 module · #5 n° leçon · #6 durée ·
% #7 description / objectifs (texte ou itemize)
% ============================================================
\newcommand{\pagegarde}[7]{%
  \thispagestyle{empty}
  \newgeometry{a4paper,margin=1.7cm,top=1.6cm,bottom=1.4cm}%
  \begin{tikzpicture}[remember picture,overlay]
    \fill[poporange!18] ([xshift=-3.2cm,yshift=-3.6cm]current page.north east) circle (4.6cm);
    \fill[poppurple!14] ([xshift=2.4cm,yshift=7.6cm]current page.south west) circle (3.8cm);
  \end{tikzpicture}%
  {\poplogo\fontsize{30}{30}\selectfont\color{popink}OnBuch\color{poporange}+}\hfill
  \raisebox{0.6ex}{\poppill{Cours signé \textbullet\ Premium}{popink}{popcream}}\par
  \vspace{1pt}\popsquiggle{poppurple}\par
  \vspace{8pt}
  \begin{tcolorbox}[enhanced,colback=poporange,colframe=popink,boxrule=2.8pt,arc=13pt,
      drop shadow=popink,left=18pt,right=18pt,top=12pt,bottom=13pt,after skip=10pt]
    \poppill{#2 \textbullet\ #3}{popink}{popcream}\hspace{5pt}\poppill{#4}{white}{popink}\par\vspace{8pt}
    {\popdisplay\fontsize{14}{14}\selectfont\color{popink}LEÇON #5}\par\vspace{4pt}
    {\raggedright\hyphenpenalty=10000\exhyphenpenalty=10000\popdisplay\fontsize{25}{27}\selectfont\color{popcream}\MakeUppercase{#1}\par}
    \vspace{8pt}
    {\popxbold\fontsize{9.5}{9.5}\selectfont\color{popink}\MakeUppercase{Durée conseillée : #6}}
  \end{tcolorbox}
  \begin{tcolorbox}[enhanced,colback=white,colframe=popink,boxrule=2.2pt,arc=10pt,
      drop shadow=popink,left=16pt,right=16pt,top=10pt,bottom=10pt,after skip=8pt]
    \poppill{Dans cette leçon}{poppurple}{white}\par\vspace{5pt}
    \popsemi\fontsize{9.3}{12.6}\selectfont\color{popink}#7
  \end{tcolorbox}
  \noindent\begin{minipage}[t]{0.487\linewidth}
    \begin{tcolorbox}[enhanced,colback=popgreen,colframe=popink,boxrule=2.2pt,arc=10pt,
        drop shadow=popink,left=11pt,right=11pt,top=10pt,bottom=10pt,height=3.1cm,valign=center]
      \tcbox[colback=white,colframe=popink,boxrule=1.3pt,arc=5pt,boxsep=0pt,nobeforeafter,
        left=3pt,right=3pt,top=3pt,bottom=3pt,valign=center]{\qrcode[height=1.9cm]{https://play.google.com/store/apps/details?id=com.luvvix.onbuch&hl=fr}}%
      \hspace{8pt}\begin{minipage}[c]{0.5\linewidth}
        {\popdisplay\fontsize{11}{12.5}\selectfont\color{white}\MakeUppercase{Télécharge OnBuch}}\par\vspace{4pt}
        {\raggedright\popsemi\fontsize{8.4}{11}\selectfont\color{white}Gratuit \textbullet\ Google Play\\Tuteur IA, annales, concours, résultats\par}
      \end{minipage}
    \end{tcolorbox}
  \end{minipage}\hfill
  \begin{minipage}[t]{0.487\linewidth}
    \begin{tcolorbox}[enhanced,colback=poppurple,colframe=popink,boxrule=2.2pt,arc=10pt,
        drop shadow=popink,left=11pt,right=11pt,top=10pt,bottom=10pt,height=3.1cm,valign=center]
      \tikz[baseline=-0.7ex]\node[circle,fill=white,draw=popink,line width=1.3pt,minimum size=1.6cm,inner sep=0]{\popdisplay\fontsize{24}{24}\selectfont\color{poppurple}+};%
      \hspace{8pt}\begin{minipage}[c]{0.6\linewidth}
        {\popdisplay\fontsize{11}{12.5}\selectfont\color{white}\MakeUppercase{Passe à OnBuch+}}\par\vspace{4pt}
        {\raggedright\popsemi\fontsize{8.4}{11}\selectfont\color{white}Tous les cours signés, Tuteur IA illimité, fiches PDF — onglet OnBuch+ de l'app\par}
      \end{minipage}
    \end{tcolorbox}
  \end{minipage}
  \vspace{8pt}
  \popcontenu
  \vfill
  \signatureonbuch
  \restoregeometry
  \newpage
}

% Rangée « Ce cours contient » (page de garde)
\newcommand{\poptile}[3]{% #1 couleur #2 gros texte #3 légende
  \begin{tcolorbox}[enhanced,colback=#1,colframe=popink,boxrule=2pt,arc=9pt,shadow={2.5pt}{-2.5pt}{0pt}{popink},
      left=6pt,right=6pt,top=9pt,bottom=9pt,halign=center,valign=center,height=1.75cm,width=\linewidth,nobeforeafter]
    {\popdisplay\fontsize{12.5}{13}\selectfont\color{white}\MakeUppercase{#2}}\par\vspace{3pt}
    {\centering\popsemi\fontsize{7.8}{9.5}\selectfont\color{white}#3\par}
  \end{tcolorbox}}
\newcommand{\popcontenu}{%
  \par\noindent{\popxboldls\fontsize{7.5}{7.5}\selectfont\color{popmuted}CE COURS SIGNÉ CONTIENT}\par\vspace{7pt}
  \noindent\begin{minipage}[t]{0.235\linewidth}\poptile{popblue}{Cours}{complet, illustré, pas à pas}\end{minipage}\hfill
  \begin{minipage}[t]{0.235\linewidth}\poptile{poporange}{Méthodes}{et exercices résolus}\end{minipage}\hfill
  \begin{minipage}[t]{0.235\linewidth}\poptile{poppink}{Exercices}{type examen + corrigés}\end{minipage}\hfill
  \begin{minipage}[t]{0.235\linewidth}\poptile{popgreen}{Fiche bilan}{l'essentiel à retenir}\end{minipage}\par}

% Sommaire
\newtcolorbox{sommaire}{enhanced,colback=white,colframe=popink,boxrule=2.2pt,arc=10pt,
  drop shadow=popink,left=18pt,right=18pt,top=13pt,bottom=13pt,before skip=6pt,after skip=20pt,
  before upper={\poppill{Sommaire}{poppurple}{white}\par\vspace{9pt}\popsemi\fontsize{10.6}{14.5}\selectfont\color{popink}}}

% Signature de fin de cours — poussée tout en bas de la dernière page
\newcommand{\findecours}{%
  \par\vfill\nopagebreak
  \begin{center}\popsquiggle{poporange}\end{center}\vspace{4pt}
  \signatureonbuch
}
\AtBeginDocument{\def\llbracket{\mathopen{[\mkern-3.5mu[}}\def\rrbracket{\mathclose{]\mkern-3.5mu]}}} % intervalles d'entiers (glyphe ⟦ absent de la police)
% Glyphes absents de Fira Math : redéfinis à partir de glyphes disponibles
\AtBeginDocument{%
  \def\setminus{\mathbin{\backslash}}\let\smallsetminus\setminus
  \def\vdots{\mathord{\vbox{\baselineskip3.2pt\lineskiplimit0pt\kern4pt\hbox{.}\hbox{.}\hbox{.}}}}%
  \def\ddots{\mathinner{\mkern1mu\raise6.5pt\hbox{.}\mkern2mu\raise3.6pt\hbox{.}\mkern2mu\raise0.7pt\hbox{.}\mkern1mu}}%
  \def\triangle{\mathord{\scalebox{0.9}{$\symup{\Delta}$}}}%
  \def\nabla{\mathord{\raisebox{\depth}{\scalebox{1}[-1]{$\symup{\Delta}$}}}}%
  \def\checkmark{\popcheck{popdark}}%
  \def\star{\mathbin{\ast}}\def\bigstar{\mathord{\ast}}%
  \def\diamond{\mathbin{\scalebox{0.6}{$\Diamond$}}}%
  \def\bigcup{\mathop{\mathchoice{\raisebox{-0.3em}{\scalebox{1.7}{$\cup$}}}{\scalebox{1.25}{$\cup$}}{\cup}{\cup}}\displaylimits}%
  \def\bigcap{\mathop{\mathchoice{\raisebox{-0.3em}{\scalebox{1.7}{$\cap$}}}{\scalebox{1.25}{$\cap$}}{\cap}{\cap}}\displaylimits}%
  \def\lesseqqgtr{\mathrel{\lessgtr}}\def\bigoplus{\mathop{\oplus}\displaylimits}%
}
\providecommand{\ssi}{si et seulement si} % abréviation fréquente
\AtBeginDocument{\providecommand{\wideparen}{\overparen}} % arc de cercle

%% FIGFIX-BEGIN v5 — correctifs globaux des figures (docs/onbuch-generation/tools/figfix)
\usepackage{adjustbox}\usepackage{etoolbox}\tcbuselibrary{raster}
% toute figure TikZ/pgfplots plus large que la ligne est réduite à la largeur disponible
\BeforeBeginEnvironment{tikzpicture}{\begin{adjustbox}{max width=\linewidth}}
\AfterEndEnvironment{tikzpicture}{\end{adjustbox}}
% légendes pgfplots lisibles par-dessus les courbes
\pgfplotsset{every axis/.append style={legend style={fill=white,fill opacity=0.92,text opacity=1,draw=black!15,inner sep=3pt}}}
% étiquettes d'arêtes (syntaxe edge["..."]) avec fond blanc
\tikzset{every edge quotes/.append style={fill=white,inner sep=1.2pt}}
%% FIGFIX-END
\begin{document}

\pagegarde{Lesson 14 — Grammar: Adverbs (for, since, ago, after, afterwards, etc.)}{Anglais}{Tle A}{Chapter 2 : Economic Life and Occupations}{EN14}{2 h}{Cette leçon de grammaire étudie les adverbes et expressions de temps for, since, ago, after et afterwards : leurs emplois, les temps verbaux qui les accompagnent et les pièges qu'ils tendent aux francophones. Observation guidée, règles illustrées d'exemples traduits et exercices d'application préparent efficacement à l'épreuve du Bac.
\vspace{4pt}
\begin{multicols}{2}\small
\begin{itemize}
\item À la fin de cette leçon, je sais distinguer for (durée) et since (point de départ) et les employer correctement.
\item À la fin de cette leçon, je sais utiliser ago avec le prétérit pour situer un événement dans le passé.
\item À la fin de cette leçon, je sais employer after comme préposition, conjonction et afterwards comme adverbe de conséquence temporelle.
\item À la fin de cette leçon, je sais choisir le bon temps verbal (present perfect, prétérit) selon l'adverbe utilisé.
\item À la fin de cette leçon, je sais éviter les erreurs fréquentes des francophones (*since 5 years, *I am here since..., *ago 3 days).
\item À la fin de cette leçon, je sais transformer des phrases en passant de for à since et inversement.
\end{itemize}\end{multicols}}

\begin{sommaire}
\begin{multicols}{2}
\begin{enumerate}
\item Observation guidée : découvrir les adverbes en contexte
\item For et Since : la durée et le point de départ
\item Ago : situer un événement dans le passé
\item After et Afterwards : la succession des événements
\item Comparaison français/anglais et pièges des francophones
\item Entraînement guidé : application et transformation
\item Activité d'intégration
\item Méthodes et exercices résolus
\item Exercices
\item Corrigés des exercices
\item Fiche bilan
\end{enumerate}
\end{multicols}
\end{sommaire}

\begin{objectifs}
\begin{itemize}
\item À la fin de cette leçon, je sais distinguer for (durée) et since (point de départ) et les employer correctement.
\item À la fin de cette leçon, je sais utiliser ago avec le prétérit pour situer un événement dans le passé.
\item À la fin de cette leçon, je sais employer after comme préposition, conjonction et afterwards comme adverbe de conséquence temporelle.
\item À la fin de cette leçon, je sais choisir le bon temps verbal (present perfect, prétérit) selon l'adverbe utilisé.
\item À la fin de cette leçon, je sais éviter les erreurs fréquentes des francophones (*since 5 years, *I am here since..., *ago 3 days).
\item À la fin de cette leçon, je sais transformer des phrases en passant de for à since et inversement.
\item À la fin de cette leçon, je sais poser et répondre à la question « How long...? » dans un dialogue sur la vie professionnelle.
\item À la fin de cette leçon, je sais rédiger un court paragraphe sur un métier ou une activité économique en utilisant ces adverbes.
\end{itemize}
\end{objectifs}

\begin{prerequis}
\begin{itemize}
\item Le prétérit simple (simple past) : formation et emploi (Première)
\item Le present perfect : formation (have/has + participe passé) et emploi (Première)
\item Les expressions de temps de base : yesterday, last week, in 2019 (Seconde/Première)
\item Le vocabulaire des métiers et de la vie économique (Chapter 2, leçons précédentes)
\item La distinction durée / point de départ en français (« depuis », « pendant », « il y a »)
\end{itemize}
\end{prerequis}

\newpage

\coursec{Observation guidée : découvrir les adverbes en contexte}

\courssub{Situation d'accroche : le piège du marché Mokolo}

Imaginez la scène : c'est samedi matin au marché Mokolo, à Yaoundé. Votre oncle, commerçant depuis des années, discute avec un client nigérian. Soudain, il se tourne vers vous et vous demande en anglais : \eng{How long have you been selling tomatoes?} — « Depuis combien de temps vends-tu des tomates ? ». Vous voulez répondre « depuis cinq ans »... mais faut-il dire \eng{for five years}, \eng{since five years} ou \eng{five years ago} ? Un instant de doute, et vous sentez que votre réponse peut tout changer.

Pendant ce temps, au Nigeria voisin, un ami vous écrit sur WhatsApp : \eng{I live here since 2020.} — une erreur typique des francophones, qui traduisent mot à mot le « depuis » français ! La phrase correcte serait : \eng{I have lived here since 2020.}

Ces petits mots — \cle{for}, \cle{since}, \cle{ago}, \cle{after}, \cle{afterwards} — sont partout dans la vie économique : durée d'un emploi, date d'ouverture d'une boutique, ancienneté dans un métier, succession d'événements dans la carrière d'un commerçant. Les maîtriser, c'est réussir les questions de grammaire du Bac \textbf{et} parler comme un vrai anglophone.

\textbf{Problématique de la leçon :} comment choisir correctement entre \eng{for}, \eng{since}, \eng{ago}, \eng{after} et \eng{afterwards}, et avec quel temps verbal les employer ?

\begin{attention}[L'objectif de cette première étape]
Avant d'apprendre la règle, nous allons \textbf{observer} : lire un texte authentique, repérer les adverbes, regarder les temps verbaux qui les accompagnent, et formuler nos propres hypothèses. C'est la démarche de l'approche par les compétences : \emph{observer} $\rightarrow$ \emph{dégager la règle} $\rightarrow$ \emph{s'entraîner}.
\end{attention}

\courssub{Lecture du texte support : Mrs Bih, trader at Mokolo Market}

Lisez le texte suivant à voix haute. Pendant la lecture, \textbf{soulignez} tous les adverbes de temps que vous rencontrez.

\begin{textebox}[Mrs Bih, a successful trader at Mokolo Market]
Mrs Bih has been a trader at Mokolo Market for fifteen years. She started her business in 2009, so she has worked there since 2009. Fifteen years ago, she was still a farmer in Bamenda, in the North-West Region. Every morning, she woke up early and went to her farm. After the harvest season, she moved to Yaoundé; afterwards, she opened her first small shop near the market.

Business was difficult at the beginning, but Mrs Bih never gave up. For many years, she sold only vegetables and fruit. Then, after five years of hard work, she saved enough money to buy a larger stall. Since that day, her business has grown steadily. She has employed three assistants for two years now, and she has been the president of the traders' association since last January.

Today, Mrs Bih is proud of her journey. “I arrived in Yaoundé fifteen years ago with almost nothing,” she says. “After all these years of effort, I can finally send my children to university.” Her story inspires many young traders at Mokolo Market.
\end{textebox}

\textbf{Glossaire (les mots difficiles) :}

\begin{center}
\begin{tabularx}{\linewidth}{|l|l|X|}
\hline
\rowcolor{popblueL} \textbf{Mot anglais} & \textbf{Nature} & \textbf{Traduction française} \\
\hline
trader & nom & commerçant(e) \\
\hline
harvest season & nom & saison des récoltes \\
\hline
to move & verbe & déménager, s'installer \\
\hline
stall & nom & étal, stand (au marché) \\
\hline
steadily & adverbe & régulièrement, progressivement \\
\hline
to give up & verbe & abandonner \\
\hline
journey & nom & parcours \\
\hline
almost & adverbe & presque \\
\hline
\end{tabularx}
\end{center}

\courssub{Repérage : souligner les adverbes, encadrer les verbes}

\begin{methode}[Comment observer une structure grammaticale en 4 étapes]
\begin{enumerate}
\item \textbf{Souligner l'adverbe} de temps (\eng{for}, \eng{since}, \eng{ago}, \eng{after}, \eng{afterwards}) et le groupe de mots qui l'accompagne.
\item \textbf{Encadrer le verbe} de la phrase et noter son temps : present perfect (\eng{has worked}), prétérit (\eng{moved}), etc.
\item \textbf{Se poser la question} : « durée ou point de départ ? » L'adverbe indique-t-il \emph{pendant combien de temps} (durée) ou \emph{à partir de quand / quand exactement} (point dans le temps) ?
\item \textbf{Traduire en français} et comparer : quel mot français correspond ? Le temps verbal est-il le même dans les deux langues ?
\end{enumerate}
\end{methode}

Appliquons cette méthode au texte. Voici les phrases clés, avec les adverbes en gras :

\begin{itemize}
\item \eng{Mrs Bih has been a trader \textbf{for fifteen years}.} — verbe : \eng{has been} (present perfect).
\item \eng{She has worked there \textbf{since 2009}.} — verbe : \eng{has worked} (present perfect).
\item \eng{Fifteen years \textbf{ago}, she was still a farmer.} — verbe : \eng{was} (prétérit).
\item \eng{\textbf{After} the harvest season, she moved to Yaoundé.} — verbe : \eng{moved} (prétérit).
\item \eng{\textbf{Afterwards}, she opened her first small shop.} — verbe : \eng{opened} (prétérit).
\item \eng{\textbf{Since} that day, her business has grown steadily.} — verbe : \eng{has grown} (present perfect).
\end{itemize}

\begin{exemplebox}[Phrases clés et leurs traductions]
\begin{itemize}
\item \eng{She has worked there since 2009.} = « Elle y travaille \textbf{depuis} 2009. »
\item \eng{She has been a trader for fifteen years.} = « Elle est commerçante \textbf{depuis} quinze ans. »
\item \eng{Fifteen years ago, she was still a farmer.} = « \textbf{Il y a} quinze ans, elle était encore agricultrice. »
\item \eng{After the harvest season, she moved to Yaoundé.} = « \textbf{Après} la saison des récoltes, elle s'est installée à Yaoundé. »
\item \eng{Afterwards, she opened her first small shop.} = « \textbf{Ensuite}, elle a ouvert sa première petite boutique. »
\end{itemize}
\textbf{Observation capitale :} le français utilise le même mot « depuis » pour \eng{for} et \eng{since}, mais l'anglais fait une distinction précise. Et remarquez : le français dit « elle \emph{travaille} » (présent) là où l'anglais dit \eng{she \emph{has worked}} (present perfect) !
\end{exemplebox}

\courssub{Observation des temps verbaux : un indice précieux}

Regardons maintenant le verbe qui accompagne chaque adverbe. Complétons ensemble ce tableau d'observation :

\begin{center}
\begin{tabularx}{\linewidth}{|l|X|l|l|}
\hline
\rowcolor{popblueL} \textbf{Adverbe} & \textbf{Exemple du texte} & \textbf{Temps verbal} & \textbf{Durée ou point ?} \\
\hline
for & has been a trader for fifteen years & present perfect & durée (15 ans) \\
\hline
since & has worked there since 2009 & present perfect & point de départ (2009) \\
\hline
ago & fifteen years ago, she was a farmer & prétérit & point dans le passé \\
\hline
after & after the harvest season, she moved & prétérit & point dans le passé \\
\hline
afterwards & afterwards, she opened a shop & prétérit & succession d'événements \\
\hline
\end{tabularx}
\end{center}

\begin{aretenir}[Premières constatations (à vérifier dans la suite de la leçon)]
\begin{itemize}
\item \eng{For} et \eng{since} s'accompagnent du \textbf{present perfect} : l'action a commencé dans le passé et continue aujourd'hui.
\item \eng{Ago}, \eng{after} et \eng{afterwards} s'accompagnent du \textbf{prétérit} : l'action est terminée, coupée du présent.
\item \eng{For} exprime une \textbf{durée} (quinze ans) ; \eng{since} exprime un \textbf{point de départ} (2009).
\end{itemize}
\end{aretenir}

\courssub{Classement : durée ou point de départ ?}

Classons maintenant les exemples en deux grands groupes. C'est la clé de toute la leçon.

\textbf{Groupe A — la durée et le lien avec le présent (present perfect) :}
\begin{itemize}
\item \eng{for fifteen years} → une période, une longueur de temps ;
\item \eng{since 2009} → le point de départ d'une action qui dure encore.
\end{itemize}

\textbf{Groupe B — les événements passés, terminés (prétérit) :}
\begin{itemize}
\item \eng{fifteen years ago} → un point situé en arrière par rapport à aujourd'hui ;
\item \eng{after the harvest season} → un événement qui en suit un autre ;
\item \eng{afterwards} → la suite, l'étape suivante du récit.
\end{itemize}

\begin{popfigure}
\begin{tikzpicture}[scale=1]
% Ligne du temps
\draw[thick, popdark, -{Stealth[length=3mm]}] (-1,0) -- (13,0);
\node[font=\small\bfseries, popdark] at (12.6,-0.45) {temps};
% Point 2009 (since)
\fill[popblue] (3,0) circle (0.12);
\node[below, font=\small, popblue] at (3,-0.15) {2009};
\node[above, font=\small\bfseries, popblue, align=center] at (3,0.9) {\eng{since 2009}\\ point de départ};
\draw[-{Stealth}, popblue, thick] (3,0.35) -- (3,0.75);
% Segment for 15 years
\draw[line width=2.5pt, popgreen] (3,0) -- (10.5,0);
\node[above, font=\small\bfseries, popgreen, align=center] at (6.75,0.15) {\eng{for fifteen years} : durée};
% Aujourd'hui
\fill[poporange] (10.5,0) circle (0.14);
\node[below, font=\small\bfseries, poporange] at (10.5,-0.2) {aujourd'hui};
% Point ago
\fill[poppurple] (6.75,-1.6) circle (0.12);
\node[font=\small, poppurple, align=center] at (6.75,-2.15) {\eng{fifteen years ago}\\ point dans le passé};
\draw[-{Stealth}, poppurple, thick] (10.5,-1.35) -- (7,-1.55);
\node[font=\small, poppurple] at (9,-1.05) {on recule de 15 ans};
\draw[dashed, poppurple] (10.5,0) -- (10.5,-1.35);
\draw[dashed, poppurple] (6.75,-1.6) -- (6.75,-0.35);
\end{tikzpicture}
\legende{Frise chronologique : \eng{since} marque le point de départ, \eng{for} mesure le segment de durée jusqu'à aujourd'hui, \eng{ago} recule dans le passé à partir d'aujourd'hui.}
\end{popfigure}

\courssub{Hypothèses des élèves : formuler la règle avant de l'apprendre}

Avant la leçon magistrale, essayez de compléter ces hypothèses avec votre voisin :

\begin{enumerate}
\item On utilise \eng{for} devant \ldots\ldots\ldots\ldots (une date ? une durée ?)
\item On utilise \eng{since} devant \ldots\ldots\ldots\ldots (une date ? une durée ?)
\item Après \eng{for} et \eng{since}, le verbe se met au \ldots\ldots\ldots\ldots
\item \eng{Ago} se place \ldots\ldots\ldots\ldots (avant ? après ?) le groupe de temps, et le verbe est au \ldots\ldots\ldots\ldots
\item \eng{Afterwards} signifie \ldots\ldots\ldots\ldots et raconte des événements au temps \ldots\ldots\ldots\ldots
\end{enumerate}

\begin{exoresolu}[Premier essai guidé]
\textbf{Énoncé.} Complétez avec \eng{for}, \eng{since} ou \eng{ago}, puis traduisez :
\begin{enumerate}
\item Mrs Bih has lived in Yaoundé \ldots\ldots\ldots fifteen years.
\item She left Bamenda fifteen years \ldots\ldots\ldots
\item She has been the president of the association \ldots\ldots\ldots last January.
\end{enumerate}
\tcblower
\textbf{Solution détaillée.}
\begin{enumerate}
\item \eng{Mrs Bih has lived in Yaoundé \textbf{for} fifteen years.} — « fifteen years » est une \textbf{durée} (on peut la compter : 1 an, 2 ans\ldots) → \eng{for}. Traduction : « Mme Bih habite à Yaoundé depuis quinze ans. »
\item \eng{She left Bamenda fifteen years \textbf{ago}.} — le verbe \eng{left} est au \textbf{prétérit} et l'événement est terminé → \eng{ago}, placé \textbf{après} le groupe de temps. Traduction : « Elle a quitté Bamenda il y a quinze ans. »
\item \eng{She has been the president \textbf{since} last January.} — « last January » est un \textbf{point de départ} (une date précise) → \eng{since}. Traduction : « Elle est présidente de l'association depuis janvier dernier. »
\end{enumerate}
\textbf{Astuce :} demandez-vous « pendant combien de temps ? » (→ \eng{for}) ou « à partir de quand ? » (→ \eng{since}).
\end{exoresolu}

\begin{attention}[Piège n°1 des francophones : le faux ami « depuis »]
En français, un seul mot — « depuis » — correspond à \textbf{deux} mots anglais différents :
\begin{itemize}
\item « depuis \textbf{cinq ans} » (durée) = \eng{for five years} ;
\item « depuis \textbf{2019} » (point de départ) = \eng{since 2019}.
\end{itemize}
Et attention au temps verbal : le français utilise le \textbf{présent} (« elle travaille depuis 2009 »), l'anglais exige le \textbf{present perfect} : \eng{She has worked there since 2009.} Ne dites jamais \eng{She works there since 2009.} !
\end{attention}

\begin{savaistu}[Ces adverbes dans la vie quotidienne au Cameroun anglophone]
Dans les régions du Nord-Ouest et du Sud-Ouest (Bamenda, Buea, Limbe), ces adverbes sont omniprésents dans les conversations quotidiennes sur le travail et le commerce : \eng{I have been a teacher here for twenty years.} (« Je suis enseignant ici depuis vingt ans. »), \eng{We have known each other since primary school.} (« Nous nous connaissons depuis l'école primaire. »), \eng{The market opened two years ago.} (« Le marché a ouvert il y a deux ans. »). Au Bac, les sujets de grammaire testent presque chaque année la distinction \eng{for}/\eng{since}/\eng{ago} : la maîtriser, c'est gagner des points faciles !
\end{savaistu}

\begin{exercice}[À vous de jouer : repérage dans le texte]
Reprenez le texte « Mrs Bih » et répondez :
\begin{enumerate}
\item Trouvez deux autres phrases contenant \eng{for} ou \eng{since} et notez le temps du verbe.
\item Pourquoi dit-on \eng{she moved to Yaoundé} (prétérit) et non \eng{she has moved} ?
\item Traduisez : \eng{After five years of hard work, she saved enough money.}
\end{enumerate}
\end{exercice}

\begin{corrige}[Exercice « À vous de jouer »]
\begin{enumerate}
\item \eng{She has employed three assistants for two years now.} (present perfect + \eng{for} + durée) ; \eng{She has been the president of the traders' association since last January.} (present perfect + \eng{since} + point de départ).
\item Parce que le déménagement est un événement \textbf{ponctuel et terminé}, situé à un moment précis du passé (\eng{after the harvest season}) : l'anglais utilise alors le prétérit, sans lien avec le présent.
\item « Après cinq années de travail acharné, elle a économisé assez d'argent. »
\end{enumerate}
\end{corrige}

Nos observations sont terminées : nous avons repéré les adverbes, noté les temps verbaux, classé durée et point de départ, et formulé nos hypothèses. Passons maintenant à la règle précise : \eng{for} et \eng{since}, la durée et le point de départ.

\coursec{For et Since : la durée et le point de départ}

Dans la section précédente, nous avons observé en contexte plusieurs adverbes et expressions de temps : \eng{for}, \eng{since}, \eng{ago}, \eng{after}, \eng{afterwards}. Nous avons notamment relevé ces deux phrases, tirées de notre texte d'observation sur la vie économique : \eng{Mr Etoundi has worked at the cocoa cooperative for fifteen years.} et \eng{The cooperative has exported cocoa to Europe since 2010.} Pourquoi l'anglais utilise-t-il deux mots différents, \eng{for} et \eng{since}, là où le français n'emploie qu'un seul mot, « depuis » ? C'est toute la question de cette section.

\courssub{Observer avant de comprendre}

\begin{textebox}[At the job centre in Yaoundé]
\textbf{Interviewer:} How long have you been looking for a job, Mr Nkeng?\par
\textbf{Mr Nkeng:} I have been unemployed for eight months. I lost my job at the bank in January.\par
\textbf{Interviewer:} So you have been unemployed since January. And before that?\par
\textbf{Mr Nkeng:} I worked at the bank for twelve years. I have known the manager since I was a student.\par
\textbf{Interviewer:} I see. Have you attended any training courses recently?\par
\textbf{Mr Nkeng:} Yes. I have attended three computer courses since March, and I have improved my English a lot. I have studied English for many years, actually.\par
\textbf{Interviewer:} Good. We have had your file on our desk for two weeks now. We will contact you soon.
\end{textebox}

\textbf{Glossaire :} \emph{to look for a job} (v.) : chercher un emploi ; \emph{unemployed} (adj.) : au chômage ; \emph{a training course} (n.) : un stage de formation ; \emph{a file} (n.) : un dossier.

\textbf{Questions d'observation :}
\begin{enumerate}
\item Relevez toutes les expressions introduites par \eng{for}. Que désigne le mot qui suit : une durée ou une date ?
\item Relevez toutes les expressions introduites par \eng{since}. Que désigne le mot qui suit : une durée ou un point de départ ?
\item Quel temps verbal utilise l'anglais dans toutes ces phrases ?
\end{enumerate}

\textbf{Réponses attendues :} \eng{for} est suivi de durées (\eng{for eight months}, \eng{for twelve years}, \eng{for many years}, \eng{for two weeks}) ; \eng{since} est suivi de points de départ (\eng{since January}, \eng{since I was a student}, \eng{since March}) ; le temps verbal est le present perfect (ou le present perfect continu).

\courssub{La règle : FOR + durée, SINCE + point de départ}

\begin{propriete}[FOR et SINCE : la distinction fondamentale]
\begin{itemize}
\item \cle{For} est suivi d'une \textbf{durée}, c'est-à-dire d'une période mesurable : \eng{for two years}, \eng{for a long time}, \eng{for six months}, \eng{for three weeks}. Il répond à la question « \textbf{pendant combien de temps ?} » (\eng{how long?}).
\item \cle{Since} est suivi d'un \textbf{point de départ}, c'est-à-dire d'une date, d'un moment précis ou d'un événement : \eng{since 2015}, \eng{since Monday}, \eng{since January}, \eng{since I left school}. Il répond à la question « \textbf{depuis quand ?} » (\eng{since when?}).
\item Les deux s'emploient généralement avec le \textbf{present perfect} (\eng{have/has + participe passé}) ou le \textbf{present perfect continu} (\eng{have/has been + V-ing}) quand l'action a commencé dans le passé et \textbf{continue encore au moment présent}.
\end{itemize}
\end{propriete}

\begin{exemplebox}[Exemples traduits]
\begin{itemize}
\item \eng{I have known him for ten years.} = Je le connais depuis dix ans.
\item \eng{They have lived in Buea since 2018.} = Ils habitent à Buéa depuis 2018.
\item \eng{He has been unemployed since he lost his job.} = Il est au chômage depuis qu'il a perdu son emploi.
\item \eng{We have been waiting for the bus for forty minutes.} = Nous attendons le bus depuis quarante minutes.
\item \eng{She has worked at the ministry since last September.} = Elle travaille au ministère depuis septembre dernier.
\item \eng{My father has had this shop for a long time.} = Mon père a cette boutique depuis longtemps.
\end{itemize}
\end{exemplebox}

\begin{popfigure}
\begin{center}
\begin{tikzpicture}[scale=1, every node/.style={font=\small}]
% Axe du temps
\draw[-{Stealth}, thick, popdark] (0,0) -- (11,0) node[right]{temps};
\draw[fill=popdark] (9.5,0) circle (0.08);
\node[above, popdark] at (9.5,0.1) {\textbf{now}};
% FOR : segment de durée
\draw[-{Stealth}, very thick, popblue] (2,0.9) -- (9.5,0.9);
\draw[fill=popblue] (2,0.9) circle (0.08);
\draw[fill=popblue] (9.5,0.9) circle (0.08);
\node[popblue, align=center] at (5.75,1.7) {\textbf{FOR + durée}\\ \eng{for ten years} (une longueur mesurable)};
\draw[popblue, dashed] (2,0.9) -- (2,0);
\draw[popblue, dashed] (9.5,0.9) -- (9.5,0);
% SINCE : point d'origine
\draw[-{Stealth}, very thick, poporange] (2,-0.9) -- (9.5,-0.9);
\node[star, star points=5, fill=poporange, inner sep=1.5pt] at (2,-0.9) {};
\node[poporange, align=center] at (5.75,-1.7) {\textbf{SINCE + point de départ}\\ \eng{since 2014} (un point précis dans le temps)};
\draw[poporange, dashed] (2,-0.9) -- (2,0);
\draw[poporange, dashed] (9.5,-0.9) -- (9.5,0);
\end{tikzpicture}
\end{center}
\legende{\eng{For} mesure une longueur de temps (un segment) ; \eng{since} désigne le point de départ (un point) à partir duquel on compte jusqu'à aujourd'hui.}
\end{popfigure}

\courssub{Le tableau des durées et des points de départ}

Le critère est simple : si l'on peut répondre à la question « combien de temps cela dure-t-il ? » par l'expression, c'est une durée, donc \eng{for}. Si l'expression désigne une date ou un moment précis du calendrier, c'est un point de départ, donc \eng{since}.

\begin{center}
\begin{tabularx}{\linewidth}{|X|X|}
\hline
\rowcolor{popblueL} \textbf{FOR + durée} & \textbf{SINCE + point de départ} \\
\hline
\eng{for two hours} (pendant deux heures) & \eng{since 8 o'clock} (depuis huit heures) \\
\hline
\eng{for three days} (pendant trois jours) & \eng{since Monday} (depuis lundi) \\
\hline
\eng{for six months} (pendant six mois) & \eng{since January} (depuis janvier) \\
\hline
\eng{for ten years} (pendant dix ans) & \eng{since 2015} (depuis 2015) \\
\hline
\eng{for a long time} (depuis longtemps) & \eng{since last week} (depuis la semaine dernière) \\
\hline
\eng{for ages} (depuis une éternité) & \eng{since Christmas} (depuis Noël) \\
\hline
\eng{for a week} (pendant une semaine) & \eng{since I was a child} (depuis mon enfance) \\
\hline
\end{tabularx}
\end{center}

\begin{attention}[L'erreur n° 1 des francophones : *since five years]
En français, un seul mot, « depuis », recouvre les deux notions : « depuis cinq ans » (durée) et « depuis 2019 » (point de départ). L'anglais, lui, \textbf{distingue obligatoirement} les deux cas. On ne dit JAMAIS *\eng{since five years}. On dit :
\begin{itemize}
\item \eng{I have lived in Douala for five years.} (durée)
\item \eng{I have lived in Douala since 2019.} (point de départ)
\end{itemize}
Avant de parler, demandez-vous : est-ce une durée ou une date ?
\end{attention}

\courssub{Le temps associé : le present perfect}

\begin{propriete}[For / since + present perfect]
Quand une action ou un état a commencé dans le passé et \textbf{se poursuit encore maintenant}, l'anglais utilise le \textbf{present perfect} (\eng{have/has + participe passé}) ou, pour insister sur la continuité de l'action, le \textbf{present perfect continu} (\eng{have/has been + V-ing}) :
\begin{itemize}
\item \eng{She has worked here for six months.} = Elle travaille ici depuis six mois.
\item \eng{She has been working here since March.} = Elle travaille ici depuis mars.
\end{itemize}
Le français, lui, utilise le \textbf{présent} dans ce cas : « elle travaille ici depuis six mois ». C'est une différence majeure entre les deux langues.
\end{propriete}

\begin{attention}[Ne jamais utiliser le présent simple]
Le français dit : « J'habite ici depuis 2020 » (présent). L'anglais interdit le présent simple avec \eng{for} et \eng{since} :
\begin{itemize}
\item FAUX : *\eng{I live here since 2020.}
\item CORRECT : \eng{I have lived here since 2020.}
\item FAUX : *\eng{He works in this factory for ten years.}
\item CORRECT : \eng{He has worked in this factory for ten years.}
\end{itemize}
Retenez : « depuis » + présent en français = \eng{for/since} + \textbf{present perfect} en anglais.
\end{attention}

\begin{aretenir}[La question-clé : How long...?]
Pour demander « depuis combien de temps ? », l'anglais utilise \eng{How long have you...?} au present perfect. La réponse se construit avec \eng{for} ou \eng{since} selon la nature de la réponse :
\begin{itemize}
\item \eng{How long have you been a teacher?} — \eng{I have been a teacher for fifteen years.} / \eng{I have been a teacher since 2009.}
\item \eng{How long have they lived in Bamenda?} — \eng{They have lived there for a long time.} / \eng{They have lived there since their marriage.}
\end{itemize}
\end{aretenir}

\courssub{Cas particuliers : for avec le prétérit, since + subordonnée}

\begin{propriete}[FOR + prétérit : une durée terminée]
\eng{For} peut aussi s'employer avec le \textbf{prétérit simple} quand la durée appartient à un passé \textbf{révolu}, sans lien avec le présent :
\begin{itemize}
\item \eng{He worked in Douala for ten years.} = Il a travaillé à Douala pendant dix ans. (\emph{il n'y travaille plus})
\item \eng{He has worked in Douala for ten years.} = Il travaille à Douala depuis dix ans. (\emph{il y travaille encore})
\end{itemize}
Le choix du temps verbal change le sens : prétérit = c'est fini ; present perfect = cela continue. Attention : \eng{since} ne s'emploie pratiquement jamais avec le prétérit simple dans la proposition principale.
\end{propriete}

\begin{propriete}[SINCE + subordonnée]
\eng{Since} peut introduire non seulement un nom (une date), mais aussi une \textbf{proposition subordonnée} complète (sujet + verbe). Le verbe de la subordonnée est alors généralement au \textbf{prétérit} :
\begin{itemize}
\item \eng{He has been unemployed since he lost his job.} = Il est au chômage depuis qu'il a perdu son emploi.
\item \eng{She has been interested in computers since she was a child.} = Elle s'intéresse à l'informatique depuis qu'elle est enfant.
\item \eng{We have been friends since we met at secondary school.} = Nous sommes amis depuis que nous nous sommes rencontrés au collège.
\end{itemize}
Structure : \textbf{present perfect + since + sujet + prétérit}.
\end{propriete}

\begin{methode}[Comment choisir entre FOR et SINCE]
\begin{enumerate}
\item Repérez l'expression de temps qui suit le blanc.
\item Posez-vous la question : est-ce une \textbf{durée} (une longueur : \eng{two years}, \eng{a long time}, \eng{six months}) ou un \textbf{point de départ} (une date, un moment : \eng{2020}, \eng{Monday}, \eng{last June}) ?
\item Durée $\rightarrow$ \eng{for}. Point de départ $\rightarrow$ \eng{since}.
\item Vérifiez le temps du verbe : si l'action continue encore, utilisez le present perfect (\eng{have/has + participe passé}).
\item Astuce de vérification : remplacez mentalement par une date. Si « depuis 2020 » fonctionne en français, c'est \eng{since} ; si « pendant cinq ans » fonctionne, c'est \eng{for}.
\end{enumerate}
\end{methode}

\courssub{La transformation for / since}

Une même réalité peut s'exprimer des deux façons, en convertissant la durée en point de départ (et inversement). Il suffit de faire le calcul à partir du moment présent :

\begin{center}
\begin{tabularx}{\linewidth}{|X|X|}
\hline
\rowcolor{popblueL} \textbf{FOR + durée} & \textbf{SINCE + point de départ} \\
\hline
\eng{I have worked here for 5 years.} & \eng{I have worked here since 2019.} \\
\hline
\eng{She has been ill for a week.} & \eng{She has been ill since last Tuesday.} \\
\hline
\eng{They have been married for 20 years.} & \eng{They have been married since 2004.} \\
\hline
\eng{We have known each other for ages.} & \eng{We have known each other since primary school.} \\
\hline
\end{tabularx}
\end{center}

\begin{exoresolu}[Exercice de transformation et de choix]
\textbf{Énoncé.} 1) Complétez avec \eng{for} ou \eng{since} : (a) \eng{My uncle has been a taxi driver ... twenty years.} (b) \eng{The shop has been closed ... last Friday.} (c) \eng{I have not seen her ... she left Yaoundé.} (d) \eng{They have played football ... two hours.} 2) Transformez avec \eng{since} : \eng{He has lived in Bafoussam for eight years.} (nous sommes en 2024).
\tcblower
\textbf{Solution détaillée.}
1) (a) \eng{for twenty years} : « vingt ans » est une durée mesurable $\rightarrow$ \eng{for}. (b) \eng{since last Friday} : « vendredi dernier » est un point de départ précis $\rightarrow$ \eng{since}. (c) \eng{since she left Yaoundé} : \eng{since} introduit ici une subordonnée (sujet \eng{she} + prétérit \eng{left}). (d) \eng{for two hours} : « deux heures » est une durée $\rightarrow$ \eng{for}.
2) 2024 moins 8 ans = 2016. La durée « huit ans » devient le point de départ « 2016 » : \eng{He has lived in Bafoussam since 2016.} Le verbe reste au present perfect car il y habite toujours.
\end{exoresolu}

\begin{exercice}[À vous de jouer]
Complétez avec \eng{for} ou \eng{since}, puis mettez le verbe au present perfect :
\begin{enumerate}
\item \eng{The students (be) in the classroom ... 7 a.m.}
\item \eng{Mrs Abena (sell) vegetables at the market ... fifteen years.}
\item \eng{I (not / hear) from my brother ... he travelled to Lagos.}
\item \eng{We (study) English ... six years.}
\item \eng{The factory (produce) soap ... 1998.}
\end{enumerate}
\end{exercice}

\begin{corrige}[Exercice « À vous de jouer »]
\begin{enumerate}
\item \eng{The students have been in the classroom since 7 a.m.} (point de départ : une heure précise).
\item \eng{Mrs Abena has sold vegetables at the market for fifteen years.} (durée).
\item \eng{I have not heard from my brother since he travelled to Lagos.} (\eng{since} + subordonnée au prétérit).
\item \eng{We have studied English for six years.} (durée).
\item \eng{The factory has produced soap since 1998.} (point de départ : une année).
\end{enumerate}
\end{corrige}

\begin{savaistu}[For et since dans la vie quotidienne anglophone]
Dans les pays anglophones, la question \eng{How long have you been here?} est l'une des premières questions qu'on pose à un nouveau collègue, à un voisin ou à un étudiant étranger. Savoir répondre correctement — \eng{I have been here for two months} ou \eng{I arrived in September, so I have been here since then} — est une compétence sociale essentielle. Au Cameroun anglophone, à Buea ou à Bamenda, vous l'entendrez chaque jour dans les administrations et les entreprises.
\end{savaistu}

\coursec{Ago : situer un événement dans le passé}

Dans la section précédente, nous avons étudié \cle{for} et \cle{since}, qui expriment une durée ou un point de départ reliés au présent, et qui s'emploient avec le \emph{present perfect}. Nous abordons maintenant un adverbe qui, lui, coupe totalement le lien avec le présent : \cle{ago}. Avec \eng{ago}, l'événement est terminé, daté, rangé dans le passé. C'est l'équivalent anglais du français « il y a » — mais attention : l'anglais ne place pas les mots dans le même ordre que le français, et c'est précisément là que se cachent les erreurs les plus fréquentes des candidats francophones au Bac.

\courssub{Observation guidée : un texte pour découvrir}

Lisons d'abord ce court texte biographique, genre très fréquent dans les épreuves de compréhension de l'écrit en Terminale.

\begin{textebox}[The story of Mama Ngo's bakery — an original text]
Mama Ngo opened her small bakery in Buea twenty-five years ago. At that time, she had only one oven and two employees. She bought her first delivery bicycle ten years ago, and her business started to grow very quickly. Five years ago, she moved to a bigger building near the market, and three years ago, her son joined the company after his studies in Douala. A few months ago, the bakery won a regional prize for the best bread in the South-West Region. Today, Mama Ngo employs fifteen people, but she still remembers the day, a long time ago, when she sold her first loaf of bread.
\end{textebox}

\textbf{Glossaire :}

\begin{center}
\begin{tabularx}{\linewidth}{|l|l|X|}\hline
\rowcolor{popblueL} \textbf{Mot anglais} & \textbf{Nature} & \textbf{Traduction} \\ \hline
bakery & nom & boulangerie \\ \hline
oven & nom & four \\ \hline
employee & nom & employé(e) \\ \hline
to grow (grew, grown) & verbe & grandir, se développer \\ \hline
prize & nom & prix, récompense \\ \hline
loaf (pl. loaves) & nom & pain, miche \\ \hline
\end{tabularx}
\end{center}

\textbf{Questions d'observation :}

\begin{enumerate}
\item Souligne dans le texte toutes les occurrences de \eng{ago}. Quel mot ou groupe de mots se trouve toujours \emph{avant} \eng{ago} ?
\item Quel temps verbal est utilisé dans toutes les phrases contenant \eng{ago} ?
\item Traduis en français : \eng{She opened her bakery twenty-five years ago.} Où se place « il y a » dans la phrase française ?
\end{enumerate}

\textbf{Réponses attendues :} (1) une durée (\eng{twenty-five years}, \eng{ten years}, \eng{a few months}, \eng{a long time}) ; (2) le prétérit simple ; (3) « Elle a ouvert sa boulangerie il y a vingt-cinq ans » — en français, « il y a » se place \emph{avant} la durée, alors qu'en anglais \eng{ago} se place \emph{après}.

\courssub{La règle : forme et emploi}

\begin{propriete}[Structure de base : durée + ago]
\cle{Ago} est un adverbe qui signifie « il y a ». Il se place \textbf{toujours après} l'expression de durée, et le verbe de la phrase est \textbf{toujours au prétérit simple} (\emph{simple past}).

\begin{center}
\textbf{Sujet + verbe au prétérit + complément + durée + ago}
\end{center}

\eng{She opened her shop two years ago.} « Elle a ouvert sa boutique il y a deux ans. »

La durée peut être :
\begin{itemize}
\item un nombre + unité de temps : \eng{two days ago}, \eng{three weeks ago}, \eng{five years ago} ;
\item un article + unité : \eng{a minute ago}, \eng{a century ago} ;
\item une expression de durée : \eng{a long time ago} (il y a longtemps), \eng{a short while ago} (il y a peu de temps), \eng{ages ago} (il y a une éternité).
\end{itemize}
\end{propriete}

\begin{exemplebox}[Exemples traduits]
\eng{My father retired five years ago.} « Mon père a pris sa retraite il y a cinq ans. »

\eng{The company was created a century ago.} « L'entreprise a été créée il y a un siècle. »

\eng{I met the manager a few minutes ago.} « J'ai rencontré le directeur il y a quelques minutes. »

\eng{They moved to Bamenda a long time ago.} « Ils se sont installés à Bamenda il y a longtemps. »

\eng{The train left ten minutes ago.} « Le train est parti il y a dix minutes. »
\end{exemplebox}

\begin{attention}[L'ordre des mots : le piège numéro un]
En français, « il y a » précède la durée ; en anglais, \eng{ago} la suit. Ne calque jamais l'ordre français !

\begin{center}
\begin{tabularx}{\linewidth}{|X|X|}\hline
\rowcolor{popblueL} \textbf{Français} & \textbf{English} \\ \hline
il y a trois jours & three days ago \\ \hline
\emph{*ago three days} & FAUX — ordre impossible \\ \hline
il y a deux ans & two years ago \\ \hline
il y a longtemps & a long time ago \\ \hline
\end{tabularx}
\end{center}
\end{attention}

\begin{attention}[Jamais de present perfect avec ago]
\eng{Ago} situe un événement \emph{terminé} à une date précise du passé : il exige le \textbf{prétérit simple}. Le present perfect est impossible.

\emph{*I have seen him two days ago.} → \eng{I saw him two days ago.} « Je l'ai vu il y a deux jours. »

\emph{*She has opened her shop two years ago.} → \eng{She opened her shop two years ago.}
\end{attention}

\begin{propriete}[La question : How long ago ... ?]
Pour demander « il y a combien de temps ? », l'anglais utilise \eng{How long ago} + auxiliaire \eng{did} + sujet + base verbale :

\eng{How long ago did she open her bakery?} « Il y a combien de temps a-t-elle ouvert sa boulangerie ? »

\eng{How long ago did they leave Yaoundé?} « Il y a combien de temps ont-ils quitté Yaoundé ? »

Attention : avec \eng{did}, le verbe revient à la base verbale — jamais \emph{*How long ago did she opened...?}
\end{propriete}

\courssub{Visualiser : la frise chronologique}

\begin{popfigure}
\begin{center}
\begin{tikzpicture}[scale=1]
\draw[-{Stealth[length=3mm]}, very thick, popdark] (0,0) -- (10.5,0);
\fill[popgreen] (9,0) circle (0.12);
\node[above, popgreen, font=\bfseries] at (9,0.15) {NOW};
\node[below, popmuted] at (9,-0.15) {today};
\fill[poporange] (2.5,0) circle (0.12);
\node[below, align=center] at (2.5,-0.15) {past event\\ \emph{finished}};
\draw[-{Stealth[length=2.5mm]}, very thick, popblue] (8.7,0.8) -- (2.8,0.8);
\node[above, popblue, font=\bfseries] at (5.75,0.85) {3 years ago};
\node[align=center, popdark] at (5.75,-1.3) {\eng{He started his business three years ago.}\\ (simple past — event completed)};
\end{tikzpicture}
\end{center}
\legende{\eng{Ago} compte à rebours à partir d'aujourd'hui (NOW) vers un événement passé terminé : le verbe est au prétérit.}
\end{popfigure}

\courssub{Ago ou before ? Le complément utile pour le Bac}

\cle{Ago} compte la durée \textbf{à partir d'aujourd'hui}, du moment où l'on parle. Mais dans un récit au passé, quand on compte à partir \textbf{d'un autre moment passé}, l'anglais utilise \cle{before} (ou \cle{earlier}), souvent avec le \emph{past perfect}.

\begin{center}
\begin{tabularx}{\linewidth}{|X|X|X|}\hline
\rowcolor{popblueL} \textbf{Situation} & \textbf{English} & \textbf{Français} \\ \hline
Compte à partir d'aujourd'hui & \eng{I met her two years ago.} & Je l'ai rencontrée il y a deux ans. \\ \hline
Compte à partir d'un moment passé & \eng{He said he had met her two years before.} & Il a dit qu'il l'avait rencontrée deux ans auparavant. \\ \hline
\end{tabularx}
\end{center}

\begin{savaistu}[Ago dans le discours rapporté : un classique du Bac]
Au discours indirect, \eng{ago} devient \eng{before} (ou \eng{earlier}), et le prétérit recule au past perfect :

\eng{“I started this job three years ago,” he said.} → \eng{He said he had started that job three years before.}

Cette transformation (\eng{ago} → \eng{before}, simple past → past perfect) est très souvent testée dans les exercices de \emph{reported speech} au baccalauréat.
\end{savaistu}

\courssub{Ago versus for : ne pas confondre}

C'est la comparaison décisive de cette leçon. Les deux expressions contiennent une durée, mais elles ne disent pas la même chose et n'emploient pas le même temps.

\begin{center}
\begin{tabularx}{\linewidth}{|X|X|X|}\hline
\rowcolor{popblueL} \textbf{Critère} & \textbf{ago} & \textbf{for} \\ \hline
Sens & il y a (événement daté, terminé) & pendant/depuis (durée qui continue) \\ \hline
Temps verbal & simple past & present perfect \\ \hline
Position & après la durée & avant la durée \\ \hline
Exemple & \eng{He started two years ago.} & \eng{He has worked here for two years.} \\ \hline
Traduction & Il a commencé il y a deux ans. & Il travaille ici depuis deux ans. \\ \hline
\end{tabularx}
\end{center}

\eng{He started two years ago} : l'action de commencer est terminée, datée. \eng{He has worked here for two years} : il travaille ici \emph{encore aujourd'hui}. Les deux phrases sont vraies en même temps — elles regardent simplement la réalité sous deux angles différents.

\courssub{Entraînement guidé}

\begin{exoresolu}[Mettre le verbe au temps correct]
Complète avec le prétérit simple ou le present perfect du verbe entre parenthèses.

1. The workers (build) \dots\ this market ten years ago.

2. She (work) \dots\ as a nurse for fifteen years, and she loves her job.

3. I (see) \dots\ the accountant a few minutes ago in his office.

4. They (live) \dots\ in Douala since 2015.

5. My uncle (buy) \dots\ his first taxi a long time ago.
\tcblower
\textbf{Corrigé détaillé :}

1. \eng{The workers built this market ten years ago.} — \eng{ago} impose le prétérit simple ; \eng{build} est irrégulier : build–built–built.

2. \eng{She has worked as a nurse for fifteen years, and she loves her job.} — \eng{for} + durée, action qui continue : present perfect.

3. \eng{I saw the accountant a few minutes ago in his office.} — \eng{ago} → prétérit ; \eng{see} est irrégulier : see–saw–seen.

4. \eng{They have lived in Douala since 2015.} — \eng{since} + point de départ : present perfect.

5. \eng{My uncle bought his first taxi a long time ago.} — \eng{ago} → prétérit ; \eng{buy} est irrégulier : buy–bought–bought.
\end{exoresolu}

\begin{exercice}[À toi de jouer]
1. Traduis en anglais : « Mon grand-père a fondé cette entreprise il y a quarante ans. »

2. Corrige la phrase fautive : \emph{*Ago two weeks, I have visited the factory.}

3. Transforme au discours indirect : \eng{“I sold my shop five years ago,” the trader said.}

4. Choisis la bonne option : \eng{The bank opened / has opened a new branch three months ago.}

5. Pose la question correspondant à la réponse : \eng{She arrived in Buea six months ago.} (interroge sur la durée)
\end{exercice}

\begin{corrige}[Exercice : À toi de jouer]
1. \eng{My grandfather founded this company forty years ago.} (durée + \eng{ago}, verbe au prétérit).

2. \eng{Two weeks ago, I visited the factory.} — deux erreurs : l'ordre (\eng{two weeks ago}, jamais \emph{*ago two weeks}) et le temps (prétérit \eng{visited}, jamais de present perfect avec \eng{ago}).

3. \eng{The trader said he had sold his shop five years before.} — \eng{ago} devient \eng{before}, le prétérit recule au past perfect.

4. \eng{The bank opened a new branch three months ago.} — \eng{ago} exige le prétérit simple.

5. \eng{How long ago did she arrive in Buea?} — \eng{How long ago} + \eng{did} + base verbale \eng{arrive}.
\end{corrige}

\begin{aretenir}[L'essentiel sur ago]
\begin{itemize}
\item \eng{Ago} = « il y a » ; il se place \textbf{après} la durée : \eng{three days ago}.
\item Le verbe est \textbf{toujours au prétérit simple} : jamais de present perfect avec \eng{ago}.
\item La question se construit avec \eng{How long ago did ... ?} + base verbale.
\item \eng{Ago} compte à partir d'aujourd'hui ; dans un récit au passé, on utilise \eng{before} avec le past perfect.
\item Opposition clé : \eng{He started two years ago} (terminé) / \eng{He has worked here for two years} (continue).
\end{itemize}
\end{aretenir}

\coursec{After et Afterwards : la succession des événements}

Dans la section précédente, nous avons vu que \cle{ago} permet de situer un événement unique dans le passé, en remontant le temps depuis le présent (\eng{two years ago} = il y a deux ans). Mais dans la vie économique comme dans la vie quotidienne, les événements ne sont pas isolés : ils se \emph{succèdent}. Un commerçant vend d'abord ses marchandises, \emph{ensuite} il rentre chez lui. Une jeune fille termine ses études, \emph{après quoi} elle cherche un emploi. Pour raconter une chaîne d'actions dans l'ordre, l'anglais dispose de deux mots que les francophones confondent souvent : \cle{after} et \cle{afterwards}. C'est l'objet de cette section.

\courssub{Observation guidée : découvrir after et afterwards en contexte}

Lisons d'abord ce court texte sur la journée d'une commerçante de Douala. Repérez les mots \eng{after} et \eng{afterwards} et observez ce qui les suit.

\begin{textebox}[A day in the life of Mrs Ngo Bell, trader in Douala]
Mrs Ngo Bell wakes up at five o'clock every morning. After a quick breakfast, she takes a taxi to the Mboppi market, where she sells tomatoes, onions and plantains. Business is usually good in the morning. After she has sold most of her goods, she counts her money carefully and puts it in her bag. After lunch, she visits her supplier in Bonabéri to negotiate prices for the next day. Afterwards, she goes back to her stall and works until six o'clock.

In the evening, after closing her stall, she meets the other members of her tontine. They discuss their savings and their projects. Afterwards, Mrs Ngo Bell returns home, cooks dinner for her family and helps her children with their homework. “I started with almost nothing,” she says. “But after I had saved enough money, I was able to rent my own stall. Afterwards, everything became easier.”
\end{textebox}

\textbf{Glossaire :}

\begin{center}
\begin{tabularx}{\linewidth}{|l|l|X|}
\hline
\rowcolor{popblueL} Mot anglais & Nature & Traduction française \\
\hline
trader & nom & commerçante \\
\hline
stall & nom & étal, stand (au marché) \\
\hline
goods & nom (pluriel) & marchandises \\
\hline
supplier & nom & fournisseur \\
\hline
to negotiate & verbe & négocier \\
\hline
tontine & nom & tontine (association d'épargne) \\
\hline
savings & nom (pluriel) & économies, épargne \\
\hline
to rent & verbe & louer \\
\hline
\end{tabularx}
\end{center}

\textbf{Questions d'observation :}

\begin{enumerate}
\item Combien de fois le mot \eng{after} apparaît-il ? Qu'est-ce qui le suit à chaque fois : un nom, un verbe en \eng{-ing}, ou une proposition complète (sujet + verbe) ?
\item Combien de fois le mot \eng{afterwards} apparaît-il ? Est-il suivi d'un nom ou d'une proposition ?
\item Dans la phrase \eng{After she has sold most of her goods, she counts her money}, quelle action se produit \emph{en premier} ?
\end{enumerate}

\textbf{Réponses attendues :} \eng{after} est toujours suivi de quelque chose (un nom : \eng{after a quick breakfast}, \eng{after lunch} ; une proposition : \eng{after she has sold...}, \eng{after I had saved...} ; un gérondif : \eng{after closing her stall}). En revanche, \eng{afterwards} est toujours seul : il termine la phrase ou est suivi d'une virgule. Enfin, la proposition introduite par \eng{after} exprime toujours l'action \emph{antérieure} (celle qui se produit en premier).

\courssub{After : préposition et conjonction}

\begin{propriete}[After est TOUJOURS suivi d'un complément]
Le mot \cle{after} ne s'emploie \textbf{jamais seul}. Il a deux natures grammaticales possibles :

\begin{enumerate}
\item \textbf{After, préposition} : suivi d'un \textbf{nom}, d'un \textbf{pronom} ou d'un \textbf{gérondif} (verbe en \eng{-ing}).

\eng{after work} (après le travail) ; \eng{after the meeting} (après la réunion) ; \eng{after finishing school} (après avoir fini l'école).

\item \textbf{After, conjonction} : suivi d'une \textbf{proposition complète} (sujet + verbe conjugué).

\eng{After he sold his goods, he went home.} (Après avoir vendu ses marchandises, il est rentré chez lui.)
\end{enumerate}
\end{propriete}

\textbf{1. After + nom ou pronom.} C'est l'emploi le plus simple, équivalent du français « après + nom ».

\eng{After the harvest, farmers sell their crops.} « Après la récolte, les agriculteurs vendent leurs récoltes. »

\eng{I will call you after the lesson.} « Je t'appellerai après le cours. »

\eng{She arrived after me.} « Elle est arrivée après moi. » (attention : \eng{after me}, pas \eng{after I} — après une préposition, l'anglais exige le pronom \emph{complément} : \eng{me}, \eng{him}, \eng{her}, \eng{us}, \eng{them}.)

\textbf{2. After + gérondif (-ing).} Quand le sujet des deux actions est le même, l'anglais préfère souvent le gérondif, exactement comme le français « après + infinitif » ou « après avoir + participe ».

\eng{After finishing school, many young Cameroonians look for a job.} « Après avoir terminé l'école, beaucoup de jeunes Camerounais cherchent un emploi. »

\eng{After selling his cocoa, the farmer paid his workers.} « Après avoir vendu son cacao, le planteur a payé ses ouvriers. »

\textbf{3. After + sujet + verbe (conjonction).} Quand les deux sujets sont différents, la proposition complète est obligatoire.

\eng{After the manager arrived, the meeting started.} « Après l'arrivée du directeur, la réunion a commencé. » (litt. : après que le directeur fut arrivé)

\eng{After she left the bank, she went to the market.} « Après avoir quitté la banque, elle est allée au marché. »

\begin{attention}[After + futur : interdit !]
Comme toutes les conjonctions de temps (\eng{when}, \eng{before}, \eng{as soon as}, \eng{until}), \cle{after} n'est \textbf{jamais suivi du futur} (\eng{will}), même quand le sens est futur. On utilise le \textbf{présent simple} à la place.

\eng{I will call you after the meeting \textbf{ends}.} (et non \eng{after the meeting will end}) « Je t'appellerai après la fin de la réunion. »

\eng{After I \textbf{finish} my exams, I will work in my uncle's shop.} « Après avoir fini mes examens, je travaillerai dans la boutique de mon oncle. »
\end{attention}

\courssub{Afterwards : l'adverbe qui se suffit à lui-même}

\begin{propriete}[Afterwards = « ensuite », toujours employé SEUL]
\cle{Afterwards} est un \textbf{adverbe} qui signifie « ensuite », « après cela », « par la suite ». Contrairement à \eng{after}, il n'est \textbf{jamais suivi d'un complément} : ni nom, ni proposition. Il se place en \textbf{fin de phrase} ou en \textbf{début de phrase suivi d'une virgule}.

\eng{We had lunch; afterwards, we visited the factory.} « Nous avons déjeuné ; ensuite, nous avons visité l'usine. »

\eng{The students wrote the test and went home afterwards.} « Les élèves ont composé et sont rentrés chez eux ensuite. »
\end{propriete}

\begin{savaistu}[Afterward ou afterwards ?]
Les deux formes existent et ont exactement le même sens. \eng{Afterwards} (avec un \emph{s}) est la forme \textbf{britannique}, celle que nous utilisons dans ce cours et qui est attendue au Baccalauréat ; \eng{afterward} (sans \emph{s}) est la forme plutôt \textbf{américaine}. Même logique pour \eng{towards/toward}, \eng{forwards/forward}.
\end{savaistu}

\begin{exemplebox}[After et afterwards dans la même histoire]
\eng{After the harvest, farmers sell their crops.} « Après la récolte, les agriculteurs vendent leurs récoltes. »

\eng{We visited the market; afterwards, we met the manager.} « Nous avons visité le marché ; ensuite, nous avons rencontré le directeur. »

\eng{After he had finished his studies, he found a job in Douala.} « Après avoir terminé ses études, il a trouvé un emploi à Douala. »

\eng{First, we interviewed the workers. Afterwards, we wrote our report.} « D'abord, nous avons interrogé les ouvriers. Ensuite, nous avons rédigé notre rapport. »
\end{exemplebox}

Remarquez la logique : \eng{after} \emph{attache} l'événement antérieur à la phrase ; \eng{afterwards} \emph{conclut} la phrase en renvoyant à ce qui vient d'être dit. On peut souvent transformer une structure en l'autre :

\eng{After the meeting, we signed the contract.} = \eng{We had a meeting; afterwards, we signed the contract.}

\courssub{Ordre logique et temps des verbes : after + past perfect}

\begin{propriete}[La proposition avec after exprime l'action ANTÉRIEURE]
Dans une phrase avec \cle{after}, l'action de la proposition subordonnée se produit \textbf{en premier}, et l'action de la proposition principale \textbf{ensuite} — même si la subordonnée est placée après dans la phrase :

\eng{She opened a shop after she had saved enough money.} = d'abord elle a épargné, ensuite elle a ouvert une boutique.

Pour marquer clairement cette \textbf{antériorité}, l'anglais emploie souvent le \textbf{past perfect} (\eng{had + participe passé}) dans la proposition introduite par \eng{after}, avec le \textbf{prétérit simple} dans la principale :

\textbf{After + sujet + had + participe passé, sujet + prétérit simple.}

\eng{After she had saved enough money, she opened a shop.} « Après avoir économisé assez d'argent, elle a ouvert une boutique. »
\end{propriete}

Le past perfect n'est pas toujours obligatoire (\eng{after} indique déjà l'ordre des événements), mais il est \emph{recommandé} et très apprécié au Bac, car il montre la maîtrise de la concordance des temps.

\begin{center}
\begin{tabularx}{\linewidth}{|X|X|}
\hline
\rowcolor{popblueL} Prétérit simple (correct) & Past perfect (plus précis, valorisé au Bac) \\
\hline
\eng{After he sold his goods, he went home.} & \eng{After he had sold his goods, he went home.} \\
\hline
\eng{After the lesson ended, we played football.} & \eng{After the lesson had ended, we played football.} \\
\hline
\eng{After she finished her studies, she found a job.} & \eng{After she had finished her studies, she found a job.} \\
\hline
\end{tabularx}
\end{center}

\begin{popfigure}
\begin{center}
\begin{tikzpicture}[font=\small]
\node[draw=popblue, fill=popblueL, rounded corners, align=center, text width=4.6cm, minimum height=1.6cm] (e1) at (0,0) {\textbf{Événement 1 (antérieur)}\\ \eng{After she had saved money...}\\ (after + nom / -ing / sujet + verbe)};
\node[draw=poporange, fill=poporangeL, rounded corners, align=center, text width=4.6cm, minimum height=1.6cm] (e2) at (7.2,0) {\textbf{Événement 2 (postérieur)}\\ \eng{...she opened a shop.}\\ (proposition principale)};
\draw[-{Stealth[length=3mm]}, line width=1.2pt, popdark] (e1) -- (e2) node[midway, above, align=center, text width=2.2cm] {temps qui passe};
\node[draw=poppurple, fill=poppinkL, rounded corners, align=center, text width=5.2cm] (e3) at (3.6,-2.4) {\textbf{Variante avec afterwards}\\ \eng{She saved money; afterwards, she opened a shop.}};
\draw[-{Stealth[length=2.5mm]}, poppurple, dashed] (e1.south) -- (e3.north west);
\draw[-{Stealth[length=2.5mm]}, poppurple, dashed] (e2.south) -- (e3.north east);
\end{tikzpicture}
\end{center}
\legende{Organigramme de la succession temporelle : \eng{after} introduit l'événement antérieur ; \eng{afterwards} conclut la seconde phrase.}
\end{popfigure}

\courssub{Synonymes de afterwards : later, then, after that}

\begin{center}
\begin{tabularx}{\linewidth}{|l|X|X|}
\hline
\rowcolor{popblueL} Mot & Emploi et nuance & Exemple \\
\hline
\eng{afterwards} & « ensuite », soutenu, écrit et oral ; jamais de complément & \eng{We had lunch; afterwards, we visited the factory.} \\
\hline
\eng{then} & « puis, ensuite », enchaîne les actions d'un récit & \eng{We sold the crops, then we counted the money.} \\
\hline
\eng{later} & « plus tard », insiste sur l'écart de temps & \eng{I will see you later.} \\
\hline
\eng{after that} & « après cela », renvoie explicitement à ce qui précède & \eng{He lost his job; after that, he started a small business.} \\
\hline
\end{tabularx}
\end{center}

\begin{attention}[Les deux erreurs classiques des francophones]
\textbf{Erreur 1 — donner un complément à afterwards.} La phrase \eng{afterwards the lesson} est \textbf{FAUSSE} : \eng{afterwards} ne prend jamais de complément. Il faut dire \eng{after the lesson} (après le cours) ou restructurer : \eng{The lesson ended; afterwards, we played football.}

\textbf{Erreur 2 — employer after seul.} La phrase \eng{I will see you after} est \textbf{FAUSSE} en anglais standard : \eng{after} exige un complément. Dites \eng{I will see you later} ou \eng{I will see you afterwards} (« je te verrai plus tard / ensuite »). De même, ne dites pas \eng{After, we went home} mais \eng{Afterwards, we went home} ou \eng{After the match, we went home}.
\end{attention}

\begin{methode}[Comment choisir entre after et afterwards]
Posez-vous une seule question : \textbf{le mot est-il suivi de quelque chose ?}

\begin{enumerate}
\item Si un \textbf{nom}, un \textbf{pronom}, un \textbf{gérondif} ou une \textbf{proposition} (sujet + verbe) suit immédiatement $\rightarrow$ utilisez \cle{after}. Ex. : \eng{after work}, \eng{after leaving school}, \eng{after he arrived}.
\item Si le mot \textbf{termine la phrase} ou est \textbf{suivi d'une virgule} (c'est-à-dire suivi de rien) $\rightarrow$ utilisez \cle{afterwards} (ou \eng{later}, \eng{then}). Ex. : \eng{We visited the factory afterwards.}
\item Vérifiez enfin la \textbf{concordance des temps} : action antérieure au past perfect (\eng{after she had saved}), action postérieure au prétérit simple (\eng{she opened a shop}) ; et jamais de futur après \eng{after}.
\end{enumerate}
\end{methode}

\begin{exoresolu}[Application guidée]
\textbf{Énoncé.} Complétez avec \eng{after} ou \eng{afterwards}, puis traduisez.

1. \eng{......... the interview, the candidate received the job.}

2. \eng{We visited the port of Douala; ........., we wrote a report.}

3. \eng{......... she had repaid her loan, she borrowed more money.}

4. \eng{The farmers worked hard and rested ......... .}

5. \eng{......... graduating from university, my sister became an accountant.}

\tcblower
\textbf{Solution détaillée.}

1. \eng{\textbf{After} the interview, the candidate received the job.} — \eng{After} car un nom (\eng{the interview}) suit. « Après l'entretien, le candidat a obtenu le poste. »

2. \eng{We visited the port of Douala; \textbf{afterwards}, we wrote a report.} — \eng{Afterwards} car le mot est suivi d'une virgule, sans complément. « Nous avons visité le port de Douala ; ensuite, nous avons rédigé un rapport. »

3. \eng{\textbf{After} she had repaid her loan, she borrowed more money.} — \eng{After} car une proposition complète (\eng{she had repaid...}) suit ; notez le past perfect pour l'antériorité. « Après avoir remboursé son prêt, elle a emprunté davantage d'argent. »

4. \eng{The farmers worked hard and rested \textbf{afterwards}.} — \eng{Afterwards} en fin de phrase, sans complément. « Les agriculteurs ont travaillé dur et se sont reposés ensuite. »

5. \eng{\textbf{After} graduating from university, my sister became an accountant.} — \eng{After} + gérondif (\eng{graduating}). « Après avoir obtenu son diplôme universitaire, ma sœur est devenue comptable. »
\end{exoresolu}

\begin{exercice}[À vous de jouer]
1. Transformez en utilisant \eng{afterwards} : \eng{After the training course, the workers received their certificates.}

2. Transformez en utilisant \eng{after} + past perfect : \eng{He saved money for three years; afterwards, he bought a bendskin.}

3. Corrigez les erreurs : \eng{Afterwards the exams, I will work in a bank. I will see you after.}

4. Traduisez : « Après avoir vendu son cacao, le planteur a payé ses dettes ; ensuite, il a ouvert un compte en banque. »
\end{exercice}

\begin{corrige}[Exercice « À vous de jouer »]
1. \eng{The workers attended a training course; afterwards, they received their certificates.} (on supprime le complément et l'on place \eng{afterwards} en début de seconde phrase, suivi d'une virgule).

2. \eng{After he had saved money for three years, he bought a bendskin.} (action antérieure au past perfect : \eng{had saved} ; action postérieure au prétérit : \eng{bought}).

3. \eng{\textbf{After} the exams, I will work in a bank. I will see you \textbf{later / afterwards}.} (\eng{afterwards} ne peut pas précéder un nom ; \eng{after} ne peut pas terminer une phrase).

4. \eng{After he had sold his cocoa (ou After selling his cocoa), the farmer paid his debts; afterwards, he opened a bank account.}
\end{corrige}

\begin{aretenir}[L'essentiel de la section]
\begin{itemize}
\item \cle{after} n'est jamais seul : \eng{after} + nom (\eng{after work}), + gérondif (\eng{after finishing}), ou + sujet + verbe (\eng{after he arrived}).
\item \cle{afterwards} = « ensuite » : toujours seul, en fin de phrase ou suivi d'une virgule.
\item La proposition avec \eng{after} exprime l'action \textbf{antérieure} ; au Bac, valorisez-la avec le \textbf{past perfect} : \eng{After she had saved enough money, she opened a shop.}
\item Jamais de futur après \eng{after} : \eng{after the meeting ends}.
\item Pièges à éviter : \eng{afterwards the lesson} et \eng{I will see you after} sont tous deux incorrects.
\end{itemize}
\end{aretenir}

\coursec{Comparaison français/anglais et pièges des francophones}

Nous venons d'étudier \eng{after} et \eng{afterwards} pour exprimer la succession des événements. Il est temps maintenant de confronter l'ensemble de ces adverbes de temps à notre langue maternelle. C'est précisément ici que les élèves francophones commettent le plus d'erreurs : le français et l'anglais ne découpent pas le temps de la même manière, et la traduction mot à mot est un piège redoutable. Observons d'abord quelques paires de phrases avant d'énoncer les règles.

\begin{exemplebox}[Observation : le même mot français, deux mots anglais]
\eng{I have lived in Douala \textbf{for} ten years.} « J'habite à Douala \textbf{depuis} dix ans. »\par
\eng{I have lived in Douala \textbf{since} 2015.} « J'habite à Douala \textbf{depuis} 2015. »\par
\eng{My uncle left Bamenda two days \textbf{ago}.} « Mon oncle a quitté Bamenda \textbf{il y a} deux jours. »\par
\eng{She finished her work; \textbf{afterwards}, she went to the market.} « Elle a fini son travail ; \textbf{ensuite}, elle est allée au marché. »
\end{exemplebox}

Que remarque-t-on ? Le mot français « depuis » se traduit tantôt par \eng{for}, tantôt par \eng{since} ; « il y a » se place en \emph{début} de groupe en français, alors que \eng{ago} se place en \emph{fin} de phrase en anglais ; et le verbe français est au présent là où l'anglais exige le \emph{present perfect}. Décortiquons chaque piège.

\courssub{« Depuis » = for ou since : durée ou point de départ ?}

\begin{propriete}[La règle d'or]
Le français utilise un seul mot, « depuis ». L'anglais, lui, fait un choix selon la nature de ce qui suit :
\begin{itemize}
\item \cle{for} + une \textbf{durée} (une quantité de temps) : \eng{for two hours}, \eng{for six months}, \eng{for ten years}, \eng{for a long time}.
\item \cle{since} + un \textbf{point de départ} (une date, un moment précis) : \eng{since Monday}, \eng{since 2019}, \eng{since last June}, \eng{since his childhood}.
\end{itemize}
Le test est simple : si l'on peut répondre à la question « pendant combien de temps ? », c'est une durée, donc \eng{for} ; si l'on répond à « à partir de quel moment ? », c'est un point de départ, donc \eng{since}.
\end{propriete}

\begin{propriete}[Since + proposition : attention au temps de la subordonnée]
\eng{Since} peut aussi introduire une \textbf{proposition entière} (sujet + verbe). Dans ce cas, le verbe de la subordonnée se met au \textbf{prétérit simple}, car il désigne le point de départ révolu, tandis que la principale reste au present perfect :\par
\eng{I have known her since she was a child.} « Je la connais depuis son enfance. »\par
\eng{We have been friends since we met at secondary school.} « Nous sommes amis depuis que nous nous sommes rencontrés au collège. »\par
Forme : \textbf{present perfect + since + sujet + prétérit simple}.
\end{propriete}

\begin{center}
\begin{tabularx}{\linewidth}{|X|X|X|}
\hline
\rowcolor{popblueL} Français & English & Remarque \\
\hline
depuis trois jours & for three days & durée : combien de temps ? \\
\hline
depuis lundi & since Monday & point de départ : à partir de quand ? \\
\hline
depuis longtemps & for a long time & durée \\
\hline
depuis son enfance & since his childhood & point de départ \\
\hline
depuis deux ans & for two years & durée (jamais \emph{since two years}) \\
\hline
depuis 2020 & since 2020 & point de départ (jamais \emph{for 2020}) \\
\hline
depuis qu'il est parti & since he left & proposition : prétérit après \eng{since} \\
\hline
\end{tabularx}
\end{center}

\courssub{Le temps du verbe : le présent français contre le present perfect anglais}

C'est le contraste le plus fondamental entre les deux langues, et le plus sanctionné au Bac.

\begin{propriete}[Présent français, present perfect anglais]
Pour une action commencée dans le passé et qui continue encore au moment où l'on parle :
\begin{itemize}
\item le français utilise le \textbf{présent} : « J'apprends l'anglais depuis cinq ans. »
\item l'anglais exige le \textbf{present perfect} (ou le present perfect continuous) : \eng{I have been learning English for five years.}
\end{itemize}
Employer le présent simple anglais avec \eng{for} ou \eng{since} est une faute grave : \emph{I learn English since five years} est doublement incorrect (\eng{since} mal choisi et présent simple interdit).
\end{propriete}

\begin{exemplebox}[Exemples traduits]
\eng{Il pleut depuis deux jours.} = \eng{It has been raining for two days.}\par
\eng{Nous connaissons le directeur depuis 2018.} = \eng{We have known the director since 2018.}\par
\eng{Elle travaille dans cette banque depuis dix ans.} = \eng{She has worked in this bank for ten years.}\par
\eng{Ils attendent le car depuis une heure.} = \eng{They have been waiting for the bus for an hour.}
\end{exemplebox}

\begin{attention}[Attention : « il y a » se renverse en anglais]
En français, « il y a » précède la durée : « Il est parti \textbf{il y a une heure}. » En anglais, l'ordre est inversé : la durée vient d'abord, \eng{ago} vient en dernier : \eng{He left \textbf{an hour ago}.} De plus, \eng{ago} situe un événement \emph{terminé} dans un passé \emph{révolu} : le verbe est donc au \textbf{prétérit simple}, jamais au present perfect.\par
\eng{Il est parti il y a une heure.} = \eng{He left an hour ago.}\par
\eng{Elle a vendu sa boutique il y a trois mois.} = \eng{She sold her shop three months ago.}\par
Forme : \textbf{sujet + verbe au prétérit + durée + ago}.
\end{attention}

\courssub{« Après » = after ; « ensuite » = afterwards}

Le français utilise « après » à la fois comme préposition et comme adverbe. L'anglais distingue deux mots :

\begin{propriete}[After ou afterwards ?]
\begin{itemize}
\item \cle{after} est une \textbf{préposition} (ou une conjonction) : il est toujours \emph{suivi} d'un nom, d'un pronom, d'un verbe en \emph{-ing} ou d'une proposition : \eng{after the match}, \eng{after lunch}, \eng{after finishing his work}, \eng{after he arrived}.
\item \cle{afterwards} est un \textbf{adverbe} : il se suffit à lui-même et n'est jamais suivi d'un complément : \eng{We had dinner; afterwards, we watched television.}
\end{itemize}
On ne peut donc jamais dire \emph{afterwards the match} : il faut \eng{after the match}. À l'inverse, on ne peut pas terminer une phrase par \eng{after} seul : \emph{We watched television after} est incorrect ; il faut \eng{afterwards} (ou \eng{after that}).
\end{propriete}

\begin{exemplebox}[Exemples traduits]
\eng{D'abord il a travaillé, ensuite il s'est reposé.} = \eng{First he worked; afterwards he rested.}\par
\eng{Après le match, les joueurs sont rentrés chez eux.} = \eng{After the match, the players went home.}\par
\eng{Après avoir fini ses devoirs, elle est sortie.} = \eng{After finishing her homework, she went out.}\par
\eng{Nous avons visité Buea, et ensuite nous sommes montés sur la montagne.} = \eng{We visited Buea, and afterwards we climbed the mountain.}
\end{exemplebox}

\begin{attention}[After + -ing, jamais after + to]
Après \eng{after}, le verbe se met à la forme en \emph{-ing} : \eng{after finishing}, \eng{after eating}, \eng{after leaving}. La construction \emph{after to finish} est un calque du français « après finir / après avoir fini » et elle est toujours fautive. Pour exprimer l'antériorité, on peut aussi utiliser \eng{after having finished}, mais \eng{after finishing} suffit dans la plupart des cas.
\end{attention}

\courssub{Carte mentale récapitulative}

\begin{popfigure}
\begin{center}\tcbox[colback=popdark,colframe=popdark,coltext=white,arc=9pt,boxsep=4pt,left=6pt,right=6pt]{\bfseries\small Adverbes de temps}\end{center}
\vspace{-2pt}
\begin{tcbraster}[raster columns=2,raster equal height=rows,raster column skip=6pt,raster row skip=6pt,enhanced,blankest]
\begin{tcolorbox}[enhanced,colback=white,colframe=popgreen,boxrule=0.9pt,arc=3pt,title={FOR durée},fonttitle=\bfseries\scriptsize,coltitle=white,colbacktitle=popgreen,left=4pt,right=4pt,top=2pt,bottom=2pt,boxsep=1pt,toptitle=1.5pt,bottomtitle=1.5pt,halign=left]
\begin{itemize}[nosep,leftmargin=*,itemsep=1pt,font=\scriptsize]
\item for two years for a long time
\item present perfect
\end{itemize}
\end{tcolorbox}
\begin{tcolorbox}[enhanced,colback=white,colframe=poporange,boxrule=0.9pt,arc=3pt,title={SINCE point de départ},fonttitle=\bfseries\scriptsize,coltitle=white,colbacktitle=poporange,left=4pt,right=4pt,top=2pt,bottom=2pt,boxsep=1pt,toptitle=1.5pt,bottomtitle=1.5pt,halign=left]
\begin{itemize}[nosep,leftmargin=*,itemsep=1pt,font=\scriptsize]
\item since 2020 since Monday
\item present perfect
\end{itemize}
\end{tcolorbox}
\begin{tcolorbox}[enhanced,colback=white,colframe=poppurple,boxrule=0.9pt,arc=3pt,title={AGO passé révolu},fonttitle=\bfseries\scriptsize,coltitle=white,colbacktitle=poppurple,left=4pt,right=4pt,top=2pt,bottom=2pt,boxsep=1pt,toptitle=1.5pt,bottomtitle=1.5pt,halign=left]
\begin{itemize}[nosep,leftmargin=*,itemsep=1pt,font=\scriptsize]
\item an hour ago (durée + ago)
\item prétérit simple
\end{itemize}
\end{tcolorbox}
\begin{tcolorbox}[enhanced,colback=white,colframe=poppink,boxrule=0.9pt,arc=3pt,title={AFTER / AFTERWARDS succession},fonttitle=\bfseries\scriptsize,coltitle=white,colbacktitle=poppink,left=4pt,right=4pt,top=2pt,bottom=2pt,boxsep=1pt,toptitle=1.5pt,bottomtitle=1.5pt,halign=left]
\begin{itemize}[nosep,leftmargin=*,itemsep=1pt,font=\scriptsize]
\item after + nom after + -ing
\item afterwards seul, sans complément
\end{itemize}
\end{tcolorbox}
\end{tcbraster}

\legende{Carte mentale : chaque adverbe de temps appelle une structure et un temps verbal précis.}
\end{popfigure}

\courssub{Tableau récapitulatif des faux pas des francophones}

\begin{center}
\begin{tabularx}{\linewidth}{|X|X|X|}
\hline
\rowcolor{popblueL} Erreur fréquente & Correction & Explication \\
\hline
\emph{since three years} & for three years & durée, pas point de départ \\
\hline
\emph{I work here since 2020} & I have worked here since 2020 & present perfect obligatoire \\
\hline
\emph{I am here since Monday} & I have been here since Monday & présent français, perfect anglais \\
\hline
\emph{ago two days} & two days ago & la durée précède \eng{ago} \\
\hline
\emph{two weeks ago I have met him} & two weeks ago I met him & \eng{ago} exige le prétérit \\
\hline
\emph{afterwards the match} & after the match & \eng{afterwards} ne prend pas de complément \\
\hline
\emph{after to finish} & after finishing & \eng{after} + verbe en \emph{-ing} \\
\hline
\emph{for 2020} & since 2020 & une date est un point de départ \\
\hline
\emph{since he has left} & since he left & prétérit dans la subordonnée en \eng{since} \\
\hline
\end{tabularx}
\end{center}

\begin{exoresolu}[Corrige et explique les erreurs]
Énoncé. Les phrases suivantes, produites par des candidats francophones, contiennent chacune une erreur. Corrige-les et justifie.\par
1. \eng{My father has his shop since fifteen years.}\par
2. \eng{I have seen a bendskin accident ago three days.}\par
3. \eng{After to sell her tomatoes, Mama Ngo Bell went home; afterwards the market was closed.}
\tcblower
Solution détaillée.\par
1. \eng{My father \textbf{has had} his shop \textbf{for} fifteen years.} Deux erreurs : « depuis quinze ans » est une \emph{durée}, donc \eng{for} et non \eng{since} ; et une action qui continue exige le present perfect, donc \eng{has had} et non \eng{has}.\par
2. \eng{I \textbf{saw} a bendskin accident \textbf{three days ago}.} Deux erreurs : la durée précède \eng{ago} (ordre inversé par rapport à « il y a trois jours »), et \eng{ago} impose le prétérit simple \eng{saw}, pas le present perfect.\par
3. \eng{After \textbf{selling} her tomatoes, Mama Ngo Bell went home; afterwards\textbf{,} the market was closed.} Deux corrections : \eng{after} est suivi du verbe en \emph{-ing} (\eng{after selling}, jamais \emph{after to sell}) ; et \eng{afterwards} est un adverbe autonome, suivi d'une virgule, jamais d'un nom — si l'on veut un complément, il faut \eng{after} : \eng{after the market closed}.
\end{exoresolu}

\begin{methode}[Vérifier une phrase au Bac en trois questions]
Face à une phrase contenant un adverbe de temps (épreuve de transformation ou de fill-in-the-blanks), procède ainsi :
\begin{enumerate}
\item \textbf{Repère l'adverbe} (ou le mot français à traduire : « depuis », « il y a », « après », « ensuite »).
\item \textbf{Durée ou point de départ ?} Pose-toi la question « pendant combien de temps ? » (durée : \eng{for}) ou « à partir de quand ? » (point de départ : \eng{since}). Pour « il y a », inverse l'ordre : durée + \eng{ago}. Pour « ensuite » seul, choisis \eng{afterwards} ; pour « après + complément », choisis \eng{after}.
\item \textbf{Vérifie le temps du verbe} : present perfect avec \eng{for} et \eng{since} (action qui continue), prétérit simple avec \eng{ago} (action révolue) et dans la subordonnée introduite par \eng{since}. C'est presque toujours là que se cache le point de grammaire testé.
\end{enumerate}
\end{methode}

\begin{savaistu}[For, since et ago au Bac A]
À l'épreuve d'anglais du Baccalauréat A, les exercices de grammaire — transformation de phrases, fill-in-the-blanks, correction d'erreurs — testent très régulièrement la paire \eng{for}/\eng{since} et le temps verbal qui l'accompagne. Un item classique consiste à transformer \eng{I started teaching in 2015} en \eng{I have been teaching since 2015}, ou à compléter \eng{She has lived in Yaoundé \ldots\ ten years} par \eng{for}. Maîtriser la distinction durée/point de départ et le couple adverbe-temps verbal, c'est sécuriser des points faciles dans la section grammar.
\end{savaistu}

\begin{exercice}[À toi de jouer]
Traduis en anglais : 1. « Nous attendons le car depuis vingt minutes. » 2. « Il a ouvert son atelier il y a six mois. » 3. « Elle étudie l'économie depuis septembre. » 4. « Après avoir payé, le client est parti ; ensuite, le commerçant a fermé la boutique. » 5. « Je connais ce commerçant depuis que je suis enfant. »
\end{exercice}

\begin{corrige}[Exercice « À toi de jouer »]
1. \eng{We have been waiting for the bus for twenty minutes.} (durée : \eng{for} + present perfect continuous).\par
2. \eng{He opened his workshop six months ago.} (ordre inversé : durée + \eng{ago} ; prétérit simple).\par
3. \eng{She has been studying economics since September.} (point de départ : \eng{since} + present perfect continuous).\par
4. \eng{After paying, the customer left; afterwards, the shopkeeper closed the shop.} (\eng{after} + \emph{-ing} ; \eng{afterwards} seul, suivi d'une virgule).\par
5. \eng{I have known this trader since I was a child.} (proposition après \eng{since} : prétérit simple \eng{was} dans la subordonnée, present perfect \eng{have known} dans la principale).
\end{corrige}

\coursec{Entraînement guidé : application et transformation}

Nous savons maintenant distinguer \eng{for} (durée), \eng{since} (point de départ), \eng{ago} (recul dans le passé) et \eng{after} / \eng{afterwards} (succession d'événements), et nous avons repéré les pièges que le français nous tend. Il est temps de passer à la pratique. Cette section vous propose un entraînement progressif, du choix de l'adverbe jusqu'à la production personnelle, exactement comme à l'épreuve du Baccalauréat. Pour chaque exercice, justifiez votre réponse par la règle : c'est la justification qui fait la différence entre un élève qui devine et un élève qui sait.

\begin{methode}[La méthode en trois questions avant de répondre]
\begin{enumerate}
\item \textbf{Qu'est-ce qui suit le blanc ?} Une durée (\eng{two years, a long time}) ou un point de départ (\eng{2015, Monday, my childhood}) ?
\item \textbf{Quel est le temps du verbe ?} Present perfect (action qui continue jusqu'à aujourd'hui) ou prétérit (action terminée, coupée du présent) ?
\item \textbf{Quelle règle s'applique ?} Durée → \eng{for} ; point de départ → \eng{since} ; recul depuis aujourd'hui → \eng{ago} ; succession d'événements → \eng{after} / \eng{afterwards}.
\end{enumerate}
\end{methode}

\begin{exemplebox}[Trois exemples traités ensemble en classe]
\begin{itemize}
\item \eng{She has been a nurse \textbf{since} 2015.} — « Elle est infirmière depuis 2015. » Ce qui suit le blanc est une \emph{année}, donc un point de départ : \cle{since}. Le present perfect est correct car elle est encore infirmière aujourd'hui.
\item \eng{They worked in the bank \textbf{for} twenty years, then retired.} — « Ils ont travaillé à la banque pendant vingt ans, puis ont pris leur retraite. » Ce qui suit est une \emph{durée} : \cle{for}. Le prétérit est correct car la période est terminée.
\item \eng{He left school two years \textbf{ago}.} — « Il a quitté l'école il y a deux ans. » On recule depuis aujourd'hui : \cle{ago}, placé \emph{après} la durée, avec un prétérit.
\end{itemize}
\end{exemplebox}

\courssub{Exercice 1 : compléter avec for ou since}

\begin{exercice}[For or since? — les métiers]
Complétez chaque phrase avec \eng{for} ou \eng{since}. Toutes parlent de la vie professionnelle.
\begin{enumerate}
\item My uncle has been a taxi driver in Douala \_\_\_\_\_ fifteen years.
\item Mrs Ngo has taught mathematics \_\_\_\_\_ 2010.
\item The carpenter has worked in this workshop \_\_\_\_\_ a long time.
\item I have wanted to become a pilot \_\_\_\_\_ my childhood.
\item The shopkeeper has owned this store \_\_\_\_\_ last March.
\item We have known our family doctor \_\_\_\_\_ ages.
\item He has been unemployed \_\_\_\_\_ the factory closed.
\item The farmers have grown cocoa in this village \_\_\_\_\_ generations.
\item She has worked as a journalist \_\_\_\_\_ she left university.
\item They have lived in Bamenda \_\_\_\_\_ 2019.
\end{enumerate}
\end{exercice}

\begin{corrige}[Exercice 1]
\begin{enumerate}
\item \eng{for} — durée (\eng{fifteen years}).
\item \eng{since} — point de départ (une année).
\item \eng{for} — durée (\eng{a long time}).
\item \eng{since} — point de départ (\eng{my childhood}).
\item \eng{since} — point de départ (\eng{last March}).
\item \eng{for} — durée (\eng{ages}).
\item \eng{since} — point de départ (une proposition au prétérit : \eng{the factory closed}).
\item \eng{for} — durée (\eng{generations}).
\item \eng{since} — point de départ (proposition au prétérit : \eng{she left university}).
\item \eng{since} — point de départ (une année).
\end{enumerate}
\textbf{Justification de la règle :} \eng{for} + durée ; \eng{since} + point de départ (date, moment, ou proposition au prétérit). Le verbe principal reste au present perfect dans les deux cas.
\end{corrige}

\courssub{Exercice 2 : choisir le bon temps selon l'adverbe}

\begin{exoresolu}[Present perfect ou prétérit ?]
\textbf{Énoncé.} Mettez le verbe entre parenthèses au temps correct (present perfect ou prétérit simple) en observant l'adverbe.
\begin{enumerate}
\item She (work) \_\_\_\_\_ as a nurse since 2015.
\item They (work) \_\_\_\_\_ in the bank for twenty years, then retired.
\item He (leave) \_\_\_\_\_ school two years ago.
\item I (not see) \_\_\_\_\_ my cousin since last Christmas.
\end{enumerate}
\tcblower
\textbf{Solution détaillée.}
\begin{enumerate}
\item \eng{She \textbf{has worked} as a nurse since 2015.} — \eng{since} + point de départ, action qui continue aujourd'hui → present perfect.
\item \eng{They \textbf{worked} in the bank for twenty years, then retired.} — la durée est \emph{terminée} (ils sont à la retraite) → prétérit, même avec \eng{for}.
\item \eng{He \textbf{left} school two years ago.} — \eng{ago} impose toujours le prétérit : l'événement est coupé du présent.
\item \eng{I \textbf{haven't seen} my cousin since last Christmas.} — \eng{since} + point de départ, situation encore vraie → present perfect.
\end{enumerate}
\textbf{Règle :} l'adverbe commande le temps. \eng{ago} → prétérit obligatoire ; \eng{since} → present perfect ; \eng{for} → present perfect si la durée continue, prétérit si elle est achevée.
\end{exoresolu}

\begin{exercice}[À vous de choisir le temps]
Mettez le verbe au temps correct.
\begin{enumerate}
\item We (live) \_\_\_\_\_ in Yaoundé for five years. We love it here.
\item My father (buy) \_\_\_\_\_ his first car ten years ago.
\item The company (employ) \_\_\_\_\_ fifty workers since it opened.
\item She (finish) \_\_\_\_\_ her training two months ago and (start) \_\_\_\_\_ her new job immediately afterwards.
\item They (be) \_\_\_\_\_ business partners for a decade before they separated.
\end{enumerate}
\end{exercice}

\begin{corrige}[Exercice 2]
\begin{enumerate}
\item \eng{have lived} — durée qui continue (\eng{We love it here}) → present perfect.
\item \eng{bought} — \eng{ago} impose le prétérit (verbe irrégulier : \eng{buy -- bought -- bought}).
\item \eng{has employed} — \eng{since} + proposition, résultat présent → present perfect.
\item \eng{finished} ; \eng{started} — \eng{ago} et \eng{afterwards} racontent des faits passés : prétérit.
\item \eng{were} — durée terminée (\eng{before they separated}) : prétérit malgré \eng{for}.
\end{enumerate}
\end{corrige}

\courssub{Exercice 3 : transformer for ↔ since}

\begin{methode}[Transformer for en since (et inversement) sans changer le sens]
\begin{enumerate}
\item \textbf{Repérez la durée} après \eng{for} : ex. \eng{for ten years}.
\item \textbf{Calculez le point de départ} : année actuelle $-$ durée. Ex. : $2024 - 10 = 2014$.
\item \textbf{Remplacez} \eng{for ten years} par \eng{since 2014}.
\item \textbf{Ne touchez à rien d'autre} : le temps du verbe reste identique.
\item Sens inverse : \eng{since 2018} → calculez la durée écoulée ($2024 - 2018 = 6$) → \eng{for six years}.
\end{enumerate}
\end{methode}

\begin{attention}[Ne changez jamais le temps verbal dans une transformation]
Dans les exercices de transformation, seuls l'adverbe et son complément changent. Le temps du verbe, le sujet et le reste de la phrase restent strictement identiques.\\
\eng{She has worked here for ten years.} → \eng{She has worked here since 2014.} (et non \emph{she worked}, ni \emph{she is working}).
\end{attention}

\begin{popfigure}
\begin{center}
\begin{tikzpicture}[node distance=1.1cm and 2.2cm]
\node[draw=popblue, fill=popblueL, rounded corners, align=center, text width=4.6cm] (start) {\eng{for ten years}\\ durée};
\node[draw=poporange, fill=poporangeL, rounded corners, align=center, text width=4.6cm, right=of start] (calc) {calcul :\\ année actuelle $-$ durée\\ $2024 - 10 = 2014$};
\node[draw=popgreen, fill=popgreenL, rounded corners, align=center, text width=4.6cm, right=of calc] (end) {\eng{since 2014}\\ point de départ};
\draw[-{Stealth}, very thick, poppurple] (start) -- (calc);
\draw[-{Stealth}, very thick, poppurple] (calc) -- (end);
\node[below=0.5cm of calc, align=center, text=popdark] {\small le temps du verbe ne change pas : \eng{has worked} $\to$ \eng{has worked}};
\end{tikzpicture}
\end{center}
\legende{Schéma de la transformation \eng{for} $\leftrightarrow$ \eng{since}}
\end{popfigure}

\begin{exercice}[Transformez sans changer le sens]
Réécrivez chaque phrase en transformant \eng{for} en \eng{since} ou l'inverse (année actuelle : 2024).
\begin{enumerate}
\item I have been learning English for six years.
\item She has been a nurse since 2015.
\item We have lived in Buea for three years.
\item He has driven a bendskin since 2020.
\item They have known each other for a decade.
\end{enumerate}
\end{exercice}

\begin{corrige}[Exercice 3]
\begin{enumerate}
\item \eng{I have been learning English since 2018.} ($2024 - 6 = 2018$)
\item \eng{She has been a nurse for nine years.} ($2024 - 2015 = 9$)
\item \eng{We have lived in Buea since 2021.} ($2024 - 3 = 2021$)
\item \eng{He has driven a bendskin for four years.} ($2024 - 2020 = 4$)
\item \eng{They have known each other since 2014.} (une décennie = dix ans)
\end{enumerate}
Dans chaque cas, le present perfect est conservé : seul le complément de temps change.
\end{corrige}

\begin{savaistu}[La transformation au Baccalauréat]
Au Bac, une phrase comme \eng{I have been learning English for six years} peut vous être donnée à transformer en \eng{I have been learning English since 2018}. La consigne exige que \emph{le sens reste identique} : si vous changez le temps du verbe ou le sujet, la réponse est fausse même si la grammaire est correcte. Entraînez-vous à vérifier mentalement le calcul de l'année de départ.
\end{savaistu}

\courssub{Exercice 4 : corriger les erreurs de francophones}

\begin{exercice}[Trouvez et corrigez l'erreur]
Chaque phrase contient une erreur typique d'un élève francophone. Corrigez-la et citez la règle.
\begin{enumerate}
\item I am living in Douala since five years.
\item He has left school two years ago.
\item She works here since 2020.
\item They have worked in the bank since twenty years.
\item We have seen him ago two days.
\item After the meeting, we afterwards went home.
\end{enumerate}
\end{exercice}

\begin{corrige}[Exercice 4]
\begin{enumerate}
\item \eng{I \textbf{have lived} in Douala \textbf{for} five years.} — durée → \eng{for} ; action qui continue → present perfect (le français « je vis depuis » trompe vers le présent).
\item \eng{He \textbf{left} school two years ago.} — \eng{ago} est incompatible avec le present perfect : prétérit obligatoire.
\item \eng{She \textbf{has worked} here since 2020.} — \eng{since} exige le present perfect, pas le présent simple.
\item \eng{They have worked in the bank \textbf{for} twenty years.} — \eng{twenty years} est une durée, pas un point de départ.
\item \eng{We \textbf{saw} him \textbf{two days ago}.} — \eng{ago} se place \emph{après} la durée et impose le prétérit (verbe irrégulier : \eng{see -- saw -- seen}).
\item \eng{After the meeting, we went home.} ou \eng{We went home afterwards.} — on ne cumule pas \eng{after} et \eng{afterwards} dans la même proposition.
\end{enumerate}
\end{corrige}

\courssub{Exercice 5 : réordonner les mots (ago, after, afterwards)}

\begin{exercice}[Remettez les mots dans l'ordre]
\begin{enumerate}
\item years / three / graduated / ago / she
\item the / after / lunch / went / we / back / office / to
\item afterwards / shopping / went / they
\item ago / long / how / did / retire / he / ?
\item interview / after / the / called / manager / me / the
\end{enumerate}
\end{exercice}

\begin{corrige}[Exercice 5]
\begin{enumerate}
\item \eng{She graduated three years ago.} — \eng{ago} en fin de groupe de durée.
\item \eng{After lunch, we went back to the office.} — \eng{after} + nom, en tête de phrase suivi d'une virgule.
\item \eng{They went shopping afterwards.} — \eng{afterwards} seul, en fin de phrase.
\item \eng{How long ago did he retire?} — dans une question, \eng{ago} reste après \eng{how long}.
\item \eng{After the interview, the manager called me.} — \eng{after} + groupe nominal.
\end{enumerate}
\end{corrige}

\courssub{Exercice 6 : production écrite personnelle}

\begin{experience}[Cinq phrases sur votre vie]
Rédigez cinq phrases vraies sur vous-même ou un proche, en utilisant chacun des outils étudiés :
\begin{enumerate}
\item une phrase avec \eng{for} (durée d'études ou d'habitation) : \eng{I have studied at this school for ...}
\item une phrase avec \eng{since} (point de départ) : \eng{I have known my best friend since ...}
\item une phrase avec \eng{ago} (événement passé) : \eng{My family moved to Douala three years ago.}
\item une phrase avec \eng{after} : \eng{After the BEPC, I ...}
\item une phrase avec \eng{afterwards} : \eng{We celebrated; afterwards, we ...}
\end{enumerate}
Échangez ensuite votre production avec un voisin : chacun vérifie l'adverbe, sa position et le temps du verbe, puis justifie ses corrections.
\end{experience}

\begin{exemplebox}[Modèle de production attendue]
\eng{I have lived in Yaoundé for twelve years. I have wanted to become an engineer since my childhood. My elder brother started his own business two years ago. After my A Levels, I will study at the university. I will work hard first and travel afterwards.}
\end{exemplebox}

\courssub{Correction collective et tableau récapitulatif final}

La correction collective se fait au tableau : un élève propose une réponse, la classe la valide ou la corrige, et l'élève \emph{justifie} par la règle (durée ou point de départ ? action continue ou terminée ? position de l'adverbe ?). Ce rituel de justification est la clé de la réussite à l'examen.

\begin{popfigure}
\begin{center}
\begin{tabularx}{\linewidth}{|l|X|X|l|X|}
\hline
\rowcolor{popblueL} \textbf{Adverbe} & \textbf{Sens} & \textbf{Ce qui suit} & \textbf{Temps du verbe} & \textbf{Exemple traduit} \\
\hline
\eng{for} & pendant / depuis (durée) & une durée : \eng{ten years, ages, a long time} & present perfect (durée continue) ou prétérit (durée terminée) & \eng{She has worked here for ten years.} « Elle travaille ici depuis dix ans. » \\
\hline
\eng{since} & depuis (point de départ) & une date, un moment, une proposition au prétérit : \eng{2015, Monday, I was born} & present perfect & \eng{He has been a driver since 2020.} « Il est chauffeur depuis 2020. » \\
\hline
\eng{ago} & il y a & une durée, placée \emph{avant} : \eng{two years ago} & prétérit simple & \eng{They left two hours ago.} « Ils sont partis il y a deux heures. » \\
\hline
\eng{after} & après & un nom ou une proposition : \eng{after lunch, after she left} & selon le contexte & \eng{After work, he went home.} « Après le travail, il est rentré. » \\
\hline
\eng{afterwards} & ensuite, après coup & rien (adverbe seul, souvent en fin de phrase) & selon le contexte & \eng{We had dinner and went out afterwards.} « Nous avons dîné puis sommes sortis ensuite. » \\
\hline
\end{tabularx}
\end{center}
\legende{Tableau récapitulatif : \eng{for, since, ago, after, afterwards}}
\end{popfigure}

\begin{aretenir}[Les cinq réflexes à garder]
\begin{itemize}
\item Durée → \cle{for} ; point de départ → \cle{since} ; recul dans le passé → \cle{ago}.
\item \eng{ago} se place après la durée et impose le prétérit.
\item \eng{since} s'associe au present perfect ; \eng{for} accepte les deux temps selon que la durée continue ou non.
\item \eng{after} + complément ; \eng{afterwards} seul : jamais ensemble.
\item Dans une transformation \eng{for} ↔ \eng{since}, seul le complément change : le temps du verbe est intouchable.
\end{itemize}
\end{aretenir}

\newpage

\coursec{Activité d'intégration}

\coursec{Lesson 14 — Grammar: Adverbs (for, since, ago, after, afterwards)}

\courssub{Objectifs de la leçon}

À l'issue de cette leçon, tu seras capable de :

\begin{itemize}
\item identifier et distinguer les adverbes et expressions de temps \cle{for}, \cle{since}, \cle{ago}, \cle{after} et \cle{afterwards} dans un texte ou un enregistrement authentique ;
\item associer chaque marqueur au temps verbal qui lui correspond (\emph{present perfect} pour \eng{for} et \eng{since} ; \emph{past simple} pour \eng{ago} et \eng{after} + proposition) ;
\item construire des phrases correctes pour exprimer la durée (\eng{for}), le point de départ (\eng{since}), l'antériorité datée (\eng{ago}) et la succession (\eng{after}, \eng{afterwards}) ;
\item éviter les erreurs classiques des francophones (\eng{since 8 years}, \eng{I work here since 2016}, \eng{ago two years}, confusion \eng{after}/\eng{afterwards}) ;
\item mobiliser ces outils dans une interaction orale (interview) et une production écrite (paragraphe biographique) en lien avec la vie économique et professionnelle.
\end{itemize}

\courssub{Observation guidée : les adverbes de temps en contexte}

Écoute (ou lis) le témoignage suivant d'un travailleur camerounais. Repère les mots en gras et note le temps verbal qui les accompagne.

\begin{textebox}[Listening script / Reading text: “A Career Path in Bamenda”]
\textbf{Mr Nfor:} “Good morning. My name is Peter Nfor. I have been a secondary school teacher \textbf{for} fifteen years now. I started teaching in 2009, so I have worked in Bamenda \textbf{since} 2009. \textbf{Ten years ago}, I taught in a small village school near Wum. \textbf{After} I obtained my Higher Teacher's Diploma, I applied for a transfer to the city. \textbf{Afterwards}, I was appointed to Government High School Bamenda. I have lived in this neighbourhood \textbf{for} a long time and I have no plan to leave.”
\end{textebox}

\begin{definition}[Glossaire du texte]
\begin{itemize}
\item \eng{secondary school teacher} (n.) : enseignant du secondaire ;
\item \eng{Higher Teacher's Diploma} (n.) : DIPES (diplôme d'enseignant du secondaire) ;
\item \eng{to apply for a transfer} (v.) : demander une mutation ;
\item \eng{to be appointed to} (v. passif) : être affecté à ;
\item \eng{neighbourhood} (n.) : quartier (orthographe britannique).
\end{itemize}
\end{definition}

\begin{exercice}[Observation et repérage]
\begin{enumerate}
\item Surligne dans le texte les cinq expressions cibles : \eng{for}, \eng{since}, \eng{ago}, \eng{after}, \eng{afterwards}.
\item Pour chacune, indique : le temps du verbe qui l'accompagne (ou qui est dans la même phrase) et la nature de l'expression (durée, point de départ, succession, etc.).
\item Traduis les cinq phrases contenant ces expressions en français.
\end{enumerate}
\end{exercice}

\begin{corrige}[Exercice 1 — Observation]
\begin{enumerate}
\item \eng{I have been a teacher \textbf{for} fifteen years} → Present perfect (\eng{have been}) + \textbf{durée}.
\item \eng{I have worked in Bamenda \textbf{since} 2009} → Present perfect (\eng{have worked}) + \textbf{point de départ}.
\item \eng{\textbf{Ten years ago}, I taught...} → Past simple (\eng{taught}) + \textbf{antériorité datée} (\eng{ago} placé \textbf{après} la durée).
\item \eng{\textbf{After} I obtained my diploma...} → Past simple (\eng{obtained}) dans la proposition introduite par \eng{after} (succession d'événements passés).
\item \eng{\textbf{Afterwards}, I was appointed...} → Past simple (\eng{was appointed}) ; \eng{afterwards} = adverbe de liaison signifiant « ensuite ».
\item Traductions : 1. « Je suis enseignant depuis quinze ans. » 2. « Je travaille à Bamenda depuis 2009. » 3. « Il y a dix ans, j'enseignais… » 4. « Après avoir obtenu mon diplôme… » 5. « Par la suite, j'ai été affecté… »
\end{enumerate}
\end{corrige}

\courssub{For et Since : durée et point de départ (Present Perfect)}

\begin{propriete}[Règle fondamentale : \cle{for} / \cle{since} + Present Perfect]
\begin{itemize}
\item \textbf{\cle{for} + durée} (\eng{for three years}, \eng{for a long time}, \eng{for ages}) : indique la \emph{longueur} de l'action. L'action a commencé dans le passé et \textbf{se poursuit encore au présent}.
  \eng{She has lived in Douala \textbf{for} ten years.} (Elle vit à Douala depuis dix ans / Elle y vit encore.)
\item \textbf{\cle{since} + point de départ} (\eng{since 2016}, \eng{since Monday}, \eng{since I arrived}, \eng{since childhood}) : indique le \emph{moment où l'action a commencé}. L'action se poursuit encore.
  \eng{They have known each other \textbf{since} childhood.} (Ils se connaissent depuis l'enfance.)
\item \textbf{Temps verbal} : \emph{Present Perfect Simple} (\eng{have/has + V-en}) ou \emph{Present Perfect Continuous} (\eng{have/has been + V-ing}) pour insister sur la durée ininterrompue.
  \eng{I have worked here \textbf{for} five years.} = \eng{I have been working here \textbf{for} five years.}
\end{itemize}
\end{propriete}

\begin{exemplebox}[Exemples contrastés \eng{for} / \eng{since}]
\begin{center}
\begin{tabularx}{\linewidth}{|X|X|X|}
\hline
\rowcolor{popblueL} Expression & Phrase anglaise & Traduction française \\
\hline
\cle{for two hours} & We have waited \textbf{for two hours}. & Nous attendons depuis deux heures. \\
\hline
\cle{since 10 a.m.} & We have waited \textbf{since 10 a.m.}. & Nous attendons depuis 10 heures. \\
\hline
\cle{for a week} & He has been sick \textbf{for a week}. & Il est malade depuis une semaine. \\
\hline
\cle{since last Monday} & He has been sick \textbf{since last Monday}. & Il est malade depuis lundi dernier. \\
\hline
\cle{for ages} & I haven't seen you \textbf{for ages}! & Je ne t'ai pas vu depuis des lustres ! \\
\hline
\end{tabularx}
\end{center}
\end{exemplebox}

\begin{attention}[Piège majeur : le piège du « depuis » français]
En français, on dit : « \emph{Je travaille ici depuis 2016} » (présent).\\
En anglais, \textbf{interdit} : \eng{I work here since 2016.} $\times$\\
\textbf{Obligatoire} : \eng{I \textbf{have worked} here \textbf{since} 2016.} (Present Perfect)\\
\textbf{Raison} : L'action a commencé dans le passé (\eng{2016}) et continue \textbf{maintenant}. Le présent simple anglais (\eng{I work}) décrit une habitude, pas une durée inachevée.
\end{attention}

\courssub{Ago : situer un événement terminé dans le passé (Past Simple)}

\begin{propriete}[Règle : \cle{ago} + Past Simple]
\begin{itemize}
\item \textbf{\cle{ago}} signifie « il y a … » (antériorité par rapport au moment de la parole).
\item \textbf{Position} : \textbf{toujours APRÈS} l'expression de durée** (\eng{five minutes ago}, \eng{two centuries ago}, \eng{a long time ago}). Jamais \eng{ago five minutes}.
\item \textbf{Temps verbal} : \emph{Past Simple} (prétérit). L'action est \textbf{terminée} et \textbf{datée}.
  \eng{He left \textbf{ten minutes ago}.} (Il est parti il y a dix minutes.)
  \eng{She started her business \textbf{twenty years ago}.} (Elle a créé son entreprise il y a vingt ans.)
\item \textbf{Repère temporel} : \eng{ago} ancre l'événement dans un passé révolu (souvent avec \eng{yesterday}, \eng{last week}, \eng{in 2010} implicites).
\end{itemize}
\end{propriete}

\begin{exemplebox}[Ago vs For / Since]
\begin{center}
\begin{tabularx}{\linewidth}{|X|X|X|}
\hline
\rowcolor{popblueL} Contexte & Structure correcte & Temps \\
\hline
Action commencée dans le passé, continue maintenant & \eng{I have known him \textbf{for} ten years / \textbf{since} 2014.} & Present Perfect \\
\hline
Action terminée à un moment précis dans le passé & \eng{I met him \textbf{ten years ago}.} & Past Simple \\
\hline
Question « Depuis quand ? » (encore vrai) & \eng{How long \textbf{have you known} him?} & Present Perfect \\
\hline
Question « Quand ? » (événement passé) & \eng{When \textbf{did you meet} him?} & Past Simple \\
\hline
\end{tabularx}
\end{center}
\end{exemplebox}

\courssub{After et Afterwards : exprimer la succession}

\begin{propriete}[Règle : \cle{after} (préposition / conjonction) vs \cle{afterwards} (adverbe)]
\begin{itemize}
\item \textbf{\cle{after} + nom / pronom / gérondif / proposition} : introduit un événement \textbf{antérieur} à un autre.
  \begin{itemize}
  \item \eng{After school}, I play football. (Préposition + nom)
  \item \eng{After saving enough money}, he bought a car. (Préposition + gérondif)
  \item \eng{After he saved enough money}, he bought a car. (Conjonction + sujet + verbe)
  \end{itemize}
  \textbf{Temps de la proposition après \eng{after}} : généralement \emph{Past Simple} (passé) ou \emph{Past Perfect} (antériorité forte : \eng{After he \textbf{had saved} enough money...}).
\item \textbf{\cle{afterwards}} (aussi \eng{afterward} en anglais américain) : adverbe \textbf{indépendant} signifiant « ensuite, par la suite, plus tard ». Il ne introduit pas de proposition.
  \textbf{Place} : en début de phrase (suivi d'une virgule) ou en fin de phrase.
  \eng{He saved money. \textbf{Afterwards}, he bought a car.} ou \eng{He saved money; he bought a car \textbf{afterwards}.}
\end{itemize}
\end{propriete}

\begin{exemplebox}[After vs Afterwards]
\begin{center}
\begin{tabularx}{\linewidth}{|X|X|}
\hline
\rowcolor{popblueL} Phrase & Analyse \\
\hline
\eng{After the meeting, we went home.} & \eng{After} = préposition + nom (\eng{the meeting}). \\
\hline
\eng{After we finished the meeting, we went home.} & \eng{After} = conjonction + proposition (\eng{we finished}). \\
\hline
\eng{We finished the meeting. Afterwards, we went home.} & \eng{Afterwards} = adverbe de liaison (nouvelle phrase). \\
\hline
\eng{We had lunch. We went for a walk afterwards.} & \eng{Afterwards} en fin de phrase. \\
\hline
\end{tabularx}
\end{center}
\end{exemplebox}

\courssub{Comparaison français/anglais et pièges des francophones}

\begin{center}
\begin{tabularx}{\linewidth}{|X|X|X|}
\hline
\rowcolor{popblueL} Français (contexte) & English (structure) & Remarque clé \\
\hline
Depuis trois ans (durée, encore vrai) & \eng{for three years} + Present Perfect & \eng{for} + durée \\
\hline
Depuis 2019 (point de départ, encore vrai) & \eng{since 2019} + Present Perfect & \eng{since} + date / événement \\
\hline
Il y a trois ans (événement terminé) & \eng{three years ago} + Past Simple & \eng{ago} \textbf{après} la durée \\
\hline
Après avoir fini / Après que j'ai fini & \eng{after finishing} / \eng{after I finished} & \eng{after} + V-ing ou proposition \\
\hline
Ensuite / Par la suite (liaison) & \eng{afterwards} (adverbe) & Pas de proposition après \\
\hline
\end{tabularx}
\end{center}

\begin{attention}[Erreurs fréquentes à bannir — Check-list]
\begin{itemize}
\item \textbf{\eng{since} + durée} : $\times$ \eng{since three years} $\rightarrow$ \checkmark \eng{for three years} OU \eng{since 2021}.
\item \textbf{Présent simple avec \eng{since/for}} (calque du français) : $\times$ \eng{I live here since 2010} $\rightarrow$ \checkmark \eng{I \textbf{have lived} here since 2010}.
\item \textbf{\eng{ago} avant la durée} : $\times$ \eng{ago two years} $\rightarrow$ \checkmark \eng{two years ago}.
\item \textbf{Confusion \eng{after} / \eng{afterwards}} : $\times$ \eng{After, I went home} (si pas de nom/proposition après) $\rightarrow$ \checkmark \eng{Afterwards, I went home} OU \eng{After that, I went home}.
\item \textbf{\eng{for} avec passé simple} (action terminée) : $\times$ \eng{I worked there for two years} (si je n'y travaille plus, c'est correct MAIS si j'y travaille encore $\rightarrow$ Present Perfect). \textbf{Test} : « Est-ce que je travaille encore là-bas ? » Oui $\rightarrow$ \eng{I have worked there for two years.} Non $\rightarrow$ \eng{I worked there for two years.}
\end{itemize}
\end{attention}

\courssub{Entraînement guidé : Mise en situation professionnelle — « Working Lives in Cameroon »}

\courssub{Situation de communication}

Tu es journaliste stagiaire à \emph{CRTV Radio} (Yaoundé). Ta rédactrice en chef prépare l'émission \emph{“Working Lives in Cameroon”}. Ta mission : interviewer un travailleur de la capitale (commerçant au Marché Mokolo, enseignant, chauffeur de taxi/benskin, couturier, mécanicien, cultivateur de cacao…) pour raconter son parcours professionnel, de ses débuts à aujourd'hui.

\begin{textebox}[CRTV Radio — Assignment Sheet]
\textbf{Assignment: “Working Lives in Cameroon”}

Dear Reporter,

Your task is to interview a worker in Yaoundé and tell the story of his or her professional life. Your interview must answer these questions:

When did the person start working? How long has he or she been in this job? What did he or she do before? What happened after important moments of his or her career?

Use time expressions such as \emph{for}, \emph{since}, \emph{ago}, \emph{after} and \emph{afterwards} to make the story clear. The best interviews will be broadcast on Friday evening.

Good luck!

The Editor
\end{textebox}

\begin{definition}[Lexique utile pour l'interview]
\begin{itemize}
\item \eng{a shopkeeper / a trader} (n.) : un commerçant ;
\item \eng{a taxi driver / a bike rider (benskin)} (n.) : un chauffeur de taxi / un conducteur de moto-taxi ;
\item \eng{a mechanic} (n.) : un mécanicien ;
\item \eng{a tailor / a seamstress} (n.) : un couturier / une couturière ;
\item \eng{a cocoa farmer} (n.) : un cultivateur de cacao ;
\item \eng{to save money} (v.) : économiser de l'argent ;
\item \eng{to start a business / to set up a business} (v.) : créer une entreprise / se mettre à son compte ;
\item \eng{to earn a living} (v.) : gagner sa vie ;
\item \eng{a customer / a client} (n.) : un client ;
\item \eng{to retire} (v.) : prendre sa retraite ;
\item \eng{apprentice} (n.) : apprenti.
\end{itemize}
\end{definition}

\courssub{Consignes de l'activité (Travail par binômes)}

\begin{enumerate}
\item \textbf{Rôles :} Élève A = Journaliste ; Élève B = Travailleur (choisis un profil). Inversez les rôles ensuite.
\item \textbf{Préparation Journaliste (5 min) :} Rédige \textbf{au moins 6 questions} variées :
  \begin{itemize}
  \item 2 questions avec \eng{How long have you...?} (Present Perfect) — ex. \eng{How long have you been a mechanic?}
  \item 2 questions avec \eng{When did you...?} (Past Simple) — ex. \eng{When did you open this shop?}
  \item 1 question avec \eng{What did you do before...?} (Past Simple) — ex. \eng{What did you do before becoming a driver?}
  \item 1 question avec \eng{What happened after...?} — ex. \eng{What happened after you got your diploma?}
  \end{itemize}
\item \textbf{Préparation Interviewé (5 min) :} Invente ton personnage (nom, âge, métier, dates clés : début d'activité, formation, événement marquant). Prépare des réponses qui intègrent \textbf{chacun des cinq adverbes} au moins une fois : \eng{for}, \eng{since}, \eng{ago}, \eng{after}, \eng{afterwards}.
\item \textbf{L'interview (5 min) :} Joue la scène à l'oral, sans lire mot à mot. Le journaliste prend des notes (mots-clés, dates, adverbes entendus).
\item \textbf{Vérification mutuelle (3 min) :} Relisez les notes ensemble. Chaque adverbe est-il utilisé avec le bon temps ? Corrigez : \eng{since ten years} $\rightarrow$ \eng{for ten years} ; \eng{I work here since...} $\rightarrow$ \eng{I have worked here since...} ; \eng{ago two years} $\rightarrow$ \eng{two years ago}.
\item \textbf{Présentation orale (Classe entière) :} 2 ou 3 binômes jouent leur interview devant la classe. Les autres élèves lèvent la main à chaque occurrence des 5 adverbes cibles.
\item \textbf{Production écrite individuelle (10 min) :} Rédige un paragraphe biographique (8–10 lignes) intitulé \eng{The Story of a Worker}. Contraintes obligatoires :
  \begin{itemize}
  \item Les 5 adverbes (\eng{for}, \eng{since}, \eng{ago}, \eng{after}, \eng{afterwards}) apparaissent au moins une fois.
  \item Au moins 2 verbes au \emph{Present Perfect}.
  \item Au moins 2 verbes au \emph{Past Simple}.
  \item Cohérence chronologique et orthographe soignée.
  \end{itemize}
\end{enumerate}

\courssub{Proposition de production et corrigé commenté}

\begin{exoresolu}[Interview et paragraphe biographique — Modèle complet]
\textbf{Énoncé.} Rédige l'interview complète (questions + réponses) puis le paragraphe biographique de 8 à 10 lignes.

\tcblower

\textbf{The Interview (Modèle)}

\textbf{Journalist:} Good morning, Madam Awah. Thank you for accepting this interview for CRTV Radio. \eng{How long have you been a shopkeeper at Mokolo Market?}

\textbf{Madam Awah:} Good morning. \eng{I have been a shopkeeper \textbf{for} twelve years now.}

\textbf{Journalist:} \eng{When did you start this business?}

\textbf{Madam Awah:} \eng{I started in 2012, so I have run this shop \textbf{since} 2012.}

\textbf{Journalist:} \eng{What did you do \textbf{before} you became a trader?}

\textbf{Madam Awah:} \eng{\textbf{Fifteen years ago}, I was an apprentice tailor in my aunt's workshop in Bafoussam. I learned sewing \textbf{for} three years.}

\textbf{Journalist:} \eng{What happened \textbf{after} your apprenticeship?}

\textbf{Madam Awah:} \eng{\textbf{After} I obtained my certificate, I moved to Yaoundé to join my brother. \textbf{Afterwards}, I saved money for two years and opened this shop.}

\textbf{Journalist:} \eng{How long have you owned this specific location?}

\textbf{Madam Awah:} \eng{I have had this spot \textbf{since} 2015. The business is growing well, thank God.}

\textbf{Journalist:} Thank you very much, Madam Awah.

\vspace{0.5cm}

\textbf{The Biographical Paragraph (Modèle)}

\emph{Madam Awah is a successful shopkeeper at Mokolo Market in Yaoundé. She has run her business \textbf{for} twelve years. She started in 2012, so she has worked at this location \textbf{since} 2012, although she has owned this specific spot \textbf{since} 2015. \textbf{Fifteen years ago}, she was an apprentice tailor in Bafoussam, where she learned sewing \textbf{for} three years. \textbf{After} she obtained her certificate, she moved to the capital. \textbf{Afterwards}, she saved money for two years before opening her own shop. Today, she earns a good living and trains two apprentices herself.}

\textbf{Commentaire linguistique détaillé}
\begin{itemize}
\item \eng{for twelve years} / \eng{for three years} / \eng{for two years} : \textbf{durée} $\rightarrow$ Present Perfect (\eng{has run}, \eng{learned} — \textit{note: learned est passé simple ici car l'apprentissage est terminé, mais "saved for two years" est une durée passée terminée $\rightarrow$ Past Simple correct}). \textbf{Attention} : \eng{I learned sewing for three years} (action passée terminée) vs \eng{I have been a shopkeeper for twelve years} (action qui continue).
\item \eng{since 2012} / \eng{since 2015} : \textbf{point de départ} $\rightarrow$ Present Perfect (\eng{has run}, \eng{has worked}, \eng{has owned}). Correct.
\item \eng{Fifteen years ago, I was...} : \textbf{ago} + \textbf{Past Simple} (\eng{was}, \eng{learned}). Position de \eng{ago} \textbf{après} la durée. Correct.
\item \eng{After I obtained...} : \eng{After} + proposition au Past Simple (antériorité simple). On pourrait aussi écrire \eng{After obtaining...} (gérondif). Correct.
\item \eng{Afterwards, I saved...} : \eng{Afterwards} = adverbe de liaison en tête de phrase, suivi d'une virgule. Correct.
\item Le paragraphe alterne \textbf{Present Perfect} (statut actuel : \eng{has run}, \eng{has worked}, \eng{has owned}, \eng{has had}) et \textbf{Past Simple} (événements datés et terminés : \eng{was}, \eng{learned}, \eng{obtained}, \eng{moved}, \eng{saved}, \eng{opened}). C'est la marque d'un récit biographique bien construit.
\end{itemize}
\end{exoresolu}

\courssub{Exercices d'application supplémentaires}

\begin{exercice}[Mise en forme : choisis la bonne forme]
\begin{enumerate}
\item My father \_\_\_\_\_ (work) in this company \_\_\_\_\_ 2005. (since / for)
\item We \_\_\_\_\_ (arrive) in Yaoundé \_\_\_\_\_ (two days ago / since two days).
\item \_\_\_\_\_ (After / Afterwards) finishing his studies, he \_\_\_\_\_ (find) a job.
\item She \_\_\_\_\_ (be) a nurse \_\_\_\_\_ (for / since) twenty years. She loves her job.
\item \_\_\_\_\_ (Ago / Since) ten years, I \_\_\_\_\_ (start) learning English.
\item They lived in Douala \_\_\_\_\_ (for / since) five years, then they moved. (Action terminée)
\item He left school in 2010. \_\_\_\_\_ (After / Afterwards), he became a driver.
\end{enumerate}
\end{exercice}

\begin{corrige}[Exercice 2]
\begin{enumerate}
\item My father \textbf{has worked} in this company \textbf{since} 2005. (Point de départ + Present Perfect)
\item We \textbf{arrived} in Yaoundé \textbf{two days ago}. (Passé simple + \eng{ago} après durée)
\item \textbf{After} finishing his studies, he \textbf{found} a job. (\eng{After} + V-ing / Past Simple)
\item She \textbf{has been} a nurse \textbf{for} twenty years. (Durée + Present Perfect — elle l'est encore)
\item \textbf{Ten years ago}, I \textbf{started} learning English. (\eng{ago} après durée + Past Simple)
\item They lived in Douala \textbf{for} five years, then they moved. (Durée passée terminée $\rightarrow$ Past Simple + \eng{for})
\item He left school in 2010. \textbf{Afterwards}, he became a driver. (Adverbe de liaison = \eng{Afterwards})
\end{enumerate}
\end{corrige}

\begin{exercice}[Transformation : réécris les phrases en corrigeant l'erreur]
\begin{enumerate}
\item I am a teacher since 2018.
\item He works here for ten years.
\item She came ago five minutes.
\item After to finish, we went out.
\item I have seen him afterwards.
\end{enumerate}
\end{exercice}

\begin{corrige}[Exercice 3]
\begin{enumerate}
\item \textbf{I have been a teacher since 2018.} (Present Perfect obligatoire avec \eng{since})
\item \textbf{He has worked here for ten years.} (Present Perfect car action qui continue ; \eng{for} pour la durée)
\item \textbf{She came five minutes ago.} (\eng{ago} après la durée)
\item \textbf{After finishing,} we went out. (\eng{After} + gérondif, pas \eng{to finish})
\item \textbf{I saw him afterwards.} (\eng{Afterwards} s'emploie au passé simple pour un événement passé ; \eng{I have seen him} impliquerait un lien avec le présent, mais \eng{afterwards} ancre dans la narration passée).
\end{enumerate}
\end{corrige}

\courssub{Bilan de la leçon}

Cette leçon t'a permis de maîtriser les cinq marqueurs temporels essentiels pour structurer un récit de vie ou une biographie professionnelle en anglais. Retiens l'essentiel :

\begin{itemize}
\item \textbf{\cle{for} + durée} et \textbf{\cle{since} + point de départ} $\rightarrow$ \textbf{Present Perfect} (action non terminée, lien avec le présent). \textit{Français : « depuis » + présent.}
\item \textbf{\cle{ago} + durée (après)} $\rightarrow$ \textbf{Past Simple} (action terminée, datée). \textit{Français : « il y a » + passé composé / imparfait.}
\item \textbf{\cle{after} + nom / V-ing / proposition} $\rightarrow$ Préposition / Conjonction de succession. Verbe souvent au Past Simple ou Past Perfect.
\item \textbf{\cle{afterwards}} $\rightarrow$ \textbf{Adverbe} autonome = « ensuite, par la suite ». Pas de proposition après.
\item \textbf{Différence clé FR/EN} : \eng{« Je travaille ici depuis 2016 »} $\neq$ \eng{I work here since 2016} $\times$ $\rightarrow$ \checkmark \eng{I \textbf{have worked} here since 2016.}
\end{itemize}

\begin{methode}[Pour réussir ton paragraphe biographique (Bac / Évaluation)]
\begin{enumerate}
\item \textbf{Chronologie} : Note les dates clés sur un brouillon (Début, Formation, Événement marquant, Situation actuelle).
\item \textbf{Codage temps} : \textcolor{popblue}{Bleu} = Present Perfect (maintenant) ; \textcolor{poporange}{Orange} = Past Simple (passé terminé).
\item \textbf{Insertion} : Place \eng{for}/\eng{since} dans les phrases « Bleu » ; \eng{ago}/\eng{after} dans les phrases « Orange » ; \eng{afterwards} pour relier deux phrases « Orange ».
\item \textbf{Relecture} : Vérifie \eng{ago} (après durée), \eng{since} (date), \eng{for} (durée), \eng{after} vs \eng{afterwards}.
\end{enumerate}
\end{methode}

\begin{savaistu}[Le saviez-vous ? — Anglais camerounais et marques de temps]
Au Cameroun anglophone (régions du Nord-Ouest et Sud-Ouest), on entend souvent des structures influencées par le pidgin ou le français : \eng{« I am here since morning »} (au lieu de \eng{I have been here since morning}) ou \eng{« Since long time I no see you »} (pidgin). À l'écrit formel (bac, administration, presse \emph{The Guardian Post}, \emph{Cameroon Tribune}), seule la norme standard britannique (\eng{I have been here since morning}) est acceptée. Maîtriser \eng{for/since/ago} est donc un marqueur de compétence \textbf{académique et professionnelle}.
\end{savaistu}

\newpage

\coursec{Méthodes et exercices résolus}

\courssub{Savoir-faire 1 : Choisir entre \eng{for} et \eng{since} pour exprimer la durée}

\begin{methode}[Choisir entre \eng{for} et \eng{since}]
\begin{enumerate}
\item \textbf{Repérer le complément de temps.} Demande-toi : est-ce une \cle{durée} (un intervalle de temps : \eng{three years, six months, a long time}) ou un \cle{point de départ} (une date, un événement, un moment précis : \eng{2019, Monday, Christmas, he left school}) ?
\item \textbf{Appliquer la règle.} Durée $\rightarrow$ \eng{for} ; point de départ $\rightarrow$ \eng{since}.
\item \textbf{Vérifier le temps du verbe.} Avec \eng{for} et \eng{since}, quand l'action continue jusqu'au présent, l'anglais exige le \cle{present perfect} (ou present perfect continuous) : \eng{She has worked here for ten years.} — jamais le présent simple comme en français (« elle travaille ici depuis dix ans »).
\item \textbf{Traduire mentalement.} \eng{for + durée} = « depuis / pendant » ; \eng{since + point de départ} = « depuis ». Si le français dit « depuis », teste : « depuis quand ? » $\rightarrow$ \eng{since} ; « depuis combien de temps ? » $\rightarrow$ \eng{for}.
\item \textbf{Cas particulier.} À la forme négative, le present perfect avec \eng{for/since} traduit souvent « cela fait ... que ... ne ... pas » : \eng{I haven't seen him since 2020.} « Je ne l'ai pas vu depuis 2020. »
\end{enumerate}
\end{methode}

\begin{propriete}[Récapitulatif : \eng{for} ou \eng{since} ?]
\begin{center}
\begin{tabularx}{\linewidth}{|X|X|X|}\hline
\rowcolor{popblueL} \eng{for} + durée & \eng{since} + point de départ & Exemple \\ \hline
\eng{for two hours} & \eng{since 8 o'clock} & \eng{The shop has been open for two hours.} \\ \hline
\eng{for three weeks} & \eng{since Monday} & \eng{He has been ill since Monday.} \\ \hline
\eng{for ten years} & \eng{since 2014} & \eng{They have lived in Yaoundé for ten years.} \\ \hline
\eng{for a long time} & \eng{since her childhood} & \eng{She has loved music since her childhood.} \\ \hline
\eng{for ages} & \eng{since he left school} & \eng{He has been a driver since he left school.} \\ \hline
\end{tabularx}
\end{center}
Remarque : \eng{since} peut être suivi d'une proposition au prétérit : \eng{since he left school} « depuis qu'il a quitté l'école ».
\end{propriete}

\begin{exoresolu}[Compléter avec \eng{for} ou \eng{since}]
Complete the sentences with \eng{for} or \eng{since}.\\
1. Mr Etoundi has been a cocoa farmer \ldots\ twenty-five years.\\
2. My aunt has lived in Douala \ldots\ 2015.\\
3. We have known each other \ldots\ primary school.\\
4. The bank has been closed \ldots\ two hours.\\
5. She has not eaten anything \ldots\ this morning.
\tcblower
\textbf{Raisonnement et solutions :}
\begin{enumerate}
\item \eng{for twenty-five years} — « twenty-five years » est une \textbf{durée} (combien de temps ?). Traduction : « M. Etoundi est planteur de cacao depuis vingt-cinq ans. »
\item \eng{since 2015} — « 2015 » est une \textbf{date}, donc un point de départ. « Ma tante habite à Douala depuis 2015. »
\item \eng{since primary school} — « primary school » est un \textbf{moment/événement} de référence, pas une durée. « Nous nous connaissons depuis l'école primaire. »
\item \eng{for two hours} — « two hours » est une durée. « La banque est fermée depuis deux heures. »
\item \eng{since this morning} — « this morning » est un point de départ précis. « Elle n'a rien mangé depuis ce matin. »
\end{enumerate}
\textbf{Justification grammaticale :} dans les cinq phrases, le verbe est au present perfect (\eng{has been, has lived, have known, has been, has not eaten}) parce que la situation commencée dans le passé se poursuit jusqu'au moment où l'on parle. Le présent simple français (« elle habite depuis 2015 ») est un piège : en anglais, \eng{*She lives here since 2015} est incorrect.
\end{exoresolu}

\courssub{Savoir-faire 2 : Employer \eng{ago} pour situer un événement dans le passé}

\begin{methode}[Construire une phrase avec \eng{ago}]
\begin{enumerate}
\item \textbf{Comprendre le sens.} \eng{ago} = « il y a » : il situe un événement \cle{terminé} en remontant dans le passé à partir d'\textbf{aujourd'hui}.
\item \textbf{Respecter la place.} \eng{ago} se place \textbf{après} la durée : \eng{two days ago} « il y a deux jours » — jamais \eng{*ago two days}.
\item \textbf{Choisir le bon temps.} \eng{ago} s'emploie toujours avec le \cle{prétérit simple} (past simple), jamais avec le present perfect : \eng{He arrived three days ago.} et non \eng{*He has arrived three days ago.}
\item \textbf{Traduire « il y a » correctement.} « Il y a deux ans que j'habite ici » (action qui continue) $\rightarrow$ \eng{I have lived here for two years} ; « Je suis arrivé il y a deux ans » (action terminée) $\rightarrow$ \eng{I arrived two years ago.}
\item \textbf{Interroger.} Pour demander « il y a combien de temps ? », on dit \eng{How long ago ...?} avec le prétérit : \eng{How long ago did he leave?}
\end{enumerate}
\end{methode}

\begin{attention}[\eng{ago} ou \eng{before} ?]
\eng{ago} compte à partir d'\textbf{aujourd'hui} : \eng{I met her two years ago} « je l'ai rencontrée il y a deux ans ». \eng{before} compte à partir d'un \textbf{autre moment du passé} (discours indirect, récit) : \eng{She said she had met him two years before} « elle dit qu'elle l'avait rencontré deux ans auparavant ». Au Bac, dans la transformation du discours direct en discours indirect, \eng{ago} devient \eng{before} : \eng{“I bought it a week ago,” he said} $\rightarrow$ \eng{He said he had bought it a week before.}
\end{attention}

\begin{exoresolu}[Traduire en anglais]
Traduisez en anglais.\\
1. Mon père a acheté cette moto il y a cinq ans.\\
2. Il y a combien de temps as-tu rencontré le directeur ?\\
3. Nous avons déménagé à Buea il y a longtemps.\\
4. Elle a obtenu son diplôme il y a quelques mois.
\tcblower
\textbf{Solutions détaillées :}
\begin{enumerate}
\item \eng{My father bought this motorbike five years ago.} — Événement terminé, situé dans le passé $\rightarrow$ prétérit \eng{bought} (verbe irrégulier \eng{buy -- bought -- bought}) + durée suivie de \eng{ago}. Erreur à éviter : \eng{*has bought ... ago}.
\item \eng{How long ago did you meet the manager?} — Question au prétérit : auxiliaire \eng{did} + base verbale \eng{meet}. Ne jamais dire \eng{*did you met}.
\item \eng{We moved to Buea a long time ago.} — \eng{a long time ago} = « il y a longtemps ». Verbe régulier : \eng{move -- moved}.
\item \eng{She got her diploma a few months ago.} — \eng{a few months ago} = « il y a quelques mois » ; \eng{get -- got -- got} (ici « obtenir »).
\end{enumerate}
\textbf{Point de méthode :} dans chaque phrase, « il y a » renvoie à un événement ponctuel et achevé $\rightarrow$ prétérit + \eng{ago} en fin de groupe de durée. Si l'action continuait jusqu'à aujourd'hui, il faudrait le present perfect avec \eng{for/since}.
\end{exoresolu}

\courssub{Savoir-faire 3 : Employer \eng{after} et \eng{afterwards} pour ordonner les événements}

\begin{methode}[Distinguer \eng{after} et \eng{afterwards}]
\begin{enumerate}
\item \textbf{Identifier la nature du mot.} \eng{after} est une \cle{préposition} ou une \cle{conjonction} : il doit être suivi d'un nom (\eng{after lunch}), d'un pronom (\eng{after that}) ou d'une proposition (\eng{after he left}). \eng{afterwards} est un \cle{adverbe} : il se suffit à lui-même et signifie « ensuite, après ».
\item \textbf{Tester la phrase.} Si un complément suit, utilise \eng{after} ; si rien ne suit, utilise \eng{afterwards} : \eng{We had lunch and afterwards we went back to work.}
\item \textbf{Placer \eng{afterwards}.} Généralement en début ou en fin de phrase : \eng{Afterwards, she apologised.} / \eng{She apologised afterwards.}
\item \textbf{Attention au temps après \eng{after}.} Quand \eng{after} introduit une proposition subordonnée de temps, on n'utilise \textbf{pas} le futur : \eng{I will call you after I finish work} (et non \eng{*after I will finish}).
\item \textbf{Variantes utiles.} \eng{soon afterwards} « peu après », \eng{shortly after} « peu de temps après », \eng{the day after} « le lendemain ».
\end{enumerate}
\end{methode}

\begin{exoresolu}[Choisir entre \eng{after} et \eng{afterwards}]
Complete with \eng{after} or \eng{afterwards}.\\
1. The workers rested \ldots\ the harvest.\\
2. We visited the factory; \ldots\ , the manager offered us some tea.\\
3. \ldots\ the meeting ended, everyone went home.\\
4. He felt sick and went to hospital shortly \ldots\ .
\tcblower
\textbf{Solutions détaillées :}
\begin{enumerate}
\item \eng{after the harvest} — \eng{after} est suivi du nom \eng{the harvest} (préposition). « Les ouvriers se sont reposés après la récolte. »
\item \eng{afterwards} — rien ne suit le mot : il faut l'adverbe. « Nous avons visité l'usine ; ensuite, le directeur nous a offert du thé. »
\item \eng{After the meeting ended} — \eng{after} est ici conjonction, suivi d'une proposition (\eng{the meeting ended}). « Après la fin de la réunion, tout le monde est rentré chez soi. »
\item \eng{afterwards} — \eng{shortly afterwards} = « peu après » ; pas de complément après le mot. « Il s'est senti malade et est allé à l'hôpital peu après. » (En anglais britannique, on dit \eng{to go to hospital}, sans article, quand on y va comme patient.)
\end{enumerate}
\textbf{Justification :} le test décisif est la présence ou non d'un complément : \eng{after} ne peut jamais terminer une phrase seul ; \eng{afterwards} ne peut jamais être suivi d'un nom. Dire \eng{*afterwards the harvest} ou \eng{*we rested after} (sans complément) est incorrect.
\end{exoresolu}

\courssub{Savoir-faire 4 : Traduire « depuis », « il y a », « après » sans calquer le français}

\begin{methode}[Éviter les calques du français]
\begin{enumerate}
\item \textbf{Identifier le sens français exact.} « Depuis » peut être \eng{for} (durée) ou \eng{since} (point de départ) ; « il y a » peut être \eng{ago} (passé) ou \eng{there is/are} (existence) ; « après » peut être \eng{after} ou \eng{afterwards}.
\item \textbf{Choisir le temps anglais en conséquence.} « Depuis » + action qui continue $\rightarrow$ present perfect ; « il y a » + action terminée $\rightarrow$ past simple.
\item \textbf{Refuser les calques.} Ne jamais traduire « depuis » par \eng{*during} ou par \eng{for} devant une date (\eng{*for 2015}), ni « il y a deux jours » par \eng{*there is two days}.
\item \textbf{Vérifier l'ordre des mots.} En anglais, la durée avec \eng{ago} se place après le verbe et son complément : \eng{I saw her two weeks ago.} L'adverbe ne se coince pas entre le verbe et son complément direct.
\item \textbf{Relire en posant la question test.} « Combien de temps ? » $\rightarrow$ \eng{for} ; « à partir de quand ? » $\rightarrow$ \eng{since} ; « quand exactement dans le passé ? » $\rightarrow$ \eng{ago}.
\end{enumerate}
\end{methode}

\begin{center}
\begin{tabularx}{\linewidth}{|X|X|X|}\hline
\rowcolor{popblueL} Français & English & Remarque \\ \hline
Elle travaille ici depuis dix ans. & \eng{She has worked here for ten years.} & durée $\rightarrow$ \eng{for} + present perfect \\ \hline
Elle travaille ici depuis 2015. & \eng{She has worked here since 2015.} & date $\rightarrow$ \eng{since} + present perfect \\ \hline
Elle est arrivée il y a dix ans. & \eng{She arrived ten years ago.} & événement terminé $\rightarrow$ prétérit + \eng{ago} \\ \hline
Il y a un cybercafé près du lycée. & \eng{There is a cybercafé near the school.} & existence $\rightarrow$ \eng{there is} \\ \hline
Après le travail, il est rentré. & \eng{After work, he went home.} & \eng{after} + nom \\ \hline
Il a déjeuné, puis il est sorti. & \eng{He had lunch and went out afterwards.} & \eng{afterwards} seul, en fin de phrase \\ \hline
\end{tabularx}
\end{center}

\begin{exoresolu}[Corriger les phrases fautives]
Each sentence contains one mistake made by a French-speaking learner. Correct it.\\
1. I live in Bamenda since 2018.\\
2. She has left the office two hours ago.\\
3. He works here for 2020.\\
4. Afterwards the exam, we celebrated.\\
5. There is three days, I met her at the market.
\tcblower
\textbf{Corrections et explications :}
\begin{enumerate}
\item \eng{I have lived in Bamenda since 2018.} — Calque du présent français (« j'habite depuis 2018 ») : l'anglais exige le present perfect pour une action qui dure encore.
\item \eng{She left the office two hours ago.} — \eng{ago} est incompatible avec le present perfect : événement achevé $\rightarrow$ prétérit \eng{left}.
\item \eng{He has worked here since 2020.} — « 2020 » est une date (point de départ), donc \eng{since}, pas \eng{for}. De plus, l'action continue $\rightarrow$ present perfect.
\item \eng{After the exam, we celebrated.} — \eng{afterwards} est un adverbe et ne peut pas introduire un nom ; il faut la préposition \eng{after}.
\item \eng{Three days ago, I met her at the market.} (ou \eng{I met her at the market three days ago.}) — « Il y a trois jours » se traduit par \eng{three days ago} ; \eng{there is} signifie « il y a » au sens d'existence (\eng{There is a book on the table}).
\end{enumerate}
\textbf{Conclusion :} chaque erreur vient d'un calque direct du français. Le réflexe salvateur : identifier d'abord si l'action est terminée (prétérit + \eng{ago}) ou encore en cours (present perfect + \eng{for/since}).
\end{exoresolu}

\courssub{Savoir-faire 5 : Rédiger un court texte en utilisant les adverbes de temps}

\begin{methode}[Rédiger un paragraphe avec \eng{for, since, ago, after, afterwards}]
\begin{enumerate}
\item \textbf{Planifier la chronologie.} Note trois repères : un point de départ (\eng{since ...}), une durée (\eng{for ...}), un événement passé (\eng{... ago}).
\item \textbf{Ordonner le récit.} Utilise \eng{after + nom/proposition} pour enchaîner deux faits et \eng{afterwards} pour conclure l'étape suivante.
\item \textbf{Accorder les temps.} Present perfect pour ce qui dure encore, prétérit pour ce qui est terminé.
\item \textbf{Varier les positions.} Place \eng{afterwards} en début de phrase avec une virgule pour dynamiser le récit : \eng{Afterwards, they opened a shop in town.}
\item \textbf{Relire avec la grille.} \eng{for} + durée ? \eng{since} + point de départ ? \eng{ago} + prétérit ? \eng{after} suivi d'un complément ? Coche chaque point avant de rendre la copie.
\end{enumerate}
\end{methode}

\begin{exoresolu}[Production écrite guidée]
Write a short paragraph (5--6 lines) about a trader in Douala, using \eng{for, since, ago, after} and \eng{afterwards} at least once each.
\tcblower
\textbf{Démarche :} on fixe d'abord les repères chronologiques (début de l'activité, durée, événement marquant), puis on rédige en respectant les temps verbaux.\\
\textbf{Modèle de réponse :}\\
\eng{Mrs Ngo Bell has sold vegetables at the Deido market for fifteen years. She has worked there since 2009, when she left her village. Three years ago, she saved enough money to rent a small shop. After the market closes every evening, she counts her money carefully. Afterwards, she goes home to cook for her children.}\\
« Mme Ngo Bell vend des légumes au marché Deido depuis quinze ans. Elle y travaille depuis 2009, quand elle a quitté son village. Il y a trois ans, elle a économisé assez d'argent pour louer une petite boutique. Chaque soir, après la fermeture du marché, elle compte soigneusement son argent. Ensuite, elle rentre à la maison cuisiner pour ses enfants. »\\
\textbf{Vérification :} \eng{for fifteen years} (durée, present perfect) ; \eng{since 2009} (point de départ, present perfect) ; \eng{three years ago} (prétérit \eng{saved}) ; \eng{after the market closes} (conjonction + proposition au présent, pas de futur) ; \eng{afterwards} (adverbe en tête de phrase). Toutes les règles sont respectées.
\end{exoresolu}

\begin{attention}[Les 5 erreurs qui coûtent des points au Bac]
\begin{enumerate}
\item \textbf{Présent simple au lieu du present perfect.} \eng{*I live here since 2019} est un calque de « j'habite ici depuis 2019 ». Correct : \eng{I have lived here since 2019.}
\item \textbf{Present perfect avec \eng{ago}.} \eng{*He has arrived two days ago} est impossible : \eng{ago} impose le prétérit $\rightarrow$ \eng{He arrived two days ago.}
\item \textbf{Confondre \eng{for} et \eng{since}.} \eng{*for 2015} est faux : une date exige \eng{since}. Retiens : \eng{for} + durée, \eng{since} + point de départ.
\item \textbf{Employer \eng{afterwards} comme préposition.} \eng{*afterwards the lesson} est incorrect : \eng{afterwards} est un adverbe sans complément. Dis \eng{after the lesson}.
\item \textbf{Futur après \eng{after} (conjonction de temps).} \eng{*I will call you after I will finish} est fautif : après \eng{after}, on met le présent $\rightarrow$ \eng{I will call you after I finish.}
\end{enumerate}
\end{attention}

\begin{aretenir}[L'essentiel de la leçon]
\begin{itemize}
\item \eng{for} + \textbf{durée} (\eng{for two years}) ; \eng{since} + \textbf{point de départ} (\eng{since 2020, since Monday, since he arrived}). Les deux s'emploient avec le \textbf{present perfect} quand l'action continue.
\item \eng{ago} = « il y a » : placé \textbf{après} la durée, toujours avec le \textbf{prétérit} : \eng{He left an hour ago.}
\item \eng{after} = préposition ou conjonction, toujours \textbf{suivi} d'un complément ; \eng{afterwards} = adverbe, \textbf{sans} complément : « ensuite ».
\item Pas de futur après les conjonctions de temps (\eng{after, when, before}) : \eng{I will rest after I finish.}
\item Question test : « combien de temps ? » $\rightarrow$ \eng{for} ; « depuis quand ? » $\rightarrow$ \eng{since} ; « quand dans le passé ? » $\rightarrow$ \eng{ago}.
\end{itemize}
\end{aretenir}

\begin{exercice}[Pour s'entraîner]
1. Complete with \eng{for}, \eng{since} or \eng{ago} : a) They have been married \ldots\ twenty years. b) The train left ten minutes \ldots\ . c) I have not heard from her \ldots\ last Christmas. d) He has been unemployed \ldots\ the factory closed.\\
2. Put the verbs in the correct tense : a) She (work) in this bank since 2016. b) We (see) that film two weeks ago. c) I will phone you after I (arrive).\\
3. Translate into English : a) « J'attends le bus depuis vingt minutes. » b) « Il a quitté le bureau il y a une heure. » c) « Après les cours, nous sommes allés au stade ; ensuite, nous avons mangé. »
\end{exercice}

\begin{corrige}[Exercice « Pour s'entraîner »]
\textbf{1.} a) \eng{for} (durée) ; b) \eng{ago} (événement passé, prétérit \eng{left}) ; c) \eng{since} (point de départ) ; d) \eng{since} (proposition au prétérit : point de départ).\\
\textbf{2.} a) \eng{She has worked in this bank since 2016.} (present perfect, action qui continue) ; b) \eng{We saw that film two weeks ago.} (prétérit avec \eng{ago}) ; c) \eng{I will phone you after I arrive.} (présent après \eng{after}, jamais de futur).\\
\textbf{3.} a) \eng{I have been waiting for the bus for twenty minutes.} (present perfect continuous, action en cours) ; b) \eng{He left the office an hour ago.} ; c) \eng{After school, we went to the stadium; afterwards, we had a meal.}
\end{corrige}

\newpage

\coursec{Exercices}

\courssub{Je vérifie mes connaissances}

\begin{exercice}[QCM : le bon adverbe]
Choisis la bonne réponse pour chaque phrase.
\begin{enumerate}
\item He has been unemployed \cle{...} he lost his job.\par
a) for \quad b) since \quad c) ago \quad d) after
\item My aunt has worked at the central market \cle{...} ten years.\par
a) since \quad b) ago \quad c) for \quad d) afterwards
\item The train left Douala two hours \cle{...}.\par
a) since \quad b) for \quad c) ago \quad d) after
\item We finished the meeting, and \cle{...} we went to a restaurant.\par
a) ago \quad b) since \quad c) for \quad d) afterwards
\item She has lived in Buea \cle{...} 2015.\par
a) for \quad b) since \quad c) ago \quad d) after
\end{enumerate}
\end{exercice}

\begin{exercice}[Vrai ou faux ?]
Dis si chaque phrase est correcte (\emph{true}) ou incorrecte (\emph{false}), puis justifie ta réponse en citant la règle de grammaire concernée.
\begin{enumerate}
\item \eng{I have been a trader since five years.}
\item \eng{They arrived in Bamenda three days ago.}
\item \eng{She has been ill since Monday.}
\item \eng{After the lesson, afterwards we played football.}
\item \eng{We have known each other for a long time.}
\item \eng{He has bought a new motorbike last month.}
\end{enumerate}
\end{exercice}

\begin{exercice}[Vocabulaire : durée ou point de départ ?]
Classe les expressions suivantes en deux colonnes : celles qui expriment une \cle{durée} (et s'emploient avec \eng{for}) et celles qui expriment un \cle{point de départ} (et s'emploient avec \eng{since}).
\begin{center}
\eng{ten years — 2010 — last January — a long time — my childhood — six months — Christmas — two weeks — I was born — ages}
\end{center}
\end{exercice}

\begin{exercice}[Conjugaison : le bon temps]
Mets le verbe entre parenthèses au temps qui convient (\emph{present perfect} ou \emph{simple past}).
\begin{enumerate}
\item I (meet) \cle{...} my business partner three years ago.
\item Since then, we (build) \cle{...} a successful company.
\item My uncle (work) \cle{...} at the post office since 2004.
\item She (leave) \cle{...} her village a long time ago.
\item They (not / see) \cle{...} each other since the end of the school year.
\item After he (finish) \cle{...} his studies, he opened a shop in Douala.
\end{enumerate}
\end{exercice}

\courssub{Je m'entraîne}

\begin{exercice}[Texte à trous : for, since ou ago]
Complète le texte suivant avec \eng{for}, \eng{since} ou \eng{ago}.

\eng{My uncle has been a civil servant \cle{...} twenty years. He got this job in 2004, so he has worked in Yaoundé \cle{...} 2004. He left his village twenty years \cle{...}. He has been married \cle{...} fifteen years, and he has lived in the same house \cle{...} his wedding. Two months \cle{...}, he was promoted, and he has been much happier \cle{...} then.}
\end{exercice}

\begin{exercice}[Traduction]
Traduis les phrases suivantes en anglais.
\begin{enumerate}
\item Mon père est chauffeur de taxi depuis vingt ans.
\item Elle a ouvert sa boutique il y a trois mois.
\item Nous habitons à Douala depuis 2018.
\item Après l'examen, nous sommes rentrés à la maison ; ensuite, nous avons fêté cela.
\item Je n'ai pas vu mon ami depuis la semaine dernière.
\item Ils attendent le car depuis une heure.
\end{enumerate}
\end{exercice}

\begin{exercice}[Transformation]
Transforme chaque phrase selon la consigne, sans changer le sens.
\begin{enumerate}
\item \eng{She has taught English for fifteen years.} $\rightarrow$ \eng{She has taught English since \cle{...}.}
\item \eng{He started working at the bank in 2012.} $\rightarrow$ \eng{He has worked at the bank since \cle{...}.}
\item \eng{They moved to Bamenda five years ago.} $\rightarrow$ \eng{They have lived in Bamenda for \cle{...}.}
\item \eng{I last saw her in June.} $\rightarrow$ \eng{I haven't seen her since \cle{...}.}
\item \eng{We finished lunch and then we visited the market.} $\rightarrow$ \eng{After lunch, \cle{...}.}
\end{enumerate}
\end{exercice}

\begin{exercice}[Appariements : la phrase et sa suite logique]
Associe chaque début de phrase (1--6) à sa suite logique (a--f).
\begin{enumerate}
\item \eng{The farmers harvested the cocoa two weeks ago,}
\item \eng{She has been a nurse}
\item \eng{He has saved money}
\item \eng{After the customers left,}
\item \eng{They signed the contract,}
\item \eng{My brother graduated last year,}
\end{enumerate}
\begin{itemize}
\item[a)] \eng{for more than ten years.}
\item[b)] \eng{and afterwards they celebrated with their partners.}
\item[c)] \eng{and they have already sold it at the market.}
\item[d)] \eng{since he opened his small business.}
\item[e)] \eng{the shopkeeper counted the day's money.}
\item[f)] \eng{and he has been looking for a job since then.}
\end{itemize}
\end{exercice}

\courssub{Je me prépare au Bac}

\begin{exercice}[Compréhension et langue — type Bac]
Lis le texte suivant, puis réponds aux questions en anglais, par des phrases complètes.

\begin{textebox}[A family business in Douala]
Mrs Etame has run a small restaurant near the Douala central market for almost twenty-five years. She opened it in the year 2000, shortly after she lost her job at a textile factory. At first, business was slow, but afterwards more and more customers started coming for her famous \emph{ndolé}. Her husband joined the business ten years ago, and since then they have worked together every day, from six in the morning until late in the evening. Their eldest daughter has helped them since she finished secondary school, and she now manages the accounts. “We have never regretted our decision,” Mrs Etame says. “This restaurant has fed our family and paid for our children's education.”
\end{textebox}

\textbf{A. Comprehension questions}
\begin{enumerate}
\item How long has Mrs Etame run her restaurant?
\item When did she open the restaurant, and why?
\item What happened after her husband joined the business?
\item What has their daughter done since she finished secondary school?
\end{enumerate}

\textbf{B. Language work}
\begin{enumerate}
\item Find in the text two adverbs or expressions of time built with \eng{for}, \eng{since}, \eng{ago}, \eng{after} or \eng{afterwards}, and explain the grammar rule each one illustrates.
\item Complete: \eng{Mrs Etame has run her restaurant \cle{...} the year 2000. She lost her factory job twenty-five years \cle{...}.}
\item Put the verbs in the correct tense: \eng{Since she (open) \cle{...} the restaurant, many customers (come) \cle{...} to taste her cooking.}
\end{enumerate}
\end{exercice}

\begin{exercice}[Traduction — type Bac]
Traduis en anglais le passage suivant (environ 60 mots). Tu veilleras au choix des adverbes de temps et des temps verbaux.

« Mon oncle est menuisier depuis trente ans. Il a appris ce métier il y a très longtemps, quand il était encore adolescent. Après son apprentissage, il a ouvert un petit atelier à Bafoussam ; ensuite, il a formé plusieurs jeunes du quartier. Depuis 2010, il fabrique aussi des meubles pour les écoles de la ville. Il n'a jamais regretté son choix. »
\end{exercice}

\begin{exercice}[Expression écrite — type Bac]
Write a composition of about \textbf{150--180 words} entitled \textbf{“My father's (or mother's) occupation”}.

Your composition must:
\begin{enumerate}
\item present the person's job and workplace;
\item say how long he or she has done this job, using \cle{for} and \cle{since} correctly;
\item mention one important past event, using \cle{ago};
\item describe what happened \cle{after} that event and what the family did \cle{afterwards};
\item use the \emph{present perfect} and the \emph{simple past} appropriately.
\end{enumerate}

\emph{Conseil méthodologique : organise ta production en trois paragraphes (présentation du métier, un souvenir marquant, la situation actuelle) et vérifie chaque adverbe de temps avant de rendre ta copie.}
\end{exercice}

\begin{methode}[Réussir la composition]
\begin{enumerate}
\item \textbf{Paragraphe 1 — présentation} : nomme le métier et le lieu de travail, puis indique la durée avec \eng{for} + durée et le point de départ avec \eng{since} + date ou événement, au \emph{present perfect}.
\item \textbf{Paragraphe 2 — un souvenir} : raconte un événement passé révolu au \emph{simple past} avec \eng{... ago}, puis enchaîne avec \eng{after} + nom ou proposition, et \eng{afterwards} pour la conséquence.
\item \textbf{Paragraphe 3 — situation actuelle} : reviens au présent ou au \emph{present perfect} pour montrer le lien avec aujourd'hui.
\item \textbf{Relecture} : vérifie que \eng{ago} suit toujours la durée, que \eng{since} n'est jamais suivi d'une durée, et qu'aucun \emph{present perfect} n'est accompagné d'un marqueur de passé révolu (\eng{yesterday}, \eng{last year}, \eng{ago}).
\end{enumerate}
\end{methode}

\newpage

\coursec{Corrigés des exercices}

\coursec{Corrigés des exercices — Leçon 14 : Adverbes de temps (for, since, ago, after, afterwards)}

\begin{corrige}[Exercice 1 — QCM : le bon adverbe]
\begin{enumerate}
\item \textbf{b) since} — \eng{He has been unemployed since he lost his job.} \eng{since} + point de départ (ici une proposition au \emph{simple past}) ; la principale est au \emph{present perfect}.
\item \textbf{c) for} — \eng{My aunt has worked at the central market for ten years.} \eng{for} + durée (\eng{ten years}).
\item \textbf{c) ago} — \eng{The train left Douala two hours ago.} \eng{ago} se place \emph{après} la durée et s'emploie avec le \emph{simple past} (\eng{left}).
\item \textbf{d) afterwards} — \eng{We finished the meeting, and afterwards we went to a restaurant.} \eng{afterwards} (« ensuite ») marque la succession de deux événements passés.
\item \textbf{b) since} — \eng{She has lived in Buea since 2015.} \eng{since} + point de départ daté (\eng{2015}).
\end{enumerate}
\end{corrige}

\begin{corrige}[Exercice 2 — Vrai ou faux ?]
\begin{enumerate}
\item \textbf{False.} \eng{five years} est une durée : il faut \eng{for}. Correction : \eng{I have been a trader for five years.}
\item \textbf{True.} \eng{ago} suit la durée (\eng{three days ago}) et s'emploie avec le \emph{simple past} (\eng{arrived}).
\item \textbf{True.} \eng{Monday} est un point de départ : \eng{since} + \emph{present perfect} (\eng{has been ill}).
\item \textbf{False.} On ne peut pas cumuler \eng{after} et \eng{afterwards} dans la même phrase. Correction : \eng{After the lesson, we played football.} ou \eng{We had a lesson; afterwards we played football.}
\item \textbf{True.} \eng{for a long time} : durée + \emph{present perfect} (\eng{have known}).
\item \textbf{False.} \eng{last month} est un moment passé révolu : il faut le \emph{simple past}. Correction : \eng{He bought a new motorbike last month.}
\end{enumerate}
\end{corrige}

\begin{corrige}[Exercice 3 — Vocabulaire : durée ou point de départ ?]
\begin{center}
\begin{tabularx}{\linewidth}{|X|X|}
\hline
\rowcolor{popblueL} \textbf{for + durée} & \textbf{since + point de départ} \\
\hline
\eng{ten years} & \eng{2010} \\
\hline
\eng{a long time} & \eng{last January} \\
\hline
\eng{six months} & \eng{my childhood} \\
\hline
\eng{two weeks} & \eng{Christmas} \\
\hline
\eng{ages} & \eng{I was born} \\
\hline
\end{tabularx}
\end{center}
Astuce : une durée répond à la question \eng{How long...?} par une longueur de temps ; un point de départ désigne un moment précis (date, événement, proposition).
\end{corrige}

\begin{corrige}[Exercice 4 — Conjugaison : le bon temps]
\begin{enumerate}
\item \eng{met} — \eng{I met my business partner three years ago.} \eng{ago} impose le \emph{simple past} (irrégulier : \emph{meet -- met -- met}).
\item \eng{have built} — \eng{Since then, we have built a successful company.} \eng{since then} impose le \emph{present perfect} (irrégulier : \emph{build -- built -- built}).
\item \eng{has worked} — \eng{My uncle has worked at the post office since 2004.} \eng{since 2004} + action qui continue : \emph{present perfect}.
\item \eng{left} — \eng{She left her village a long time ago.} \eng{a long time ago} : \emph{simple past} (irrégulier : \emph{leave -- left -- left}).
\item \eng{haven't seen} (ou \eng{have not seen}) — \eng{They haven't seen each other since the end of the school year.} \eng{since} + négation : \emph{present perfect}.
\item \eng{finished} (ou \eng{had finished}) — \eng{After he finished his studies, he opened a shop in Douala.} Après \eng{after}, le \emph{simple past} est correct ; le \emph{past perfect} insiste sur l'antériorité.
\end{enumerate}
\end{corrige}

\begin{corrige}[Exercice 5 — Texte à trous : for, since ou ago]
\eng{My uncle has been a civil servant \textbf{for} twenty years. He got this job in 2004, so he has worked in Yaoundé \textbf{since} 2004. He left his village twenty years \textbf{ago}. He has been married \textbf{for} fifteen years, and he has lived in the same house \textbf{since} his wedding. Two months \textbf{ago}, he was promoted, and he has been much happier \textbf{since} then.}\par
Justification : \eng{for} devant les durées (\eng{twenty years}, \eng{fifteen years}) ; \eng{since} devant les points de départ (\eng{2004}, \eng{his wedding}, \eng{then}) ; \eng{ago} après les durées avec un verbe au \emph{simple past} (\eng{left}, \eng{was promoted}).
\end{corrige}

\begin{corrige}[Exercice 6 — Traduction]
\begin{enumerate}
\item \eng{My father has been a taxi driver for twenty years.}
\item \eng{She opened her shop three months ago.}
\item \eng{We have lived in Douala since 2018.}
\item \eng{After the exam, we went home; afterwards, we celebrated.}
\item \eng{I haven't seen my friend since last week.}
\item \eng{They have been waiting for the bus for an hour.}
\end{enumerate}
Remarque : le « depuis » français se traduit par \eng{for} + durée ou \eng{since} + point de départ, toujours avec le \emph{present perfect} (ou \emph{present perfect continuous} pour insister sur l'action en cours, phrase 6). Attention à ne jamais traduire « il y a trois mois » par \eng{for three months} : c'est \eng{three months ago} + \emph{simple past}.
\end{corrige}

\begin{corrige}[Exercice 7 — Transformation]
\begin{enumerate}
\item \eng{She has taught English since 2010.} (ou toute année cohérente : la durée est transformée en point de départ)
\item \eng{He has worked at the bank since 2012.}
\item \eng{They have lived in Bamenda for five years.}
\item \eng{I haven't seen her since June.}
\item \eng{After lunch, we visited the market.}
\end{enumerate}
Point de méthode : passer de \eng{for} à \eng{since} (ou l'inverse) exige de convertir la durée en point de départ, sans changer le \emph{present perfect} de la principale.
\end{corrige}

\begin{corrige}[Exercice 8 — Appariements : la phrase et sa suite logique]
1 -- c ; 2 -- a ; 3 -- d ; 4 -- e ; 5 -- b ; 6 -- f.\par
Phrases complètes :
\begin{enumerate}
\item \eng{The farmers harvested the cocoa two weeks ago, and they have already sold it at the market.}
\item \eng{She has been a nurse for more than ten years.}
\item \eng{He has saved money since he opened his small business.}
\item \eng{After the customers left, the shopkeeper counted the day's money.}
\item \eng{They signed the contract, and afterwards they celebrated with their partners.}
\item \eng{My brother graduated last year, and he has been looking for a job since then.}
\end{enumerate}
Vérification grammaticale : \eng{ago} appelle le \emph{simple past} (1) ; \eng{for} + durée suit un \emph{present perfect} (2) ; \eng{since} + proposition au passé suit un \emph{present perfect} (3, 6) ; \eng{after} introduit une subordonnée au passé (4) ; \eng{afterwards} relie deux événements passés (5).
\end{corrige}

\begin{corrige}[Exercice 9 — Compréhension et langue — type Bac]
\textbf{A. Comprehension}
\begin{enumerate}
\item \eng{She has run her restaurant for almost twenty-five years.}
\item \eng{She opened it in the year 2000, shortly after she lost her job at a textile factory.}
\item \eng{After her husband joined the business, they worked together every day, from six in the morning until late in the evening.}
\item \eng{Since she finished secondary school, their daughter has helped them, and she now manages the accounts.}
\end{enumerate}
\textbf{B. Language work}
\begin{enumerate}
\item Exemples possibles : \eng{for almost twenty-five years} (\eng{for} + durée, avec le \emph{present perfect}) ; \eng{ten years ago} (\eng{ago} placé après la durée, avec le \emph{simple past}) ; \eng{since she finished secondary school} (\eng{since} + proposition au \emph{simple past}, principale au \emph{present perfect}) ; \eng{afterwards} (marque la succession d'événements passés).
\item \eng{Mrs Etame has run her restaurant \textbf{since} the year 2000. She lost her factory job twenty-five years \textbf{ago}.}
\item \eng{Since she \textbf{opened} the restaurant, many customers \textbf{have come} to taste her cooking.}
\end{enumerate}
Barème indicatif (sur 20) : compréhension 8 pts (2 pts par réponse exacte et bien formulée) ; travail sur la langue 12 pts (4 pts pour les deux exemples expliqués, 4 pts pour la complétion, 4 pts pour la conjugaison).
\end{corrige}

\begin{corrige}[Exercice 10 — Traduction — type Bac]
\eng{My uncle has been a carpenter for thirty years. He learnt this trade a very long time ago, when he was still a teenager. After his apprenticeship, he opened a small workshop in Bafoussam; afterwards, he trained several young people from the neighbourhood. Since 2010, he has also been making furniture for the town's schools. He has never regretted his choice.}\par
Points de vigilance :
\begin{itemize}
\item « depuis trente ans » = \eng{for thirty years} + \emph{present perfect} (jamais \eng{since thirty years}) ;
\item « il y a très longtemps » = \eng{a very long time ago} + \emph{simple past} ;
\item « ensuite » = \eng{afterwards} (jamais \eng{after} seul en fin de phrase) ;
\item « depuis 2010 » = \eng{since 2010} + \emph{present perfect} (continuous possible pour insister sur l'activité) ;
\item « menuisier » = \eng{carpenter} ; « apprentissage » = \eng{apprenticeship} ; « atelier » = \eng{workshop}.
\end{itemize}
Barème indicatif (sur 10) : exactitude des adverbes de temps 4 pts ; choix des temps verbaux 3 pts ; vocabulaire 2 pts ; orthographe et fluidité 1 pt.
\end{corrige}

\begin{corrige}[Exercice 11 — Expression écrite — type Bac]
Proposition de rédaction modèle :\par
\eng{My mother is a seamstress. She has worked in her own workshop near the Mokolo market for more than fifteen years. She has been passionate about sewing since she was a young girl, and she opened her workshop in 2009, just after she finished her training.}\par
\eng{Ten years ago, something important happened: a famous tailor from Yaoundé noticed her dresses and ordered twenty outfits for a wedding. After that order, her business grew very quickly. She hired two apprentices, and afterwards she bought a second sewing machine. Our whole family was very proud of her.}\par
\eng{Since then, my mother has trained many young girls from our neighbourhood. She has never complained about her long working days, because her job has paid for our education. For me, she is a perfect example of courage and hard work.}\par
(environ 140 mots ; à adapter à la situation personnelle de l'élève)\par
Points de vigilance :
\begin{itemize}
\item \eng{for more than fifteen years} (durée) et \eng{since she was a young girl} (point de départ) avec le \emph{present perfect} ;
\item \eng{ten years ago} avec le \emph{simple past} (\eng{happened}, \eng{noticed}, \eng{ordered}) ;
\item \eng{after that order} (préposition + nom) et \eng{afterwards} (adverbe, « ensuite ») bien distingués ;
\item aucun \emph{present perfect} avec un marqueur de passé révolu.
\end{itemize}
Barème indicatif (sur 20) : respect des consignes et de la longueur 4 pts ; correction des adverbes de temps 5 pts ; emploi des temps verbaux 5 pts ; vocabulaire et cohérence 4 pts ; orthographe et présentation 2 pts.
\end{corrige}

\newpage

\coursec{Fiche bilan}

\begin{popfigure}
\begin{center}\tcbox[colback=popblue,colframe=popblue,coltext=white,arc=9pt,boxsep=4pt,left=6pt,right=6pt]{\bfseries\small Adverbs of Time}\end{center}
\vspace{-2pt}
\begin{tcbraster}[raster columns=3,raster equal height=rows,raster column skip=6pt,raster row skip=6pt,enhanced,blankest]
\begin{tcolorbox}[enhanced,colback=poporange,colframe=poporange,coltext=white,boxrule=0.9pt,arc=3pt,left=4pt,right=4pt,top=3pt,bottom=3pt,boxsep=1pt,halign=left]\bfseries\scriptsize FOR + durée (for two years)\end{tcolorbox}
\begin{tcolorbox}[enhanced,colback=popgreen,colframe=popgreen,coltext=white,boxrule=0.9pt,arc=3pt,left=4pt,right=4pt,top=3pt,bottom=3pt,boxsep=1pt,halign=left]\bfseries\scriptsize SINCE + point de départ (since 2020)\end{tcolorbox}
\begin{tcolorbox}[enhanced,colback=poppurple,colframe=poppurple,coltext=white,boxrule=0.9pt,arc=3pt,left=4pt,right=4pt,top=3pt,bottom=3pt,boxsep=1pt,halign=left]\bfseries\scriptsize AGO + durée, passé (two days ago)\end{tcolorbox}
\begin{tcolorbox}[enhanced,colback=poppink,colframe=poppink,coltext=white,boxrule=0.9pt,arc=3pt,left=4pt,right=4pt,top=3pt,bottom=3pt,boxsep=1pt,halign=left]\bfseries\scriptsize AFTER + nom / phrase (after work)\end{tcolorbox}
\begin{tcolorbox}[enhanced,colback=popgold,colframe=popgold,coltext=white,boxrule=0.9pt,arc=3pt,left=4pt,right=4pt,top=3pt,bottom=3pt,boxsep=1pt,halign=left]\bfseries\scriptsize AFTERWARDS seul, en fin de phrase\end{tcolorbox}
\begin{tcolorbox}[enhanced,colback=popdark,colframe=popdark,coltext=white,boxrule=0.9pt,arc=3pt,left=4pt,right=4pt,top=3pt,bottom=3pt,boxsep=1pt,halign=left]\bfseries\scriptsize Temps verbaux present perfect past simple\end{tcolorbox}
\end{tcbraster}

\legende{Carte mentale : les adverbes de temps \eng{for, since, ago, after, afterwards}.}
\end{popfigure}

\begin{aretenir}[L'essentiel de la leçon]
\textbf{Lexique clé (English / Français) :}
\begin{itemize}
  \item \eng{for} : pendant / depuis (durée) ; \eng{since} : depuis (point de départ)
  \item \eng{ago} : il y a ; \eng{after} : après ; \eng{afterwards} : ensuite, après
  \item \eng{duration} (n.) : durée ; \eng{starting point} : point de départ
  \item \eng{so far / up to now} : jusqu'à présent ; \eng{ever since} : depuis lors
\end{itemize}

\textbf{Règles à connaître par cœur :}
\begin{center}
\begin{tabularx}{\linewidth}{|l|X|X|}\hline
\rowcolor{popblueL} \textbf{Adverbe} & \textbf{Emploi + temps verbal} & \textbf{Exemple} \\ \hline
\eng{for} & + durée ; present perfect (ou past simple si l'action est terminée) & \eng{She has worked here for ten years.} \\ \hline
\eng{since} & + point de départ (date, événement) ; present perfect & \eng{He has lived in Douala since 2015.} \\ \hline
\eng{ago} & + durée, placé après ; toujours past simple & \eng{They left Buea two years ago.} \\ \hline
\eng{after} & + nom ou proposition ; jamais seul en fin de phrase & \eng{After the meeting, we went home.} \\ \hline
\eng{afterwards} & seul, en début ou en fin de phrase ; past simple & \eng{We had lunch and afterwards we rested.} \\ \hline
\end{tabularx}
\end{center}

\textbf{Méthodes express :}
\begin{enumerate}
  \item Durée (\eng{two hours, five years}) $\rightarrow$ \eng{for} ; date ou événement de départ (\eng{Monday, 2019, I was born}) $\rightarrow$ \eng{since}.
  \item « Il y a + durée » $\rightarrow$ durée + \eng{ago} + verbe au past simple.
  \item \eng{after} exige toujours un complément ; sinon, utilise \eng{afterwards}.
  \item Action commencée dans le passé et qui continue $\rightarrow$ present perfect + \eng{for/since}.
\end{enumerate}
\end{aretenir}

\begin{exemplebox}[Exemples modèles à retenir]
\begin{itemize}
  \item \eng{The cocoa farmers have used this cooperative for fifteen years.} « Les planteurs de cacao utilisent cette coopérative depuis quinze ans. »
  \item \eng{My uncle has been a taxi driver in Bamenda since 2010.} « Mon oncle est chauffeur de taxi à Bamenda depuis 2010. »
  \item \eng{The new market opened three months ago.} « Le nouveau marché a ouvert il y a trois mois. »
  \item \eng{After the interview, the manager offered her the job.} « Après l'entretien, le directeur lui a offert le poste. »
  \item \eng{We visited the factory; afterwards, we wrote our reports.} « Nous avons visité l'usine ; ensuite, nous avons rédigé nos rapports. »
\end{itemize}
\end{exemplebox}

\begin{attention}[Pièges fréquents des francophones]
\begin{itemize}
  \item Ne dis jamais \eng{I am here since 2020} : le présent simple est impossible avec \eng{since} pour une action qui continue. Dis \eng{I have been here since 2020.}
  \item \eng{ago} se place \emph{après} la durée : \eng{two weeks ago}, jamais \eng{ago two weeks}.
  \item Pas de present perfect avec \eng{ago} : \eng{He arrived two days ago}, jamais \eng{He has arrived two days ago}.
  \item « Depuis » ne se traduit pas par \eng{from} : \eng{since Monday}, sauf dans la structure \eng{from Monday to Friday} (« du lundi au vendredi »).
  \item \eng{after} ne peut pas terminer une phrase seul : \eng{We ate and went home afterwards}, jamais \eng{...went home after}.
  \item Attention au faux ami : \eng{actually} signifie « en fait », pas « actuellement » (\eng{at present, currently}).
\end{itemize}
\end{attention}

\begin{exoresolu}[Transformation guidée]
Énoncé : transforme la phrase \eng{I came to Yaoundé in 2019.} en utilisant \eng{since}, puis \eng{for} (nous sommes en 2024).
\tcblower
Solution détaillée :
\begin{enumerate}
  \item L'action a commencé en 2019 et continue aujourd'hui : il faut le \cle{present perfect} (\eng{have been}).
  \item Avec un point de départ : \eng{I have been in Yaoundé since 2019.} « Je suis à Yaoundé depuis 2019. »
  \item Avec une durée (2024 $-$ 2019 = 5 ans) : \eng{I have been in Yaoundé for five years.} « Je suis à Yaoundé depuis cinq ans. »
\end{enumerate}
\end{exoresolu}

\begin{aretenir}[Checklist avant le Bac]
\begin{itemize}
  \item $\square$ Je distingue \eng{for} (+ durée) et \eng{since} (+ point de départ) sans hésiter.
  \item $\square$ Je place \eng{ago} après la durée : \eng{three days ago}, jamais \eng{ago three days}.
  \item $\square$ J'utilise le past simple avec \eng{ago}, jamais le present perfect.
  \item $\square$ J'utilise le present perfect avec \eng{for} et \eng{since} pour une action qui continue.
  \item $\square$ Je sais que \eng{after} ne peut pas terminer une phrase seul ; je dis \eng{afterwards}.
  \item $\square$ Je ne traduis pas « depuis » par \eng{from} : \eng{since Monday}, pas \eng{from Monday} (sauf \eng{from Monday to Friday}).
  \item $\square$ Je construis la question de durée avec \eng{How long have you...?}
  \item $\square$ Je sais que \eng{since} peut introduire une proposition au past simple : \eng{since I arrived}.
  \item $\square$ J'évite le faux ami \eng{actually} (= en fait) pour « actuellement » (= \eng{at present}).
  \item $\square$ Je peux transformer une phrase : \eng{I came here in 2020.} $\rightarrow$ \eng{I have been here since 2020.}
  \item $\square$ Je vérifie l'accord du temps verbal après chaque adverbe dans mes rédactions.
\end{itemize}
\end{aretenir}

\findecours

\end{document}
